% Generated mechanically from frozen input inputs/p11/ja/artin.tex.
% Do not edit this generated file; edit the bound locale source instead.

% OK, start here.
%

\setcounter{chapter}{97}
\chapter{Artin の公理}


\phantomsection
\label{artin-section-phantom}





\section{序論}
\label{artin-section-introduction}

\noindent
本章では、関手が代数空間により表現されるための Artin の公理を論じる。
参考文献として、論文
\cite{ArtinI}, \cite{ArtinII}, \cite{ArtinVersal}
を挙げる。

\medskip\noindent
本章の記法・規約・用語の一部は扱いにくく、経験を積んだ読者には
向きが逆に見えるかもしれない。これは意図的である。
説明については『Quot』の節 \ref{quot-section-conventions} を参照せよ。

\medskip\noindent
$S$ を局所 Noether な基礎スキームとする。
$$
p : \mathcal{X} \longrightarrow (\Sch/S)_{fppf}
$$
をグルーポイドをファイバーとする圏とする。$S$ 上有限型の体 $k$ 上の
$\mathcal{X}$ の対象を $x_0$ とする。本章を通じて重要な役割を果たすのは、
$k$ 上の $x_0$ に付随する前変形圏
（『形式的変形理論』の定義
\ref{formal-defos-definition-predeformation-category} を参照）
$$
\mathcal{F}_{\mathcal{X}, k, x_0}
\longrightarrow
\{\text{Artinian local }S\text{-algebras with residue field }k\}
$$
である。$\mathcal{X}$ に対する Rim--Schlessinger 条件 (RS) を導入し、
これにより $\mathcal{F}_{\mathcal{X}, k, x_0}$ が変形圏、すなわち
$\mathcal{F}_{\mathcal{X}, k, x_0}$ 自身が (RS) を満たすことを示す。
また、$k$ を有限次拡大で置き換えたときに
$\mathcal{F}_{\mathcal{X}, k, x_0}$ がどう変わるかを論じ、
接空間についても論じる。

\medskip\noindent
次に、$\mathcal{X}$ の形式対象 $\xi = (\xi_n)$ を論じる。これは、
剰余体が $S$ 上有限型である Noether 完備局所 $S$-代数 $R$ に対し、
その極大イデアルの累乗による商 $R/\mathfrak m^n$ 上にある対象の逆系である。
これは、ある $x_0$ と $k$ に対して
$\mathcal{F}_{\mathcal{X}, k, x_0}$ の形式対象をもつことと同じである。
$R$ 上の $\mathcal{X}$ の対象でこの逆系を生じるものが存在するとき、
形式対象は有効であるという。$\mathcal{X}$ の形式対象が半普遍的であるとは、
それが $\mathcal{F}_{\mathcal{X}, k, x_0}$ の半普遍的形式対象を生じることをいう。
最後に、有限型 $S$-スキーム $U$、$U$ 上の $\mathcal{X}$ の対象 $x$、
および閉点 $u_0 \in U$ が与えられたとき、完備局所環
$\mathcal{O}_{U, u_0}^\wedge$ 上に誘導される形式対象が半普遍的ならば、
$x$ は $u_0$ において半普遍的であるという。

\medskip\noindent
以上を準備すると、Artin の著名な定理を述べられる。
次が成り立つならば、$\mathcal{X}$ は代数スタックである。
\begin{enumerate}
\item すべての $s \in S$ に対して $\mathcal{O}_{S, s}$ は G-環である
\footnote{これは $S$ が G-スキームであることを意味する。『性質』の節
\ref{properties-section-G-scheme} を参照せよ。}
\item $\Delta : \mathcal{X} \to \mathcal{X} \times \mathcal{X}$ は
代数空間により表現可能である。
\item $\mathcal{X}$ はエタール位相に関するスタックである。
\item $\mathcal{X}$ は極限保持的である。
\item $\mathcal{X}$ は (RS) を満たす。
\item 変形圏 $\mathcal{F}_{\mathcal{X}, k, x_0}$ の接空間と
無限小自己同型の空間は有限次元である。
\item 形式対象は有効である。
\item $\mathcal{X}$ は半普遍性の開性を満たす。
\end{enumerate}
これは補題 \ref{artin-lemma-diagonal-representable} である。
わずかに強い結果については命題
\ref{artin-proposition-second-diagonal-representable} も参照せよ。
どの関手 $F : (\Sch/S)_{fppf}^{opp} \to \textit{Sets}$ が
代数空間であるかを特徴づける類似の命題については、
第 \ref{artin-section-algebraic-spaces} 節を参照せよ。

\medskip\noindent
Artin の定理の証明の概略を述べる。
まず、(RS) と接空間および自己同型空間の有限次元性を用いて、
十分多くの半普遍的形式対象が存在することを示す。たとえば
『形式的変形理論』の補題
\ref{formal-defos-lemma-minimal-groupoid-in-functors-construction}
を参照せよ。これらの形式対象は仮定により有効である。
有効な形式対象は、有限型 $S$-スキーム $U$ 上の対象 $x$ により
「近似」できる。補題 \ref{artin-lemma-approximate} を参照せよ。
この近似には、$S$ の局所環が G-環であることと
$\mathcal{X}$ が極限保持的であることを用いる。
これはおそらく証明の最も難しい部分であり、一般 N\'eron 脱特異化を用いて、
Noether 局所 G-環上の代数方程式の形式解をヘンゼル化における解で
近似することに依存する。
次に、半普遍性の開性により、$U$ を縮小した後、
$x$ が $U$ のすべての閉点で半普遍的であると仮定できる。
以上の準備のもとで、$U \to \mathcal{X}$ が滑らかな射であることを示す。
十分多くの $U \to \mathcal{X}$ を取ることにより、
$\mathcal{X}$ の「滑らかなアトラス」が得られ、これから
$\mathcal{X}$ が代数スタックであることが従う。

\medskip\noindent
グルーポイドをファイバーとする圏 $\mathcal{X}$ に対して
Artin の公理を確認するとき、最も難しい段階はしばしば
半普遍性の開性を検証することである。以下の議論では、
$\mathcal{X}/S$ が上に挙げた他の条件をすでに満たすと仮定する。
本章では、$\mathcal{X}$ が半普遍性の開性を満たすことを示すための
二つの方法を提示する。
\begin{enumerate}
\item 第一の方法では、(RS*) と呼ばれる、より強い Rim--Schlessinger 条件と、
より強い形式的有効性を仮定する。後者は本質的に、
無限小厚化の逆系上の対象が有効であることを要求する。
これらの仮定のもとでは半普遍性の開性が自動的に従う。
補題 \ref{artin-lemma-SGE-implies-openness-versality} を参照せよ。
ここでは、$\mathcal{X}$ が Noether スキームの圏だけでなく、
$S$ 上のすべてのスキームの圏で定義されていることを
本質的に用いている点に注意されたい。
\item Artin に従う第二の方法では、$\mathcal{X}$ に障害理論が備わることを
要求する。この障害理論が適切な意味で「積と可換」ならば、
$\mathcal{X}$ は半普遍性の開性を満たす。
補題 \ref{artin-lemma-get-openness-obstruction-theory} を参照せよ。
\end{enumerate}
障害理論は多様な仕方で公理化でき、実際、多くの変種
（しばしば特定のモジュライ・スタックに合わせたもの）が文献に見られる。
導来圏を用いる一つの変種を補題 \ref{artin-lemma-dual-openness} で説明する。
これは、文献中の変形理論の計算から自然に現れることが多い。

\medskip\noindent
第 \ref{artin-section-algebraic-spaces-noetherian} 節では、
$S$ 上の局所 Noether スキームの圏
$(\textit{Noetherian}/S)_\etale$ 上で定義された関手に対して
議論を成り立たせるために何を修正すべきかを論じる。

\medskip\noindent
本章の最終節では、Artin の公理の応用として
縮約の存在に関する Artin の定理を証明する。
第 \ref{artin-section-contractions} 節を参照せよ。
定理は大まかに次を主張する。$S$ 上分離的かつ有限型の代数空間 $X'$、
閉部分集合 $T' \subset |X'|$、および形式的修正
$$
\mathfrak{f} : X'_{/T'} \longrightarrow \mathfrak{X}
$$
が与えられ、$\mathfrak{X}$ が $S$ 上の Noether 形式代数空間であるとする。
このとき、「縮約を実現する」固有射 $f : X' \to X$ が存在する。
すなわち、$\mathfrak{X} = X_{/T}$ という同一視であって、
$T = f(T')$ かつ
$\mathfrak{f} = f_{/T'} : X'_{/T'} \to X_{/T}$ となるものが存在し、
さらに $f$ は $X \setminus T$ 上で同型である。
証明は $S$ 上の局所 Noether スキームの圏上の関手 $F$ を定義し、
$F$ に対して Artin の公理を証明することで進む。
興味深いことに、この Artin の公理の応用では、
半普遍性の開性が最も証明しにくい条件ではない。
むしろ $F$ が極限保持的であることの証明に、
多くの作業と準備結果を要する。





\section{規約}
\label{artin-section-conventions}

\noindent
本章で用いる規約は『代数スタック』の章のものと同じである。
『代数スタック』の節 \ref{algebraic-section-conventions} を参照せよ。
本章では基礎スキーム $S$ はしばしば局所 Noether である
（ただし、結果を述べる際には常にこの条件を改めて明記する）。





\section{前変形圏}
\label{artin-section-predeformation-categories}

\noindent
$S$ を局所 Noether な基礎スキームとする。
$$
p : \mathcal{X} \longrightarrow (\Sch/S)_{fppf}
$$
をグルーポイドをファイバーとする圏とする。$k$ を体とし、
$\Spec(k) \to S$ を有限型射とする
（『射』の補題 \ref{morphisms-lemma-point-finite-type} を参照）。
簡単に、{\it $k$ は $S$ 上有限型の体である}ということもある。
$\Spec(k)$ 上にある $\mathcal{X}$ の対象を $x_0$ とする。
$S$, $\mathcal{X}$, $k$, $x_0$ が与えられたとき、
『形式的変形理論』の定義
\ref{formal-defos-definition-predeformation-category}
にいう前変形圏を構成する。この構成は『形式的変形理論』の注意
\ref{formal-defos-remark-localize-cofibered-groupoid}
の構成に似ている。

\medskip\noindent
まず、『射』の補題 \ref{morphisms-lemma-point-finite-type} により、
アフィン開集合 $\Spec(\Lambda) \subset S$ を、
$\Spec(k) \to S$ が $\Spec(\Lambda)$ を経由し、
付随する環準同型 $\Lambda \to k$ が有限となるように選べる。
これにより圏 $\mathcal{C}_\Lambda$ が得られる。
『形式的変形理論』の定義
\ref{formal-defos-definition-CLambda} を参照せよ。
圏 $\mathcal{C}_\Lambda$ は、標準的同値を除いて、
$S$ のアフィン開集合 $\Spec(\Lambda)$ の選択に依存しない。
実際、$\mathcal{C}_\Lambda$ は、構造射の分解
\begin{equation}
\label{artin-equation-factor}
\Spec(k) \to \Spec(A) \to S
\end{equation}
であって、$A$ が Artin 局所環であり、
$\Spec(k) \to \Spec(A)$ が $k$ を $A$ の剰余体と同一視する環準同型
$A \to k$ に対応するものからなる圏の反対圏と同値である。

\medskip\noindent
$\mathcal{F} = \mathcal{F}_{\mathcal{X}, k, x_0}$ を次の圏とする。
\begin{enumerate}
\item 対象は、$p(x) = \Spec(A)$ であって $A$ が Artin 局所環であり、
$p(x_0) \to p(x) \to S$ が (\ref{artin-equation-factor}) のような分解となる、
$\mathcal{X}$ の射 $x_0 \to x$ である。
\item 射 $(x_0 \to x) \to (x_0 \to x')$ は、
$\mathcal{X}$ における可換図式
$$
\xymatrix{
x & & x' \ar[ll] \\
& x_0 \ar[lu] \ar[ru]
}
$$
である（矢印の向きが反転していることに注意）。
\end{enumerate}
$x_0 \to x$ が $\mathcal{F}$ の対象ならば、$p(x) = \Spec(A)$ と書くと、
$\mathcal{C}_\Lambda$ の対象 $A$ が得られる。
しばしば $x_0 \to x$ または $x$ は $A$ 上にあるという。
$A$ 上にある $x_0 \to x$ と $A'$ 上にある $x_0 \to x'$ の間の
$\mathcal{F}$ の射は、$\mathcal{X}$ の射 $x' \to x$ に対応し、
したがって射
$p(x' \to x) : \Spec(A') \to \Spec(A)$ に対応する。
さらにこれは環準同型 $A \to A'$ に対応する。
$\mathcal{X}$ は $S$ 上のスキームの圏上の圏なので、
$A \to A'$ は $\Lambda$-代数準同型である。したがって関手
\begin{equation}
\label{artin-equation-predeformation-category}
p : \mathcal{F} = \mathcal{F}_{\mathcal{X}, k, x_0}
\longrightarrow
\mathcal{C}_\Lambda.
\end{equation}
を得る。$\mathcal{C}_\Lambda$ の対象 $A$ 上のファイバー圏を
$\mathcal{F}(A)$ と記す。$\mathcal{F}(A)$ の対象は、
$x$ が $\Spec(A)$ 上にあり、$x_0 \to x$ が
$\Spec(k) \to \Spec(A)$ 上にあるような、$\mathcal{X}$ の射
$x_0 \to x$ にほかならない。

\begin{lemma}
\label{artin-lemma-predeformation-category}
上で定義した関手 $p : \mathcal{F} \to \mathcal{C}_\Lambda$ は
前変形圏である。
\end{lemma}

\begin{proof}
示すべきことは、(a) $\mathcal{F}$ が $\mathcal{C}_\Lambda$ 上
グルーポイドをファイバーとする余ファイバー圏であることと、
(b) $\mathcal{F}(k)$ が、ただ一つの対象とただ一つの射をもつ圏と
同値であることである。

\medskip\noindent
(a) の証明。$\mathcal{X}$ のファイバー圏がグルーポイドなので、
$\mathcal{C}_\Lambda$ 上の $\mathcal{F}$ のファイバー圏も
グルーポイドである。$A \to A'$ を $\mathcal{C}_\Lambda$ の射とし、
$x_0 \to x$ を $\mathcal{F}(A)$ の対象とする。
$\mathcal{X}$ はグルーポイドをファイバーとするから、
$\Spec(A') \to \Spec(A)$ 上にある射 $x' \to x$ を見いだせる。
合成 $A \to A' \to k$ は与えられた写像 $A \to k$ に等しいので、
引き戻しの同型を除く一意性により、$\Spec(k) \to \Spec(A')$ による
$x'$ の引き戻しは $x_0$ である。すなわち、
$\Spec(k) \to \Spec(A')$ 上にある射 $x_0 \to x'$ で、
$x_0 \to x$ および $x' \to x$ と両立するものが存在する。
これにより $\mathcal{F}$ が順像をもつことが分かる。
最後に『圏』の補題
\ref{categories-lemma-fibred-groupoids} の双対を用いる。

\medskip\noindent
(b) の証明。$A = k$ ならば $\Spec(k) = \Spec(A)$ であり、
$\mathcal{X}$ は $(\Sch/S)_{fppf}$ 上グルーポイドをファイバーとするから、
$\mathcal{F}(k)$ の任意の対象 $x_0 \to x$ に対して
射 $x_0 \to x$ は同型である。したがって
$\mathcal{F}(k)$ のすべての対象は $x_0 \to x_0$ と同型である。
明らかに、$\mathcal{F}$ における $x_0 \to x_0$ の自己射は
恒等射だけである。
\end{proof}


\noindent
$S$ を局所 Noether な基礎スキームとする。
$(\Sch/S)_{fppf}$ 上グルーポイドをファイバーとする圏の間の $1$-射を
$F : \mathcal{X} \to \mathcal{Y}$ とする。
$k$ を $S$ 上有限型の体とする。
$\Spec(k)$ 上にある $\mathcal{X}$ の対象を $x_0$ とし、
$y_0 = F(x_0)$ とおく。これは $\Spec(k)$ 上にある
$\mathcal{Y}$ の対象である。このとき $F$ は、
$\mathcal{C}_\Lambda$ 上の余ファイバー圏の関手
\begin{equation}
\label{artin-equation-functoriality}
F :
\mathcal{F}_{\mathcal{X}, k, x_0}
\longrightarrow
\mathcal{F}_{\mathcal{Y}, k, y_0}
\end{equation}
を誘導する。実際、$\mathcal{F}_{\mathcal{X}, k, x_0}(A)$ の対象
$x_0 \to x$ に、$\mathcal{F}_{\mathcal{Y}, k, y_0}(A)$ の対象
$F(x_0) \to F(x)$ を対応させる。

\begin{lemma}
\label{artin-lemma-formally-smooth-on-deformation-categories}
$S$ を局所 Noether スキームとする。
$F : \mathcal{X} \to \mathcal{Y}$ を
$(\Sch/S)_{fppf}$ 上グルーポイドをファイバーとする圏の $1$-射とする。
次のいずれかを仮定する。
\begin{enumerate}
\item $F$ は対象について形式的に滑らかである
（『表現可能性の判定条件』の節
\ref{criteria-section-formally-smooth}）。
\item $F$ は代数空間により表現可能かつ形式的に滑らかである。
\item $F$ は代数空間により表現可能かつ滑らかである。
\end{enumerate}
このとき、$S$ 上の任意の有限型の体 $k$ と $k$ 上の
$\mathcal{X}$ の対象 $x_0$ に対し、関手
(\ref{artin-equation-functoriality}) は『形式的変形理論』の定義
\ref{formal-defos-definition-smooth-morphism}
の意味で滑らかである。
\end{lemma}

\begin{proof}
(1) は定義を展開すれば分かる。
仮定 (2) は『表現可能性の判定条件』の補題
\ref{criteria-lemma-representable-by-spaces-formally-smooth}
により (1) を含意する。
仮定 (3) は『空間の射についてさらに』の補題
\ref{spaces-more-morphisms-lemma-smooth-formally-smooth}
と『代数スタック』の補題
\ref{algebraic-lemma-representable-transformations-property-implication}
の原理により (2) を含意する。
\end{proof}

\begin{lemma}
\label{artin-lemma-fibre-product-deformation-categories}
$S$ を局所 Noether スキームとする。
$$
\xymatrix{
\mathcal{W} \ar[d] \ar[r] & \mathcal{Z} \ar[d] \\
\mathcal{X} \ar[r] & \mathcal{Y}
}
$$
を $(\Sch/S)_{fppf}$ 上グルーポイドをファイバーとする圏の
$2$-ファイバー積とする。$k$ を $S$ 上有限型の体とし、
$w_0$ を $k$ 上の $\mathcal{W}$ の対象とする。
$x_0, z_0, y_0$ を、図式の射による $w_0$ の像とする。このとき
$$
\xymatrix{
\mathcal{F}_{\mathcal{W}, k, w_0} \ar[d] \ar[r] &
\mathcal{F}_{\mathcal{Z}, k, z_0} \ar[d] \\
\mathcal{F}_{\mathcal{X}, k, x_0} \ar[r] & \mathcal{F}_{\mathcal{Y}, k, y_0}
}
$$
は前変形圏のファイバー積である。
\end{lemma}

\begin{proof}
定義を展開すれば分かる。詳細は省略する。
\end{proof}





\section{押し出しとスタック}
\label{artin-section-pushouts}

\noindent
本節では、代数スタックがある種の押し出しに関してよく振る舞うことを示す。
本節の結果は任意の基礎スキーム上で成り立つ。

\medskip\noindent
次の補題は $Y$, $X'$, $X$, $Y'$ が代数空間の場合にも正しい
（将来の参照をここに挿入）。

\begin{lemma}
\label{artin-lemma-pushout}
\begin{slogan}
代数スタックは（強）Rim--Schlessinger 条件を満たす
\end{slogan}
$S$ をスキームとする。
$$
\xymatrix{
X \ar[r] \ar[d] & X' \ar[d] \\
Y \ar[r] & Y'
}
$$
を $S$ 上のスキームの圏における押し出しとし、$X \to X'$ は無限小厚化、
$X \to Y$ はアフィンであるとする。『射についてさらに』の補題
\ref{more-morphisms-lemma-pushout-along-thickening} を参照せよ。
$\mathcal{Z}$ を $S$ 上の代数スタックとする。このときファイバー圏の関手
$$
\mathcal{Z}_{Y'}
\longrightarrow
\mathcal{Z}_Y \times_{\mathcal{Z}_X} \mathcal{Z}_{X'}
$$
は圏の同値である。
\end{lemma}

\begin{proof}
$y'$ を左辺の対象とする。$Y'$ 上のスキームの圏上の層
$\mathit{Isom}(y', y')$ は、$Y'$ 上の代数空間 $I$ により表現可能である。
『代数スタック』の補題
\ref{algebraic-lemma-representable-diagonal} を参照せよ。
$Y'$ はスキームの圏だけでなく代数空間の圏でも押し出しであるから、
補題の関手は充満忠実であると結論できる。『空間の押し出し』の補題
\ref{spaces-pushouts-lemma-pushout-along-thickening-schemes}
を参照せよ。

\medskip\noindent
$(y, x', f)$ を右辺の対象とする。ここで
$f : y|_X \to x'|_X$ は同型である。証明を完了するには、
$Y$ と $X'$ への制限が、$f$ と両立する仕方でそれぞれ
$y$ と $x'$ に一致するような $\mathcal{Z}_{Y'}$ の対象 $y'$ を
構成しなければならない。実際、$Y'$ 上 fppf 局所的に $y'$ を
構成すれば十分である。『スタック』の補題
\ref{stacks-lemma-characterize-essentially-surjective-when-ff}
を参照せよ。表現可能な代数スタック $\mathcal{W}$ と、
全射かつ滑らかな射 $\mathcal{W} \to \mathcal{Z}$ を選ぶ。このとき
$$
(\Sch/Y)_{fppf} \times_{y, \mathcal{Z}} \mathcal{W}
\quad\text{and}\quad
(\Sch/X')_{fppf} \times_{x', \mathcal{Z}} \mathcal{W}
$$
は、それぞれ $Y$ と $X'$ 上滑らかな代数空間 $V$ と $U'$ により
表現される代数スタックである。同型 $f$ は、$X$ 上の同型
$\varphi : V \times_Y X \to U' \times_{X'} X$ を誘導する。
『空間の押し出し』の補題
\ref{spaces-pushouts-lemma-pushout-along-thickening} および
\ref{spaces-pushouts-lemma-equivalence-categories-spaces-pushout-flat}
により、押し出し $V' = V \amalg_{V \times_Y X} U'$ は
$Y'$ 上滑らかな代数空間であり、その $Y$ と $X'$ への基底変換は、
$\varphi$ と両立する仕方で $V$ と $U'$ を復元する。

\medskip\noindent
$\mathcal{W}$ を表現する代数空間を $W$ とする。射影
$V \to W$ と $U' \to W$ は
$V \times_Y X \cong U' \times_{X'} X$ 上の射として一致するので、
押し出しの普遍性により代数空間の射 $V' \to W$ が定まる。
スキーム $Y_1'$ と全射エタール射 $Y_1' \to V'$ を選ぶ。
$Y_1 = Y \times_{Y'} Y_1'$,
$X_1' = X' \times_{Y'} Y_1'$,
$X_1 = X \times_{Y'} Y_1'$ とおく。合成
$$
(\Sch/Y_1') \to (\Sch/V') \to (\Sch/W) = \mathcal{W} \to \mathcal{Z}
$$
は $2$-Yoneda の補題により、$Y_1'$ 上の $\mathcal{Z}$ の対象 $y_1'$ に
対応する。その $Y_1$ と $X_1'$ への制限は、$f|_{X_1}$ と両立する仕方で
$y|_{Y_1}$ と $x'|_{X_1'}$ に一致する。
こうして $Y'$ 上滑らか局所的に所望の対象を構成でき、証明が完了する。
\end{proof}








\section{Rim--Schlessinger 条件}
\label{artin-section-RS}

\noindent
次の定義の動機は、補題 \ref{artin-lemma-pushout}、および
『形式的変形理論』の定義 \ref{formal-defos-definition-RS} と
補題 \ref{formal-defos-lemma-RS-2-categorical} に由来する。

\begin{definition}
\label{artin-definition-RS}
$S$ を局所 Noether スキームとする。$\mathcal{Z}$ を
$(\Sch/S)_{fppf}$ 上のグルーポイドをファイバーとする圏とする。
$S$ 上のスキームの圏における任意の押し出し
$$
\xymatrix{
X \ar[r] \ar[d] & X' \ar[d] \\
Y \ar[r] & Y' = Y \amalg_X X'
}
$$
であって、
\begin{enumerate}
\item $X$, $X'$, $Y$, $Y'$ が局所 Artin 環のスペクトルであり、
\item $X$, $X'$, $Y$, $Y'$ が $S$ 上有限型であり、かつ
\item $X \to X'$（したがって $Y \to Y'$）が閉埋め込みである
\end{enumerate}
ものすべてに対して、ファイバー圏の関手
$$
\mathcal{Z}_{Y'}
\longrightarrow
\mathcal{Z}_Y \times_{\mathcal{Z}_X} \mathcal{Z}_{X'}
$$
が圏同値であるとき、$\mathcal{Z}$ は {\it 条件 (RS)} を満たすという。
\end{definition}

\noindent
$A$ を剰余体 $k$ をもつ Artin 局所環とすると、任意の射
$\Spec(A) \to S$ がアフィンかつ有限型であることと、誘導される射
$\Spec(k) \to S$ が有限型であることは同値である。『射』の補題
\ref{morphisms-lemma-Artinian-affine} および
\ref{morphisms-lemma-artinian-finite-type} を参照せよ。

\begin{lemma}
\label{artin-lemma-algebraic-stack-RS}
$\mathcal{X}$ を局所 Noether な基礎 $S$ 上の代数スタックとする。
このとき $\mathcal{X}$ は (RS) を満たす。
\end{lemma}

\begin{proof}
定義と補題 \ref{artin-lemma-pushout} から直ちに従う。
\end{proof}

\begin{lemma}
\label{artin-lemma-fibre-product-RS}
$S$ をスキームとする。$p : \mathcal{X} \to \mathcal{Y}$ と
$q : \mathcal{Z} \to \mathcal{Y}$ を、$(\Sch/S)_{fppf}$ 上の
グルーポイドをファイバーとする圏の $1$-射とする。
$\mathcal{X}$, $\mathcal{Y}$, $\mathcal{Z}$ が (RS) を満たすならば、
$\mathcal{X} \times_\mathcal{Y} \mathcal{Z}$ も (RS) を満たす。
\end{lemma}

\begin{proof}
これは形式的である。定義 \ref{artin-definition-RS} に現れる形の図式
$$
\xymatrix{
X \ar[r] \ar[d] & X' \ar[d] \\
Y \ar[r] & Y' = Y \amalg_X X'
}
$$
をとる。示すべきことは、
$$
(\mathcal{X} \times_{\mathcal{Y}} \mathcal{Z})_{Y'}
\longrightarrow
(\mathcal{X} \times_{\mathcal{Y}} \mathcal{Z})_Y
\times_{(\mathcal{X} \times_{\mathcal{Y}} \mathcal{Z})_X}
(\mathcal{X} \times_{\mathcal{Y}} \mathcal{Z})_{X'}
$$
が同値であることである。$2$-ファイバー積の定義を用いると、これは
\begin{equation}
\label{artin-equation-RS-fibre-product}
\mathcal{X}_{Y'} \times_{\mathcal{Y}_{Y'}} \mathcal{Z}_{Y'}
\longrightarrow
(\mathcal{X}_Y \times_{\mathcal{Y}_Y} \mathcal{Z}_Y)
\times_{(\mathcal{X}_X \times_{\mathcal{Y}_X} \mathcal{Z}_X)}
(\mathcal{X}_{X'} \times_{\mathcal{Y}_{X'}} \mathcal{Z}_{X'}).
\end{equation}
となる。関手
$$
\mathcal{X}_{Y'} \to \mathcal{X}_Y \times_{\mathcal{Y}_Y} \mathcal{Z}_Y,
\quad
\mathcal{Y}_{Y'} \to \mathcal{X}_X \times_{\mathcal{Y}_X} \mathcal{Z}_X,
\quad
\mathcal{Z}_{Y'} \to
\mathcal{X}_{X'} \times_{\mathcal{Y}_{X'}} \mathcal{Z}_{X'}
$$
がそれぞれ同値であることが与えられている。
(\ref{artin-equation-RS-fibre-product}) の右辺の対象は、系
$$
((x_Y, z_Y, \phi_Y), (x_{X'}, z_{X'}, \phi_{X'}), (\alpha, \beta)).
$$
である。このとき $(x_Y, x_{Y'}, \alpha)$ は
$\mathcal{X}_{Y'}$ の対象 $x_{Y'}$ の像と同型であり、
$(z_Y, z_{Y'}, \beta)$ は $\mathcal{Z}_{Y'}$ の対象 $z_{Y'}$ の像と
同型である。射の組 $(\phi_Y, \phi_{X'})$ は、$\mathcal{Y}_{Y'}$ における
$x_{Y'}$ と $z_{Y'}$ の像の間の射 $\psi$ に対応する。
したがって $(x_{Y'}, z_{Y'}, \psi)$ は
(\ref{artin-equation-RS-fibre-product}) の左辺の対象であり、与えられた右辺の
対象へ写る。これにより (\ref{artin-equation-RS-fibre-product}) が本質的全射で
あることが示された。充満忠実であることの証明は省略する。
\end{proof}




\section{変形圏}
\label{artin-section-deformation-categories}

\noindent
上で導入した記法を『形式的変形理論』の記法と対応させる。

\begin{lemma}
\label{artin-lemma-deformation-category}
$S$ を局所 Noether スキームとする。$\mathcal{X}$ を
$(\Sch/S)_{fppf}$ 上のグルーポイドをファイバーとする圏で、(RS) を
満たすものとする。$S$ 上有限型な任意の体 $k$ と、$k$ 上にある
$\mathcal{X}$ の任意の対象 $x_0$ に対して、前変形圏
$p : \mathcal{F}_{\mathcal{X}, k, x_0} \to \mathcal{C}_\Lambda$
（\ref{artin-equation-predeformation-category}）は変形圏である。『形式的変形理論』
の定義 \ref{formal-defos-definition-deformation-category} を参照せよ。
\end{lemma}

\begin{proof}
$\mathcal{F} = \mathcal{F}_{\mathcal{X}, k, x_0}$ とおく。
$f_1 : A_1 \to A$ と $f_2 : A_2 \to A$ を $\mathcal{C}_\Lambda$ における
環準同型で、$f_2$ は全射であるとする。関手
$$
\mathcal{F}(A_1 \times_A A_2)
\longrightarrow
\mathcal{F}(A_1) \times_{\mathcal{F}(A)} \mathcal{F}(A_2)
$$
が同値であることを示さなければならない。『形式的変形理論』の補題
\ref{formal-defos-lemma-RS-2-categorical} を参照せよ。
$X = \Spec(A)$, $X' = \Spec(A_2)$, $Y = \Spec(A_1)$,
$Y' = \Spec(A_1 \times_A A_2)$ とおく。スキームの圏で
$Y' = Y \amalg_X X'$ であることに注意せよ。『射についてさらに』の補題
\ref{more-morphisms-lemma-pushout-along-thickening} を参照せよ。
ファイバー圏の関手の図式
$$
\xymatrix{
\mathcal{X}_{Y'} \ar[r] \ar[d] &
\mathcal{X}_Y \times_{\mathcal{X}_X} \mathcal{X}_{X'} \ar[d] \\
\mathcal{X}_{\Spec(k)} \ar@{=}[r] & \mathcal{X}_{\Spec(k)}
}
$$
において、上の水平矢印は定義 \ref{artin-definition-RS} により同値である。
$\mathcal{F}(B)$ は、$k$ 上で $x_0$ との同一視を備えた
$\mathcal{X}_{\Spec(B)}$ の対象の圏なので、主張が従う。
\end{proof}

\begin{remark}
\label{artin-remark-deformation-category-implies}
$S$ を局所 Noether スキームとし、$\mathcal{X}$ を
$(\Sch/S)_{fppf}$ 上のグルーポイドをファイバーとする圏とする。
$k$ を $S$ 上有限型な体、$x_0$ を $k$ 上の $\mathcal{X}$ の対象とする。
$p : \mathcal{F} \to \mathcal{C}_\Lambda$ を
(\ref{artin-equation-predeformation-category}) のものとする。
$\mathcal{F}$ が変形圏、すなわち $\mathcal{F}$ が
Rim--Schlessinger 条件 (RS) を満たすならば、『形式的変形理論』の補題
\ref{formal-defos-lemma-RS-implies-S1-S2} により、$\mathcal{F}$ は
Schlessinger の条件 (S1) と (S2) を満たす。
$\overline{\mathcal{F}}$ を同型類の関手とする。『形式的変形理論』の注意
\ref{formal-defos-remarks-cofibered-groupoids}
（\ref{formal-defos-item-associated-functor-isomorphism-classes}）を参照せよ。
このとき $\overline{\mathcal{F}}$ も (S1) と (S2) を満たす。
『形式的変形理論』の補題
\ref{formal-defos-lemma-S1-S2-associated-functor} を参照せよ。
特に、これは補題 \ref{artin-lemma-deformation-category} の状況で成り立つ。
\end{remark}




\section{体の変更}
\label{artin-section-change-of-field}

\noindent
本節は『形式的変形理論』の節
\ref{formal-defos-section-change-of-field} に対応する。
そこで指摘したように、体の変更のもとで何が起こるかを論じるには、
$\mathcal{C}_\Lambda$ の代わりに $\mathcal{C}_{\Lambda, k}$ と書く必要がある。
次の補題では、『形式的変形理論』の状況
\ref{formal-defos-situation-change-of-fields} で導入した記法
$\mathcal{F}_{l/k}$ を用いる。

\begin{lemma}
\label{artin-lemma-change-of-field}
$S$ を局所 Noether スキームとする。$\mathcal{X}$ を
$(\Sch/S)_{fppf}$ 上のグルーポイドをファイバーとする圏とする。
$k$ を $S$ 上有限型な体とし、$l/k$ を有限拡大とする。
$x_0$ を $\Spec(k)$ 上にある $\mathcal{F}$ の対象とする。
$x_{l, 0}$ で $x_0$ の $\Spec(l)$ への制限を表す。このとき
$\mathcal{C}_{\Lambda, l}$ 上のグルーポイドをファイバーとする余ファイバー圏の
標準的な関手
$$
(\mathcal{F}_{\mathcal{X}, k , x_0})_{l/k}
\longrightarrow
\mathcal{F}_{\mathcal{X}, l, x_{l, 0}}
$$
が存在する。$\mathcal{X}$ が (RS) を満たすならば、この関手は同値である。
\end{lemma}

\begin{proof}
(\ref{artin-equation-factor}) に現れる形の分解
$$
\Spec(l) \to \Spec(B) \to S
$$
を考える。定義により
$$
(\mathcal{F}_{\mathcal{X}, k , x_0})_{l/k}(B) =
\mathcal{F}_{\mathcal{X}, k, x_0}(B \times_l k)
$$
である。『形式的変形理論』の状況
\ref{formal-defos-situation-change-of-fields} を参照せよ。
したがって、その対象は、射
$\Spec(k) \to \Spec(B \times_l k)$ 上にある $\mathcal{X}$ の射
$x_0 \to x$ である。$\mathcal{X}$ の引き戻し関手を選ぶと、
$x_0 \to x$ に射 $x_{l, 0} \to x_B$ を対応させられる。ここで $x_B$ は、
$B \times_l k \subset B$ から得られる射
$\Spec(B) \to \Spec(B \times_l k)$ を介した $x$ の $\Spec(B)$ への制限である。
この構成は $B$ について関手的で、射と両立する。

\medskip\noindent
次に、$\mathcal{X}$ が (RS) を満たすと仮定する。図式
$$
\vcenter{
\xymatrix{
l & B \ar[l] \\
k \ar[u] & B \times_l k \ar[l] \ar[u]
}
}
\quad\text{and}\quad
\vcenter{
\xymatrix{
\Spec(l) \ar[d] \ar[r] & \Spec(B) \ar[d] \\
\Spec(k) \ar[r] & \Spec(B \times_l k)
}
}
$$
を考える。左の図式は環のファイバー積である。右の図式はスキームの圏に
おける押し出しである。『射についてさらに』の補題
\ref{more-morphisms-lemma-pushout-along-thickening} を参照せよ。
これらのスキームはいずれも $S$ 上有限型である（定義
\ref{artin-definition-RS} に続く注意を参照せよ）。したがって (RS) により、
ファイバー圏の同値
$$
\mathcal{X}_{\Spec(B \times_l k)}
\longrightarrow
\mathcal{X}_{\Spec(k)}
\times_{\mathcal{X}_{\Spec(l)}}
\mathcal{X}_{\Spec(B)}
$$
が得られる。これは、上で定義した関手がファイバー圏の同値を与えることを
意味する。したがって『圏』の補題
\ref{categories-lemma-equivalence-fibred-categories} の双対により、この関手は
グルーポイドをファイバーとする余ファイバー圏の同値である。
\end{proof}








\section{接空間}
\label{artin-section-tangent-spaces}

\noindent
$S$ を局所 Noether スキームとする。$\mathcal{X}$ を
$(\Sch/S)_{fppf}$ 上のグルーポイドをファイバーとする圏とする。
$k$ を $S$ 上有限型な体とし、$x_0$ を $k$ 上の $\mathcal{X}$ の対象とする。
『形式的変形理論』の節 \ref{formal-defos-section-tangent-spaces} では、
前変形圏 $\mathcal{F}_{\mathcal{X}, k, x_0}$ の {\it 接空間}
\begin{equation}
\label{artin-equation-tangent-space}
T\mathcal{F}_{\mathcal{X}, k, x_0} =
\left\{
\begin{matrix}
\text{isomorphism classes of morphisms}\\
x_0 \to x\text{ over }\Spec(k) \to \Spec(k[\epsilon])
\end{matrix}
\right\}
\end{equation}
を定義した。『形式的変形理論』の節
\ref{formal-defos-section-infinitesimal-automorphisms} では、さらに
\begin{equation}
\label{artin-equation-infinitesimal-automorphisms}
\text{Inf}(\mathcal{F}_{\mathcal{X}, k, x_0}) =
\Ker\left(
\text{Aut}_{\Spec(k[\epsilon])}(x'_0) \to \text{Aut}_{\Spec(k)}(x_0)
\right)
\end{equation}
を定義した。ここで $x_0'$ は $x_0$ の $\Spec(k[\epsilon])$ への引き戻しである。
$\mathcal{X}$ が Rim--Schlessinger 条件 (RS) を満たすならば、
『形式的変形理論』の補題
\ref{formal-defos-lemma-tangent-space-vector-space} により、
$T\mathcal{F}_{\mathcal{X}, k, x_0}$ は自然な $k$-ベクトル空間の構造をもつ
（仮定が成り立つことは補題 \ref{artin-lemma-deformation-category} と注意
\ref{artin-remark-deformation-category-implies} による）。さらに『形式的変形理論』の
補題 \ref{formal-defos-lemma-infaut-vector-space} によれば、
$\text{Inf}(\mathcal{F}_{\mathcal{X}, k, x_0})$ は、加法が自己同型の合成と
一致する自然な $k$-ベクトル空間の構造をもつ。これらのベクトル空間が
有限次元であることを要求するのが自然な条件である。

\medskip\noindent
次の補題は、$\mathcal{X}$ が $S$ 上局所有限型ならば、この条件が成り立つことを
述べる（『スタックの射』の節
\ref{stacks-morphisms-section-finite-type} を参照せよ）。

\begin{lemma}
\label{artin-lemma-finite-dimension}
$S$ を局所 Noether スキームとする。次を仮定する。
\begin{enumerate}
\item $\mathcal{X}$ は代数スタックである。
\item $U$ は $S$ 上局所有限型なスキームである。
\item $(\Sch/U)_{fppf} \to \mathcal{X}$ は滑らかで全射な射である。
\end{enumerate}
このとき、節 \ref{artin-section-predeformation-categories} に現れる任意の
$\mathcal{F} = \mathcal{F}_{\mathcal{X}, k, x_0}$ に対して、接空間
$T\mathcal{F}$ と無限小自己同型空間 $\text{Inf}(\mathcal{F})$ は
$k$ 上有限次元である。
\end{lemma}

\begin{proof}
$\mathcal{U} = (\Sch/U)_{fppf}$ と書く。代数スタックの定義により、
$1$-射 $\mathcal{U} \to \mathcal{X}$ は代数空間により表現可能である。
したがって、特に 2-ファイバー積
$$
\mathcal{U}_{x_0} = (\Sch/\Spec(k))_{fppf} \times_\mathcal{X} \mathcal{U}
$$
は $\Spec(k)$ 上の代数空間 $U_{x_0}$ により表現される。このとき
$U_{x_0} \to \Spec(k)$ は滑らかで全射である（特に $U_{x_0}$ は空でない）。
『体上の空間』の補題
\ref{spaces-over-fields-lemma-smooth-separable-closed-points-dense} により、
有限拡大 $l/k$ と $k$ 上の点 $\Spec(l) \to U_{x_0}$ をとれる。
補題 \ref{artin-lemma-change-of-field} と、$\mathcal{X}$ が (RS) を満たすことにより、
$$
(\mathcal{F}_{\mathcal{X}, k , x_0})_{l/k} =
\mathcal{F}_{\mathcal{X}, l, x_{l, 0}}
$$
である。したがって、『形式的変形理論』の補題
\ref{formal-defos-lemma-tangent-space-change-of-field} および
\ref{formal-defos-lemma-inf-aut-change-of-field} により、
$$
T\mathcal{F} \otimes_k l \cong T\mathcal{F}_{\mathcal{X}, l, x_{l, 0}}
\quad\text{and}\quad
\text{Inf}(\mathcal{F}) \otimes_k l \cong
\text{Inf}(\mathcal{F}_{\mathcal{X}, l, x_{l, 0}})
$$
であることが分かる（これらを適用できることは、補題
\ref{artin-lemma-algebraic-stack-RS} と \ref{artin-lemma-deformation-category}、および注意
\ref{artin-remark-deformation-category-implies} による）。ゆえに、
$T\mathcal{F}_{\mathcal{X}, l, x_{l, 0}}$ と
$\text{Inf}(\mathcal{F}_{\mathcal{X}, l, x_{l, 0}})$ が $l$ 上有限次元であることを
示せば十分である。$x_{l, 0}$ は $l$ 上の $\mathcal{U}$ の点 $u_0$ から
得られることに注意せよ。

\medskip\noindent
議論の流れをいったん中断して、無限小自己同型に関する主張が接空間に関する
主張から従うことを示す。すなわち、
$\mathcal{R} = \mathcal{U} \times_\mathcal{X} \mathcal{U}$ とおく。
$r_0$ を $\mathcal{R}$ の $l$-値点
$(u_0, u_0, \text{id}_{x_0})$ とする。補題
\ref{artin-lemma-fibre-product-deformation-categories} と『形式的変形理論』の補題
\ref{formal-defos-lemma-deformation-functor-diagonal} を合わせると、
$$
\text{Inf}(\mathcal{F}_{\mathcal{X}, l, x_{l, 0}})
\subset
T\mathcal{F}_{\mathcal{R}, l, r_0}
$$
を得る。$\mathcal{R}$ は代数スタックであることに注意せよ。『代数スタック』の
補題 \ref{algebraic-lemma-2-fibre-product-general} を参照せよ。
また、$\mathcal{R}$ は $U$ 上滑らかな代数空間 $R$ により表現される
（いずれの射影を用いてもよい。『代数スタック』の補題
\ref{algebraic-lemma-stack-presentation} を参照せよ）。したがって、スキーム
$U'$ と全射エタール射 $U' \to R$ を選べば、$U'$ は $U$ 上滑らかであり、
したがって $S$ 上局所有限型である。$(\Sch/U')_{fppf} \to \mathcal{R}$ は
全射かつ滑らかなので、問題は接空間の場合に帰着された。

\medskip\noindent
関手 (\ref{artin-equation-functoriality})
$$
\mathcal{F}_{\mathcal{U}, l, u_0}
\longrightarrow
\mathcal{F}_{\mathcal{X}, l, x_{l, 0}}
$$
は補題 \ref{artin-lemma-formally-smooth-on-deformation-categories} により滑らかである。
誘導される接空間の写像
$$
T\mathcal{F}_{\mathcal{U}, l, u_0}
\longrightarrow
T\mathcal{F}_{\mathcal{X}, l, x_{l, 0}}
$$
は $l$-線形であり（『形式的変形理論』の補題
\ref{formal-defos-lemma-k-linear-differential} による）、全射である
（前変形圏の滑らかな写像が接空間上の全射を誘導することは、
『形式的変形理論』の補題
\ref{formal-defos-lemma-smooth-morphism-essentially-surjective} による）。
ゆえに、点 $u_0$ における、表現可能な代数スタック $\mathcal{U}$ に付随する
変形空間の接空間が有限次元であることを示せば十分である。
$u_0 : \Spec(l) \to U$ が $\Spec(R)$ を経由し、かつ
$\Spec(R) \to S$ が $\Spec(\Lambda) \subset S$ を経由するように
$\Spec(R) \subset U$ を選ぶ。$u_0$ に対応する $\Lambda$-代数準同型
$\varphi_0 : R \to l$ の核を $\mathfrak m_R \subset R$ とする。
$R$ は Noether 環 $\Lambda$ 上有限型なので Noether 環である。したがって
$\mathfrak m_R = (f_1, \ldots, f_n)$ は有限生成イデアルである。さらに
$$
T\mathcal{F}_{\mathcal{U}, l, u_0}
=
\{\varphi : R \to l[\epsilon] \mid
\varphi \text{ is a } \Lambda\text{-algebra map and }
\varphi \bmod \epsilon = \varphi_0\}
$$
である。右辺の元は $f_1, \ldots, f_n$ における値によって決まるので、
次元は高々 $n$ であり、主張が従う。いくつかの詳細は省略する。
\end{proof}

\begin{lemma}
\label{artin-lemma-fibre-product-tangent-spaces}
$S$ を局所 Noether スキームとする。$p : \mathcal{X} \to \mathcal{Y}$ と
$q : \mathcal{Z} \to \mathcal{Y}$ を、$(\Sch/S)_{fppf}$ 上の
グルーポイドをファイバーとする圏の $1$-射とする。
$\mathcal{X}$, $\mathcal{Y}$, $\mathcal{Z}$ が (RS) を満たすと仮定する。
$k$ を $S$ 上有限型な体とし、$w_0$ を $k$ 上の
$\mathcal{W} = \mathcal{X} \times_\mathcal{Y} \mathcal{Z}$ の対象とする。
$w_0$ から得られる $\mathcal{X}, \mathcal{Y}, \mathcal{Z}$ の対象を
$x_0, y_0, z_0$ と表す。このとき、$k$-ベクトル空間の $6$-項完全列
$$
\xymatrix{
0 \ar[r] &
\text{Inf}(\mathcal{F}_{\mathcal{W}, k, w_0}) \ar[r] &
\text{Inf}(\mathcal{F}_{\mathcal{X}, k, x_0}) \oplus
\text{Inf}(\mathcal{F}_{\mathcal{Z}, k, z_0}) \ar[r] &
\text{Inf}(\mathcal{F}_{\mathcal{Y}, k, y_0}) \ar[lld] \\
 &
T\mathcal{F}_{\mathcal{W}, k, w_0} \ar[r] &
T\mathcal{F}_{\mathcal{X}, k, x_0} \oplus
T\mathcal{F}_{\mathcal{Z}, k, z_0} \ar[r] &
T\mathcal{F}_{\mathcal{Y}, k, y_0}
}
$$
が存在する。
\end{lemma}

\begin{proof}
補題 \ref{artin-lemma-fibre-product-RS} により、$\mathcal{W}$ は (RS) を満たすので、
補題の主張は意味をもつ。主張を示すには、補題
\ref{artin-lemma-fibre-product-deformation-categories} と
\ref{artin-lemma-deformation-category}、および『形式的変形理論』の補題
\ref{formal-defos-lemma-deformation-categories-fiber-product-morphisms} を適用する。
\end{proof}






\section{形式対象}
\label{artin-section-formal-objects}

\noindent
本節では、『形式的変形理論』で既に定義したいくつかの概念を現在の状況へ
移す。以下では、$R$ がスキームの射 $\Spec(R) \to S$ を備えた環であることを
「$R$ は $S$-代数である」と表す。

\begin{definition}
\label{artin-definition-formal-objects}
$S$ を局所 Noether スキームとし、
$p : \mathcal{X} \to (\Sch/S)_{fppf}$ をグルーポイドをファイバーとする圏とする。
\begin{enumerate}
\item $\mathcal{X}$ の {\it 形式対象} $\xi = (R, \xi_n, f_n)$ とは、
Noether 完備局所 $S$-代数 $R$、$\Spec(R/\mathfrak m_R^n)$ 上にある
$\mathcal{X}$ の対象 $\xi_n$、および
$\Spec(R/\mathfrak m^n) \to \Spec(R/\mathfrak m^{n + 1})$ 上にある
$\mathcal{X}$ の射 $f_n : \xi_n \to \xi_{n + 1}$ からなり、
$R/\mathfrak m$ が $S$ 上有限型な体であるものをいう。
\item {\it 形式対象の射}
$a : \xi = (R, \xi_n, f_n) \to \eta = (T, \eta_n, g_n)$ とは、
各 $n$ に対して図式
$$
\xymatrix{
\xi_n \ar[r]_{f_n} \ar[d]_{a_n} & \xi_{n + 1} \ar[d]^{a_{n + 1}} \\
\eta_n \ar[r]^{g_n} & \eta_{n + 1}
}
$$
が可換となるような射 $a_n : \xi_n \to \eta_n$ により与えられる。
関手 $p$ を適用すると、両立する射の族
$\Spec(R/\mathfrak m_R^n) \to \Spec(T/\mathfrak m_T^n)$ が得られ、したがって
$S$ 上の射 $a_0 : \Spec(R) \to \Spec(T)$ が得られる。$a$ は
$a_0$ の {\it 上にある} という。
\end{enumerate}
\end{definition}

\noindent
このようにして $\mathcal{X}$ の形式対象の圏を得る。

\begin{remark}
\label{artin-remark-formal-objects-match}
$S$ を局所 Noether スキームとし、
$p : \mathcal{X} \to (\Sch/S)_{fppf}$ をグルーポイドをファイバーとする圏とする。
$\xi = (R, \xi_n, f_n)$ を形式対象とする。$k = R/\mathfrak m$ および
$x_0 = \xi_1$ とおく。形式対象 $\xi$ は、前変形圏
$\mathcal{F}_{\mathcal{X}, k, x_0}$ の形式対象 $\xi$ を定める。
これは、上の定義 \ref{artin-definition-formal-objects}、『形式的変形理論』の定義
\ref{formal-defos-definition-formal-objects}、および節
\ref{artin-section-predeformation-categories} における前変形圏
$\mathcal{F}_{\mathcal{X}, k, x_0}$ の構成から直ちに従う。
\end{remark}

\noindent
$F : \mathcal{X} \to \mathcal{Y}$ が $(\Sch/S)_{fppf}$ 上の
グルーポイドをファイバーとする圏の $1$-射ならば、$F$ は形式対象の圏の
間の関手も誘導する。

\begin{lemma}
\label{artin-lemma-smooth-lift-formal}
$S$ を局所 Noether スキームとする。$F : \mathcal{X} \to \mathcal{Y}$ を
$(\Sch/S)_{fppf}$ 上のグルーポイドをファイバーとする圏の $1$-射とする。
$\eta = (R, \eta_n, g_n)$ を $\mathcal{Y}$ の形式対象とし、$\xi_1$ を
$F(\xi_1) \cong \eta_1$ を満たす $\mathcal{X}$ の対象とする。
$F$ が対象について形式的に滑らかならば（『表現可能性の判定条件』の節
\ref{criteria-section-formally-smooth} を参照せよ）、$F(\xi) \cong \eta$ を
満たす $\mathcal{X}$ の形式対象 $\xi = (R, \xi_n, f_n)$ が存在する。
\end{lemma}

\begin{proof}
各射 $\Spec(R/\mathfrak m^n) \to \Spec(R/\mathfrak m^{n + 1})$ は、
$S$ 上のアフィンスキームの一階無限小厚化であることに注意せよ。
したがって $F$ に関する仮定により、$\xi_1$ を $\mathcal{X}$ の対象
$\xi_2, \xi_3, \ldots$ へ逐次持ち上げ、両立する同型
$\eta_n|_{\Spec(R/\mathfrak m^{n - 1})} \cong \eta_{n - 1}$ および
$F(\eta_n) \cong \xi_n$ を備えさせることができる。
\end{proof}

\noindent
$S$ を局所 Noether スキームとし、
$p : \mathcal{X} \to (\Sch/S)_{fppf}$ をグルーポイドをファイバーとする圏とする。
$x$ を $R$ 上の $\mathcal{X}$ の対象と仮定する。ここで $R$ は、$S$ 上有限型な
剰余体をもつ Noether 完備局所 $S$-代数である。このとき、制限の系
$\xi_n = x|_{\Spec(R/\mathfrak m^n)}$ と、制限の推移性から得られる自然な射
$\xi_1 \to \xi_2 \to \ldots$ を考えられる。したがって
$\xi = (R, \xi_n, \xi_n \to \xi_{n + 1})$ は $\mathcal{X}$ の形式対象である。
この構成は対象 $x$ について関手的である。ゆえに関手
\begin{equation}
\label{artin-equation-approximation}
\left\{
\begin{matrix}
\text{objects }x\text{ of }\mathcal{X} \text{ such that }p(x) = \Spec(R) \\
\text{where }R\text{ is Noetherian complete local}\\
\text{with }R/\mathfrak m\text{ of finite type over }S
\end{matrix}
\right\}
\longrightarrow
\left\{
\begin{matrix}
\text{formal objects of }\mathcal{X}
\end{matrix}
\right\}
\end{equation}
を得る。正確には、左辺は表示された対象からなる $\mathcal{X}$ の充満部分圏で、
右辺は定義 \ref{artin-definition-formal-objects} における $\mathcal{X}$ の形式対象の
圏である。

\begin{definition}
\label{artin-definition-effective}
$S$ を局所 Noether スキームとする。$\mathcal{X}$ を
$(\Sch/S)_{fppf}$ 上のグルーポイドをファイバーとする圏とする。
$\mathcal{X}$ の形式対象 $\xi = (R, \xi_n, f_n)$ が関手
(\ref{artin-equation-approximation}) の本質像に属するとき、これは
{\it 有効} であるという。
\end{definition}

\noindent
グルーポイドをファイバーとする圏が代数スタックならば、次の補題により
すべての形式対象は有効である。

\begin{lemma}
\label{artin-lemma-effective}
$S$ を局所 Noether スキームとし、$\mathcal{X}$ を $S$ 上の代数スタックとする。
関手 (\ref{artin-equation-approximation}) は同値である。
\end{lemma}

\begin{proof}
場合 I：$\mathcal{X}$ が（スキームにより）表現可能である場合。
$S$ 上のあるスキーム $X$ に対して
$\mathcal{X} = (\Sch/X)_{fppf}$ であるとする。定義を展開すると、次を示せばよい。
$R/\mathfrak m$ が $S$ 上有限型である Noether 完備局所 $S$-代数 $R$ に対して、
$$
\Mor_S(\Spec(R), X) \longrightarrow \lim \Mor_S(\Spec(R/\mathfrak m^n), X)
$$
は全単射である。これは『形式空間』の補題
\ref{formal-spaces-lemma-map-into-scheme} から従う。

\medskip\noindent
場合 II：$\mathcal{X}$ が代数空間により表現可能である場合。
$\mathcal{X}$ は $X$ により表現されるとする。この場合も、上のような $R$ に
対して
$$
\Mor_S(\Spec(R), X) \longrightarrow \lim \Mor_S(\Spec(R/\mathfrak m^n), X)
$$
が全単射であることを示さなければならない。これは『形式空間』の補題
\ref{formal-spaces-lemma-map-into-algebraic-space} である。

\medskip\noindent
場合 III：代数スタックの一般の場合。一般的な注意として、
(\ref{artin-equation-approximation}) の左辺と右辺は、$S$ 上のアフィンスキームのうち、
$S$ 上有限型な剰余体をもつ Noether 完備局所環のスペクトルであるものの圏上で、
グルーポイドをファイバーとする圏である。以下の証明では、それらがこの圏上の
ある位相に関するスタックをなすことも分かる。

\medskip\noindent
まず充満忠実性を証明する。$R$ を $k = R/\mathfrak m$ が $S$ 上有限型である
Noether 完備局所 $S$-代数とする。$x, x'$ を $R$ 上の $\mathcal{X}$ の対象とする。
$\mathcal{X}$ は代数スタックなので、$\mathit{Isom}(x, x')$ は
$\Spec(R)$ 上の代数空間 $I$ により表現される。『代数スタック』の補題
\ref{algebraic-lemma-representable-diagonal} を参照せよ。
$\Spec(R)$ 上の $I$ に場合 II を適用すると、
(\ref{artin-equation-approximation}) が $\Spec(R)$ 上のファイバー圏で充満忠実であることが
直ちに従う。したがって『圏』の補題
\ref{categories-lemma-equivalence-fibred-categories} により、この関手は充満忠実である。

\medskip\noindent
本質的全射性を示す。$\xi = (R, \xi_n, f_n)$ を $\mathcal{X}$ の形式対象とする。
$S$ 上のスキーム $U$ と、全射かつ滑らかな射
$f : (\Sch/U)_{fppf} \to \mathcal{X}$ を選ぶ。各 $n$ に対してファイバー積
$$
(\Sch/\Spec(R/\mathfrak m^n))_{fppf}
\times_{\xi_n, \mathcal{X}, f}
(\Sch/U)_{fppf}
$$
を考える。仮定により、これは $\Spec(R/\mathfrak m^n)$ 上全射かつ滑らかな
代数空間 $V_n$ により表現される。射 $f_n : \xi_n \to \xi_{n + 1}$ は、
代数空間の Cartesian 図式
$$
\xymatrix{
V_{n + 1} \ar[d] & V_n \ar[d] \ar[l] \\
\Spec(R/\mathfrak m^{n + 1}) & \Spec(R/\mathfrak m^n) \ar[l]
}
$$
を誘導する。『体上の空間』の補題
\ref{spaces-over-fields-lemma-smooth-separable-closed-points-dense} により、
有限分離拡大 $k'/k$ と $k$ 上の点 $v'_1 : \Spec(k') \to V_1$ をとれる。
$R \subset R'$ を、剰余体拡大が $k'/k$ である有限エタール拡大とする
（これは『代数』の補題 \ref{algebra-lemma-henselian-cat-finite-etale} および
\ref{algebra-lemma-complete-henselian} により存在し、一意である）。
$n = 2, 3, 4, \ldots$ に対して $V_n \to \Spec(R/\mathfrak m^n)$ に
滑らかさの無限小持ち上げ判定法（『空間の射についてさらに』の補題
\ref{spaces-more-morphisms-lemma-smooth-formally-smooth} を参照せよ）を適用すると、
$\Spec(R/\mathfrak m^n)$ 上の射
$v'_n : \Spec(R'/(\mathfrak m')^n) \to V_n$ で、可換図式
$$
\xymatrix{
\Spec(R'/(\mathfrak m')^{n + 1}) \ar[d]_{v'_{n + 1}} &
\Spec(R'/(\mathfrak m')^n) \ar[d]^{v'_n} \ar[l] \\
V_{n + 1} & V_n \ar[l]
}
$$
に入るものを逐次とれる。射影 $V_n \to U$ と合成して、両立する射の系
$u'_n : \Spec(R'/(\mathfrak m')^n) \to U$ を得る。場合 I により、族
$(u'_n)$ は一意な射 $u' : \Spec(R') \to U$ から得られる。
$u'$ に $1$-射 $f$ を適用して得られる $\Spec(R')$ 上の
$\mathcal{X}$ の対象を $x'$ と表す。構成により、$\Spec(R') \to \Spec(R)$ 上にある
形式対象の射
$$
(\ref{artin-equation-approximation})(x') =
(R', x'|_{\Spec(R'/(\mathfrak m')^n)}, \ldots)
\longrightarrow
(R, \xi_n, f_n)
$$
が存在する。$R' \otimes_R R'$ は、この議論を適用できる Noether 完備局所環の
スペクトルの有限直積であることに注意せよ。二つの射影を
$p_0, p_1 : \Spec(R' \otimes_R R') \to \Spec(R')$ と表す。
上で示した充満忠実性により、標準的な同型
$\varphi : p_0^*x' \to p_1^*x'$ が存在する。実際、
$\Spec((R' \otimes_R R')/\mathfrak m^n(R' \otimes_R R'))$ 上ではそのような
同型が存在する。$\varphi$ が余サイクル条件を満たすことの証明は省略する
（『スタック』の定義 \ref{stacks-definition-descent-data} を参照せよ）。
$\{\Spec(R') \to \Spec(R)\}$ は fppf 被覆なので、$x'$ は $\Spec(R)$ 上の
$\mathcal{X}$ の対象 $x$ へ降下する。$x_n$ が $x$ の
$\Spec(R/\mathfrak m^n)$ への制限であることの証明は省略する。
\end{proof}

\begin{lemma}
\label{artin-lemma-fibre-product-effective}
$S$ をスキームとする。$p : \mathcal{X} \to \mathcal{Y}$ と
$q : \mathcal{Z} \to \mathcal{Y}$ を、$(\Sch/S)_{fppf}$ 上の
グルーポイドをファイバーとする圏の $1$-射とする。
関手 (\ref{artin-equation-approximation}) が $\mathcal{X}$, $\mathcal{Y}$,
$\mathcal{Z}$ に対して同値ならば、$\mathcal{X} \times_\mathcal{Y} \mathcal{Z}$ に
対しても同値である。
\end{lemma}

\begin{proof}
$\mathcal{X} \times_\mathcal{Y} \mathcal{Z}$ に対する
(\ref{artin-equation-approximation}) の左辺と右辺は、それぞれ単に、$\mathcal{Y}$ 上での
$\mathcal{X}$, $\mathcal{Z}$ に対する (\ref{artin-equation-approximation}) の左辺と右辺の
$2$-ファイバー積である。圏同値は $2$-ファイバー積をとることと両立するので、
主張が従う。『圏』の補題
\ref{categories-lemma-equivalence-2-fibre-product} を参照せよ。
\end{proof}






\section{近似}
\label{artin-section-approximation}

\noindent
Michael Artin の基本的洞察は、極限を保つスタックの対象を近似できるという
ことである。すなわち、Noether 完備局所環上のスタックの対象 $x$ が与えられると、
$x$ に「近い」代数的な環上の対象 $x_A$ を見つけることができる。ここで代数的な
環とは有限型 $S$-代数を意味し、近いとは進位的に近いことを意味する。本節では、
これを単純だが一般的な形で提示する。

\medskip\noindent
結果を定式化するため、Stacks project のさまざまな箇所にある定義をまとめる
必要がある。まず『表現可能性の判定条件』の節
\ref{criteria-section-limit-preserving} では、スキームの圏上の
グルーポイドをファイバーとする圏の $1$-射について
{\it 対象上で極限を保つ} ことを導入した。『代数についてさらに』の定義
\ref{more-algebra-definition-G-ring} では {\it G-環} の概念を定義した。
$S$ を局所 Noether スキームとし、$A$ を $S$-代数とする。
$\Spec(A) \to S$ が有限型であるとき、$A$ は {\it $S$ 上有限型}、または
{\it 有限型 $S$-代数} であるという。この場合、$A$ は Noether 環である。
最後に、環 $A$ とイデアル $I$ に対して
$\text{Gr}_I(A) = \bigoplus I^n/I^{n + 1}$ と表す。

\begin{lemma}
\label{artin-lemma-approximate}
$S$ を局所 Noether スキームとする。
$p : \mathcal{X} \to (\Sch/S)_{fppf}$ をグルーポイドをファイバーとする圏とする。
$x$ を $\Spec(R)$ 上にある $\mathcal{X}$ の対象とする。ここで $R$ は、
$S$ 上有限型な剰余体 $k$ をもつ Noether 完備局所環である。
$s \in S$ を $\Spec(k) \to S$ の像とする。(a) $\mathcal{O}_{S, s}$ が
G-環であり、(b) $p$ が対象上で極限を保つと仮定する。このとき、任意の整数
$N \geq 1$ に対して、次が存在する。
\begin{enumerate}
\item 有限型 $S$-代数 $A$、
\item 極大イデアル $\mathfrak m_A \subset A$、
\item $\Spec(A)$ 上の $\mathcal{X}$ の対象 $x_A$、
\item $S$-同型 $R/\mathfrak m_R^N \cong A/\mathfrak m_A^N$、
\item (4) と両立する同型
$x|_{\Spec(R/\mathfrak m_R^N)} \cong x_A|_{\Spec(A/\mathfrak m_A^N)}$、および
\item 次数付き $k$-代数の同型
$\text{Gr}_{\mathfrak m_R}(R) \cong \text{Gr}_{\mathfrak m_A}(A)$。
\end{enumerate}
\end{lemma}

\begin{proof}
$k$ が有限 $\Lambda$-代数となるようなアフィン開集合
$\Spec(\Lambda) \subset S$ を選ぶ。『射』の補題
\ref{morphisms-lemma-point-finite-type} を参照せよ。$S$ を
$\Spec(\Lambda)$ に置き換えてよく、実際そのようにする。

\medskip\noindent
$R$ を有向余極限 $R = \colim C_j$ と書ける。ここで各 $C_j$ は有限型
$\Lambda$-代数である（『代数』の補題
\ref{algebra-lemma-ring-colimit-fp} を参照せよ）。仮定 (b) により、対象 $x$ は
ある $C_j$ 上の対象の制限と同型である。したがって、有限型 $\Lambda$-代数 $C$、
$\Lambda$-代数準同型 $C \to R$、および $\Spec(C)$ 上の $\mathcal{X}$ の対象
$x_C$ で $x = x_C|_{\Spec(R)}$ を満たすものを選べる。
$C$ の選択は記帳のための工夫であり、省くこともできる。後で用いるため、
$C = \Lambda[y_1, \ldots, y_u]/(f_1, \ldots, f_v)$ と書き、
$C \to R$ による $y_i$ の像を $\overline{a}_i \in R$ と表す。
$\mathfrak m_C = C \cap \mathfrak m_R$ とおく。

\medskip\noindent
$\Lambda$-代数の全射 $\Lambda[x_1, \ldots, x_s] \to k$ を選び、その核を
$\mathfrak m'$ と表す。多項式環の普遍性により、これを $\Lambda$-代数準同型
$\Lambda[x_1, \ldots, x_s] \to R$ へ持ち上げられる。
$\mathfrak m_R$ の生成元へ写る変数をいくつか追加する、すなわち $s$ を少し
増やす。こうすると $\Lambda[x_1, \ldots, x_s] \to R/\mathfrak m_R^2$ は全射に
なる。すると、たとえば『形式的変形理論』の補題
\ref{formal-defos-lemma-surjective-cotangent-space} により、
\begin{equation}
\label{artin-equation-surjection}
P = \Lambda[x_1, \ldots, x_s]_{\mathfrak m'}^\wedge \longrightarrow R
\end{equation}
は Noether 完備局所環の全射である。

\medskip\noindent
上で得た $\overline{a}_i$ の持ち上げ $a_i \in P$ を選ぶ。
(\ref{artin-equation-surjection}) の核の生成元 $b_1, \ldots, b_r \in P$ を選ぶ。
これまでの選択により可能であるように、$P$ において
$$
f_j(a_1, \ldots, a_u) = \sum c_{ji} b_i
$$
を満たす $c_{ji} \in P$ を選ぶ。生成元
$$
k_1, \ldots, k_t \in
\Ker(P^{\oplus r} \xrightarrow{(b_1, \ldots, b_r)} P)
$$
を選び、$k_i = (k_{i1}, \ldots, k_{ir})$ および $K = (k_{ij})$ と書く。
すると
$$
P^{\oplus t} \xrightarrow{K}
P^{\oplus r} \xrightarrow{(b_1, \ldots, b_r)}
P \to R \to 0
$$
は $P$-加群の完全列となる。特に $\sum k_{ij} b_j = 0$ である。
必要なら $N$ を増やして、この完全列の最初の二つの写像について
Artin--Rees の補題で $N - 1$ が使えると仮定してよい（用語については
『代数についてさらに』の節 \ref{more-algebra-section-artin-rees} を参照せよ）。

\medskip\noindent
仮定により $\mathcal{O}_{S, s} = \Lambda_{\Lambda \cap \mathfrak m'}$ は G-環である。
したがって『代数についてさらに』の命題
\ref{more-algebra-proposition-finite-type-over-G-ring} により、環
$\Lambda[x_1, \ldots, x_s]_{\mathfrak m'}$ は $G$-環である。
ゆえに『環準同型の平滑化』の定理
\ref{smoothing-theorem-approximation-property-variant} により、エタール環準同型
$$
\Lambda[x_1, \ldots, x_s]_{\mathfrak m'} \to B,
$$
$\mathfrak m'$ 上にある $B$ の極大イデアル $\mathfrak m_B$、および元
$a'_i, b'_i, c'_{ij}, k'_{ij} \in B'$ で、次を満たすものが存在する。
\begin{enumerate}
\item $\kappa(\mathfrak m') = \kappa(\mathfrak m_B)$。これは
$\Lambda[x_1, \ldots, x_s]_{\mathfrak m'} \subset B_{\mathfrak m_B}
\subset P$ であり、$P$ が $\mathfrak m_B$ における $B$ の完備化と同一視される
ことを意味する。『環準同型の平滑化』の定理
\ref{smoothing-theorem-approximation-property-variant} の直前の注意を参照せよ。
\item $a_i - a'_i, b_i - b'_i, c_{ij} - c'_{ij}, k_{ij} - k'_{ij} \in
(\mathfrak m')^N P$。
\item $f_j(a'_1, \ldots, a'_u) = \sum c'_{ji} b'_i$ かつ
$\sum k'_{ij}b'_j = 0$。
\end{enumerate}
$A = B/(b'_1, \ldots, b'_r)$ とおき、$A$ における $\mathfrak m_B$ の像を
$\mathfrak m_A$ と表す。（$A$ は $\Lambda$ 上本質的有限型であり、証明の最後で
$\Lambda$ 上有限型な $A$ を得る方法を示す。）$a'_i$ は
$(b'_1, \ldots, b'_r)$ を法として所望の方程式を満たすので、
$y_i \mapsto a'_i$ で定まる環準同型 $C \to A$ が存在する。
上の性質 (2) により、$P = B^\wedge$ の商として
$A/\mathfrak m_A^N = R/\mathfrak m_R^N$ であることに注意せよ。
$x_A = x_C|_{\Spec(A)}$ とおく。写像
$$
C \to A \to A/\mathfrak m_A^N \cong R/\mathfrak m_R^N
\quad\text{and}\quad
C \to R \to R/\mathfrak m_R^N
$$
は等しいので、同型 $A/\mathfrak m_A^N = R/\mathfrak m_R^N$ を介して、
$x_A$ と $x$ は $\mathfrak m_R^N$ を法として一致する。ここまでで、補題の
主張の性質 (1) -- (5) を示した。(6) を示すため、
$$
P^{\oplus t} \xrightarrow{K}
P^{\oplus r} \xrightarrow{(b_1, \ldots, b_r)}
P
\quad\text{and}\quad
P^{\oplus t} \xrightarrow{K'}
P^{\oplus r} \xrightarrow{(b'_1, \ldots, b'_r)}
P
$$
が、$(\mathfrak m')^N$ を法として合同な二つの $P$-加群の複体であり、最初の
ものは完全であることに注意する。上での $N$ の選択により、
『代数についてさらに』の補題
\ref{more-algebra-lemma-approximate-complex-graded} から、
$R = P/(b_1, \ldots, b_r)$ と
$P/(b'_1, \ldots, b'_r) = B^\wedge/(b'_1, \ldots, b'_r) = A^\wedge$ は
同型な付随次数付き代数をもつ。これが示すべきことであった。

\medskip\noindent
証明の最後のこの段落では、$A$ が $S$ 上本質的有限型だが、まだ有限型ではない
という問題を処理する。上の構成により $A = B/(b'_1, \ldots, b'_r)$ かつ
$\mathfrak m_A \subset A$ であり、$B$ は
$\Lambda[x_1, \ldots, x_s]_{\mathfrak m'}$ 上エタールである。したがって $A$ は
Noether 環 $\Lambda[x_1, \ldots, x_s]_{\mathfrak m'}$ 上有限型である。
ゆえに、ある有限型 $\Lambda[x_1, \ldots, x_n]$-代数 $A_0$ に対して
$A = (A_0)_{\mathfrak m'}$ と書ける。このとき『代数』の補題
\ref{algebra-lemma-localization-colimit} により、
$A = \colim (A_0)_f$ である。ここで
$f \in \Lambda[x_1, \ldots, x_n] \setminus \mathfrak m'$ である。
$p : \mathcal{X} \to (\Sch/S)_{fppf}$ は対象上で極限を保つので、上のような
ある $f$ に対して、$x_A$ は $\Spec((A_0)_f)$ 上のある対象
$x_{(A_0)_f}$ から得られる。$A$ を $(A_0)_f$ に、$x_A$ を
$x_{(A_0)_f}$ に、$\mathfrak m_A$ を $(A_0)_f \cap \mathfrak m_A$ に置き換えれば、
証明は完了する。
\end{proof}





\section{極限を保つこと}
\label{artin-section-limits}

\noindent
射 $p : \mathcal{X} \to (\Sch/S)_{fppf}$ が、『表現可能性の判定条件』の節
\ref{criteria-section-limit-preserving} で定義した意味で対象上で極限を保つとは、
下の定義における関手が本質的全射であることをいう。しかし『例』の節
\ref{examples-section-limit-preserving} の例が示すように、これは極限を保つことと
同値ではない。

\begin{definition}
\label{artin-definition-limit-preserving}
$S$ をスキームとする。$\mathcal{X}$ を $(\Sch/S)_{fppf}$ 上の
グルーポイドをファイバーとする圏とする。$S$ 上のアフィンスキームの有向逆系
$T_i$ の極限 $T = \lim T_i$ である任意の $S$ 上のアフィンスキーム $T$ に対して、
ファイバー圏の同値
$$
\colim \mathcal{X}_{T_i} \longrightarrow \mathcal{X}_T
$$
が成り立つとき、$\mathcal{X}$ は {\it 極限を保つ} という。
\end{definition}

\noindent
これが何を意味するか詳しく述べる。まず、$T_i$ 上の $\mathcal{X}$ の対象
$x, y$ が与えられたとき、
$$
\Mor_{\mathcal{X}_T}(x|_T, y|_T) =
\colim_{i' \geq i} \Mor_{\mathcal{X}_{T_{i'}}}(x|_{T_{i'}}, y|_{T_{i'}})
$$
が成り立つべきである。次に、$\mathcal{X}_T$ のすべての対象は、ある $i$ に
ついて $T_i$ 上の対象の制限と同型であるべきである。第一の条件は、前層
$\mathit{Isom}_\mathcal{X}(x, y)$（『スタック』の定義
\ref{stacks-definition-mor-presheaf} を参照せよ）が極限を保つことを意味する。

\begin{lemma}
\label{artin-lemma-fibre-product-limit-preserving}
$S$ をスキームとする。$p : \mathcal{X} \to \mathcal{Y}$ と
$q : \mathcal{Z} \to \mathcal{Y}$ を、$(\Sch/S)_{fppf}$ 上の
グルーポイドをファイバーとする圏の $1$-射とする。
\begin{enumerate}
\item $\mathcal{X} \to (\Sch/S)_{fppf}$ と
$\mathcal{Z} \to (\Sch/S)_{fppf}$ が対象上で極限を保ち、$\mathcal{Y}$ が極限を
保つならば、$\mathcal{X} \times_\mathcal{Y} \mathcal{Z} \to
(\Sch/S)_{fppf}$ は対象上で極限を保つ。
\item $\mathcal{X}$, $\mathcal{Y}$, $\mathcal{Z}$ が極限を保つならば、
$\mathcal{X} \times_\mathcal{Y} \mathcal{Z}$ も極限を保つ。
\end{enumerate}
\end{lemma}

\begin{proof}
これは形式的である。(1) を証明する。$T = \lim_{i \in I} T_i$ を
$S$ 上のアフィンスキーム $T_i$ の有向極限とする。関手
$\colim \mathcal{X}_{T_i} \to \mathcal{X}_T$ が本質的全射であることを示す。
$T$ 上のファイバー積の対象は四つ組 $(T, x, z, \alpha)$ であることを思い出そう。
ここで $x$ は $T$ 上にある $\mathcal{X}$ の対象、$z$ は $T$ 上にある
$\mathcal{Z}$ の対象、$\alpha : p(x) \to q(z)$ は $T$ 上の
$\mathcal{Y}$ のファイバー圏における射である。
$\mathcal{X}$ と $\mathcal{Z}$ に関する仮定により、ある $i$ と $T_i$ 上の対象
$x_i$, $z_i$ で、$x_i|_T \cong T$ および $z_i|_T \cong z$ を満たすものを
とれる。このとき $\alpha$ は同型
$p(x_i)|_T \to q(z_i)|_T$ に対応し、$\mathcal{Y}$ に関する仮定により、これは
同型 $\alpha_{i'} : p(x_i)|_{T_{i'}} \to q(z_i)|_{T_{i'}}$ から得られる。
$i$ を $i'$ に、$x_i$ を $x_i|_{T_{i'}}$ に、$z_i$ を
$z_i|_{T_{i'}}$ に置き換えると、$(T_i, x_i, z_i, \alpha_i)$ は $T_i$ 上の
ファイバー積の対象であり、その $T$ への制限は所望どおり
$(T, x, z, \alpha)$ と同型である。

\medskip\noindent
(2) において $\colim \mathcal{X}_{T_i} \to \mathcal{X}_T$ が充満忠実であることを
示す議論は省略する。
\end{proof}

\begin{lemma}
\label{artin-lemma-limit-preserving-algebraic-space}
$S$ をスキームとし、$\mathcal{X}$ を $S$ 上の代数スタックとする。このとき
次は同値である。
\begin{enumerate}
\item $\mathcal{X}$ は集合亜群のスタックであり、
$\mathcal{X} \to (\Sch/S)_{fppf}$ は対象上で極限を保つ。
\item $\mathcal{X}$ は集合亜群のスタックであり、極限を保つ。
\item $\mathcal{X}$ は局所有限表示な代数空間により表現可能である。
\end{enumerate}
\end{lemma}

\begin{proof}
三つの仮定のいずれのもとでも、$\mathcal{X}$ は $S$ 上の代数空間 $X$ により
表現可能である。『代数スタック』の命題
\ref{algebraic-proposition-algebraic-stack-no-automorphisms} を参照せよ。
集合亜群の間の関手は、同型類上で全射であることと圏同値であることが同値なので、
(1) と (2) が同値であることは明らかである。最後に、(1) と (3) の同値性は
『空間の極限』の命題
\ref{spaces-limits-proposition-characterize-locally-finite-presentation} による。
\end{proof}

\begin{lemma}
\label{artin-lemma-diagonal}
$S$ をスキームとする。$\mathcal{X}$ を $(\Sch/S)_{fppf}$ 上の
グルーポイドをファイバーとする圏とする。
$\Delta : \mathcal{X} \to \mathcal{X} \times \mathcal{X}$ が代数空間により
表現可能で、$\mathcal{X}$ が極限を保つと仮定する。このとき $\Delta$ は
局所有限型である。
\end{lemma}

\begin{proof}
『表現可能性の判定条件』の補題
\ref{criteria-lemma-check-property-limit-preserving} を適用する。
$V$ を $S$ 上局所有限表示なアフィンスキーム $V$ とし、$\theta$ を $V$ 上の
$\mathcal{X} \times \mathcal{X}$ の対象とする。
$\mathcal{X} \times_{\Delta, \mathcal{X} \times \mathcal{X}, \theta}
(\Sch/V)_{fppf}$ を表現する代数空間を $F_\theta$ とし、標準射を
$f_\theta : F_\theta \to V$ とする（『代数スタック』の節
\ref{algebraic-section-morphisms-representable-by-algebraic-spaces} を参照せよ）。
$F_\theta \to V$ が対応する性質をもつことを示せば十分である。補題
\ref{artin-lemma-fibre-product-limit-preserving} と
\ref{artin-lemma-limit-preserving-algebraic-space} により、$F_\theta \to S$ は
局所有限表示である。したがって『空間の射』の補題
\ref{spaces-morphisms-lemma-permanence-finite-type} により、
$F_\theta \to V$ は局所有限型である。
\end{proof}






\section{半普遍性}
\label{artin-section-versality}

\noindent
前節では、完備局所環上の対象を代数的な対象により近似する方法を説明した。
しかし、スタック $\mathcal{X}$ が代数スタックであることを示すには、スキームから
$\mathcal{X}$ への滑らかな $1$-射を見つける必要がある。$\mathcal{X}$ が表現可能な
対角射をもつことをあらかじめ仮定しないので、$\mathcal{X}$ への滑らかな射について
語ることすらできない。そこで、変形理論から用語を借りて、半普遍的対象を導入する。

\begin{definition}
\label{artin-definition-versal-formal-object}
$S$ を局所 Noether スキームとする。
$p : \mathcal{X} \to (\Sch/S)_{fppf}$ をグルーポイドをファイバーとする圏とする。
$\xi = (R, \xi_n, f_n)$ を形式対象とする。$k = R/\mathfrak m$ および
$x_0 = \xi_1$ とおく。$\xi$ が前変形圏
$\mathcal{F}_{\mathcal{X}, k, x_0}$ の形式対象として（注意
\ref{artin-remark-formal-objects-match}）、『形式的変形理論』の定義
\ref{formal-defos-definition-versal} の意味で半普遍的であるとき、
$\xi$ は {\it 半普遍的} であるという。
\end{definition}

\noindent
これが何を意味するか簡潔に述べる。定義の記法を用い、閉埋め込み
$\Spec(k) \to \Spec(A) \to \Spec(B)$ 上にある $\mathcal{X}$ の射
$\xi_1 = x_0 \to y \to z$ が与えられていると仮定する。ここで $A, B$ は
剰余体 $k$ をもつ Artin 局所環である。さらに $n \geq 1$ と可換図式
$$
\vcenter{
\xymatrix{
& y \ar[ld] \\
\xi_n & \xi_1 \ar[u] \ar[l]
}
}
\quad\text{lying over}\quad
\vcenter{
\xymatrix{
& \Spec(A) \ar[ld] \\
\Spec(R/\mathfrak m^n) & \Spec(k) \ar[u] \ar[l]
}
}
$$
が与えられているとする。半普遍性とは、上のような任意のデータに対して、
$m \geq n$ と可換図式
$$
\vcenter{
\xymatrix{
& & z \ar[lldd] \\
& & y \ar[ld] \ar[u] \\
\xi_m & \xi_n \ar[l] & \xi_1 \ar[u] \ar[l]
}
}
\:\text{lying over}\:
\vcenter{
\xymatrix{
& & \Spec(B) \ar[lldd] \\
& & \Spec(A) \ar[ld] \ar[u] \\
\Spec(R/\mathfrak m^m) &
\Spec(R/\mathfrak m^n) \ar[l] &
\Spec(k) \ar[u] \ar[l]
}
}
$$
が存在することをいう。『形式的変形理論』の注意
\ref{formal-defos-remark-versal-object} と比較せよ。

\medskip\noindent
$S$ を局所 Noether スキームとする。$U$ を $S$ 上のスキームで、その構造射
$U \to S$ が局所有限型であるものとする。$u_0 \in U$ を $U$ の有限型点とする。
『射』の定義 \ref{morphisms-definition-finite-type-point} を参照せよ。
$k = \kappa(u_0)$ とおく。合成 $\Spec(k) \to S$ も有限型であることに注意せよ。
『射』の補題 \ref{morphisms-lemma-composition-finite-type} を参照せよ。
$p : \mathcal{X} \to (\Sch/S)_{fppf}$ をグルーポイドをファイバーとする圏とする。
$x$ を $U$ 上にある $\mathcal{X}$ の対象とする。$u_0$ による $x$ の引き戻しを
$x_0$ と表す。$2$-Yoneda の補題により、$x$ は $1$-射
$$
x : (\Sch/U)_{fppf} \longrightarrow \mathcal{X},
$$
に対応する。『代数スタック』の節
\ref{algebraic-section-2-yoneda} を参照せよ。$\mathcal{C}_\Lambda$ 上の
前変形圏の射
\begin{equation}
\label{artin-equation-hat-x}
\hat x :
\mathcal{F}_{(\Sch/U)_{fppf}, k, u_0}
\longrightarrow
\mathcal{F}_{\mathcal{X}, k, x_0},
\end{equation}
を得る。(\ref{artin-equation-functoriality}) を参照せよ。

\begin{definition}
\label{artin-definition-versal}
$S$ を局所 Noether スキームとする。$\mathcal{X}$ を
$(\Sch/S)_{fppf}$ 上のグルーポイドをファイバーとする圏とする。
$U$ を $S$ 上局所有限型なスキームとする。$x$ を $U$ 上にある
$\mathcal{X}$ の対象とし、$u_0$ を $U$ の有限型点とする。射 $\hat x$
（\ref{artin-equation-hat-x}）が滑らかであるとき、$x$ は $u_0$ で
{\it 半普遍的} であるという。『形式的変形理論』の定義
\ref{formal-defos-definition-smooth-morphism} を参照せよ。
\end{definition}

\noindent
この定義は $\mathcal{X}$ の形式対象に対する半普遍性の概念と一致する。

\begin{lemma}
\label{artin-lemma-versality-matches}
定義 \ref{artin-definition-versal} の記法を用いる。
$R = \mathcal{O}_{U, u_0}^\wedge$ とする。
$\xi$ を、$x|_{\Spec(R)}$ に付随する $R$ 上の $\mathcal{X}$ の形式対象とする。
(\ref{artin-equation-approximation}) を参照せよ。このとき
$$
x\text{ is versal at }u_0
\Leftrightarrow
\xi\text{ is versal}
$$
である。
\end{lemma}

\begin{proof}
$\mathcal{O}_{U, u_0}$ は剰余体 $k$ をもつ Noether 局所 $S$-代数である。
したがって $R = \mathcal{O}_{U, u_0}^\wedge$ は
$\mathcal{C}_\Lambda^\wedge$ の対象である。『形式的変形理論』の定義
\ref{formal-defos-definition-completion-CLambda} を参照せよ。
$\xi$ が半普遍的であるとは、
$\underline{\xi} : \underline{R}|_{\mathcal{C}_\Lambda} \to
\mathcal{F}_{\mathcal{X}, k, x_0}$ が滑らかであることであり、$x$ が $u_0$ で
半普遍的であるとは、$\hat x : \mathcal{F}_{(\Sch/U)_{fppf}, k, u_0}
\to \mathcal{F}_{\mathcal{X}, k, x_0}$ が滑らかであることである。
前変形圏の同一視
$$
\underline{R}|_{\mathcal{C}_\Lambda}
=
\mathcal{F}_{(\Sch/U)_{fppf}, k, u_0},
$$
がある。記法については『形式的変形理論』の注意
\ref{formal-defos-remark-formal-objects-yoneda} を参照せよ。実際、剰余体が $k$ と
同一視された Artin 局所 $S$-代数 $A$ が与えられると、
$$
\Mor_{\mathcal{C}_\Lambda^\wedge}(R, A) =
\{\varphi \in \Mor_S(\Spec(A), U) \mid \varphi|_{\Spec(k)} = u_0\}
$$
である。定義を展開すると、得られる写像
$$
\underline{R}|_{\mathcal{C}_\Lambda} =
\mathcal{F}_{(\Sch/U)_{fppf}, k, u_0}
\xrightarrow{\hat x}
\mathcal{F}_{\mathcal{X}, k, x_0},
$$
が $\underline{\xi}$ に等しいことを読者は確認できる。したがって補題が従う。
\end{proof}

\noindent
ここで健全性確認を一つ行う。

\begin{lemma}
\label{artin-lemma-versal-implies-smooth}
$S$ を局所 Noether スキームとする。$f : U \to V$ を $S$ 上局所有限型な
スキームの射とし、$u_0 \in U$ を有限型点とする。次は同値である。
\begin{enumerate}
\item $f$ は $u_0$ で滑らかである。
\item $U$ 上の $(\Sch/V)_{fppf}$ の対象とみなした $f$ は $u_0$ で
半普遍的である。
\end{enumerate}
\end{lemma}

\begin{proof}
これは『射についてさらに』の補題
\ref{more-morphisms-lemma-lifting-along-artinian-at-point} の言い換えである。
\end{proof}

\noindent
この概念は体の拡大に関してよく振る舞うことが分かる。

\begin{lemma}
\label{artin-lemma-versal-change-of-field}
$S$, $\mathcal{X}$, $U$, $x$, $u_0$ を定義 \ref{artin-definition-versal} のものとする。
$l$ を体とし、$u_{l, 0} : \Spec(l) \to U$ を、像が $u_0$ で
$l/k = \kappa(u_0)$ が有限となる射とする。
$x_{l, 0} = x_0|_{\Spec(l)}$ とおく。$\mathcal{X}$ が (RS) を満たし、$x$ が
$u_0$ で半普遍的ならば、
$$
\mathcal{F}_{(\Sch/U)_{fppf}, l, u_{l, 0}}
\longrightarrow
\mathcal{F}_{\mathcal{X}, l, x_{l, 0}}
$$
は滑らかである。
\end{lemma}

\begin{proof}
補題 \ref{artin-lemma-algebraic-stack-RS} により $(\Sch/U)_{fppf}$ は (RS) を満たす。
したがって、補題の関手は $\hat x$ に付随する関手
$$
(\mathcal{F}_{(\Sch/U)_{fppf}, k , u_0})_{l/k}
\longrightarrow
(\mathcal{F}_{\mathcal{X}, k , x_0})_{l/k}
$$
である。補題 \ref{artin-lemma-change-of-field} を参照せよ。ゆえに主張は
『形式的変形理論』の補題
\ref{formal-defos-lemma-change-of-fields-smooth} から従う。
\end{proof}

\noindent
次の補題も健全性確認の一つである。これは概ね、定義
\ref{artin-definition-versal} の意味で $x$ が $u_0$ において半普遍的ならば、$U$ から
$\mathcal{X}$ への射とみなした $x$ は、スキームによる任意の基底変換の後に
滑らかであることを意味する。

\begin{lemma}
\label{artin-lemma-base-change-versal}
$S$, $\mathcal{X}$, $U$, $x$, $u_0$ を定義 \ref{artin-definition-versal} のものとする。
次を仮定する。
\begin{enumerate}
\item $\Delta : \mathcal{X} \to \mathcal{X} \times \mathcal{X}$ は代数空間により
表現可能である。
\item $\Delta$ は局所有限型である（たとえば $\mathcal{X}$ が極限を保つ場合）。
\item $\mathcal{X}$ は (RS) を満たす。
\end{enumerate}
$V$ を $S$ 上局所有限型なスキームとし、$y$ を $V$ 上の $\mathcal{X}$ の対象とする。
$2$-ファイバー積
$$
\xymatrix{
\mathcal{Z} \ar[r] \ar[d] & (\Sch/U)_{fppf} \ar[d]^x \\
(\Sch/V)_{fppf} \ar[r]^y & \mathcal{X}
}
$$
を作る。$\mathcal{Z}$ を表現する代数空間を $Z$ とし、$z_0 \in |Z|$ を
$u_0$ 上にある有限型点とする。$x$ が $u_0$ で半普遍的ならば、射
$Z \to V$ は $z_0$ で滑らかである。
\end{lemma}

\begin{proof}
（主張中の括弧内の注意は補題 \ref{artin-lemma-diagonal} による。）仮定 (1) と
『代数スタック』の補題 \ref{algebraic-lemma-representable-diagonal} により
$Z$ が存在することに注意せよ。仮定 (2) により $Z \to V \times_S U$ は
局所有限型である。スキーム $W$、閉点 $w_0 \in W$、および $w_0$ を $z_0$ へ
写すエタール射 $W \to Z$ を選ぶ。『空間の射』の定義
\ref{spaces-morphisms-definition-finite-type-point} を参照せよ。このとき $W$ は
$S$ 上局所有限型であり、$w_0$ は $W$ の有限型点である。
$l = \kappa(z_0)$ とする。$\Spec(l) = w_0$ への引き戻しによって得られる
$\Spec(l)$ 上の $\mathcal{Z}$, $(\Sch/V)_{fppf}$,
$(\Sch/U)_{fppf}$, $\mathcal{X}$ の対象を、それぞれ
$z_{l, 0}$, $v_{l, 0}$, $u_{l, 0}$, $x_{l, 0}$ と表す。図式
$$
\xymatrix{
\mathcal{F}_{(\Sch/W)_{fppf}, l, w_0} \ar[r] &
\mathcal{F}_{\mathcal{Z}, l, z_{l, 0}} \ar[d] \ar[r] &
\mathcal{F}_{(\Sch/U)_{fppf}, l, u_{l, 0}} \ar[d] \\
& \mathcal{F}_{(\Sch/V)_{fppf}, l, v_{l, 0}} \ar[r] &
\mathcal{F}_{\mathcal{X}, l, x_{l, 0}}
}
$$
を考える。補題 \ref{artin-lemma-fibre-product-deformation-categories} により、この正方形は
前変形圏のファイバー積である。補題 \ref{artin-lemma-versal-change-of-field} により、右の
垂直矢印は滑らかである。『形式的変形理論』の補題
\ref{formal-defos-lemma-smooth-properties} により、左の垂直矢印は滑らかである。
補題 \ref{artin-lemma-formally-smooth-on-deformation-categories} により、左の水平矢印は
滑らかである。したがって『形式的変形理論』の補題
\ref{formal-defos-lemma-smooth-properties} により、写像
$$
\mathcal{F}_{(\Sch/W)_{fppf}, l, w_0} \to
\mathcal{F}_{(\Sch/V)_{fppf}, l, v_{l, 0}}
$$
は滑らかである。ゆえに『射についてさらに』の補題
\ref{more-morphisms-lemma-lifting-along-artinian-at-point} により、$W \to V$ は
$w_0$ で滑らかである。これはちょうど $Z \to V$ が $z_0$ で滑らかであることを
意味し、証明が完了する。
\end{proof}

\noindent
近似の結果を半普遍的対象の言葉で言い換える。

\begin{lemma}
\label{artin-lemma-approximate-versal}
$S$ を局所 Noether スキームとする。
$p : \mathcal{X} \to (\Sch/S)_{fppf}$ をグルーポイドをファイバーとする圏とする。
$\xi = (R, \xi_n, f_n)$ を $\mathcal{X}$ の形式対象とし、$\xi_1$ は像が
$s \in S$ である $\Spec(k) \to S$ 上にあるとする。次を仮定する。
\begin{enumerate}
\item $\xi$ は半普遍的である。
\item $\xi$ は有効である。
\item $\mathcal{O}_{S, s}$ は G-環である。
\item $p : \mathcal{X} \to (\Sch/S)_{fppf}$ は対象上で極限を保つ。
\end{enumerate}
このとき、有限型射 $U \to S$、剰余体が $k$ である有限型点
$u_0 \in U$、および $U$ 上の $\mathcal{X}$ の対象 $x$ で、$x$ は $u_0$ で
半普遍的であり、かつ
$x|_{\Spec(\mathcal{O}_{U, u_0}/\mathfrak m_{u_0}^n)} \cong \xi_n$ を
満たすものが存在する。
\end{lemma}

\begin{proof}
$\Spec(R)$ 上にある $\mathcal{X}$ の対象 $x_R$ で、付随する形式対象が
$\xi$ であるものを選ぶ。$N = 2$ とし、補題 \ref{artin-lemma-approximate} を適用する。
$A, \mathfrak m_A, x_A, \ldots$ を得る。
$\eta = (A^\wedge, \eta_n, g_n)$ を
$x_A|_{\Spec(A^\wedge)}$ に付随する形式対象とする。図式
$$
\vcenter{
\xymatrix{
& \eta \ar[d] \\
\xi \ar[r] \ar@{..>}[ru] & \xi_2 = \eta_2
}
}
\quad\text{lying over}\quad
\vcenter{
\xymatrix{
& A^\wedge \ar[d] \\
R \ar[r] \ar@{..>}[ru] & R/\mathfrak m_R^2 = A/\mathfrak m_A^2
}
}
$$
がある。$\xi$ の半普遍性は、図式中の点線の矢印を見つけられることにほかならない。
実際、『形式的変形理論』の注意 \ref{formal-defos-remark-versal-object} により、
射 $\xi \to \eta_3$, $\xi \to \eta_4$ などを逐次見つけられる。
対応する環準同型 $R \to A^\wedge$ は『形式的変形理論』の補題
\ref{formal-defos-lemma-surjective-cotangent-space} により全射である。
一方、構成により、すべての $n$ に対して
$\dim_k \mathfrak m_R^n/\mathfrak m_R^{n + 1} =
\dim_k \mathfrak m_A^n/\mathfrak m_A^{n + 1}$ である。したがって、長さの加法性と
『形式的変形理論』の補題 \ref{formal-defos-lemma-length} により、
$R/\mathfrak m_R^n$ と $A/\mathfrak m_A^n$ は $\Lambda$-加群として同じ有限の
長さをもつ。ゆえに、すべての $n$ に対して
$R/\mathfrak m_R^n \to A/\mathfrak m_A^n$ は同型であり、したがって
$R \to A^\wedge$ は同型である。よって $\eta$ は半普遍的対象と同型であり、
それ自身も半普遍的である。補題 \ref{artin-lemma-versality-matches} により、$x_A$ は
$\mathfrak m_A$ に対応する $U = \Spec(A)$ の点 $u_0$ において半普遍的である。
\end{proof}

\begin{example}
\label{artin-example-approximate-versal-implies}
この例では、補題 \ref{artin-lemma-approximate-versal} の結果が成り立つためには局所環
$\mathcal{O}_{S, s}$ が G-環でなければならないことを示す。実際、$\Lambda$ を
Noether 環、$\mathfrak m$ を $\Lambda$ の極大イデアルとする。
$R = \Lambda_\mathfrak m^\wedge$ とおく。$C$ が $\Lambda$ 上有限型であるような分解
$\Lambda \to C \to R$ をとる。
$S = \Spec(\Lambda)$, $U = S \setminus \{\mathfrak m\}$,
$S' = U \amalg \Spec(C)$ とおく。規則
$$
F(T) = 
\left\{
\begin{matrix}
* & \text{if }T \to S\text{ factors through }S' \\
\emptyset & \text{else}
\end{matrix}
\right.
$$
で定まる関手 $F : (\Sch/S)_{fppf}^{opp} \to \textit{Sets}$ を考える。
$\mathcal{X} = \mathcal{S}_F$ を $F$ に付随する集合をファイバーとする圏とする。
『代数スタック』の節 \ref{algebraic-section-split} を参照せよ。このとき
$\mathcal{X} \to (\Sch/S)_{fppf}$ は対象上で極限を保ち、$R$ 上の有効かつ
半普遍的な形式対象 $\xi$ が存在する。したがって、補題
\ref{artin-lemma-approximate-versal} の結論が $\mathcal{X}$ に対して成り立つならば、
有限型環準同型 $\Lambda \to A$ と、$\mathfrak m$ 上にある極大イデアル
$\mathfrak m_A$ で、次を満たすものが存在する。
\begin{enumerate}
\item $\kappa(\mathfrak m) = \kappa(\mathfrak m_A)$。
\item $\Lambda \to A$ と $\mathfrak m_A$ は『代数』の補題
\ref{algebra-lemma-smooth-test-artinian} の条件 (4) を満たす。
\item $\Lambda$-代数準同型 $C \to A$ が存在する。
\end{enumerate}
したがって、引用した補題により $\Lambda \to A$ は $\mathfrak m_A$ で滑らかである。
$A$ を切断することにより、$\Lambda \to A$ が $\mathfrak m_A$ でエタールであると
仮定してよい。たとえば『射についてさらに』の補題
\ref{more-morphisms-lemma-slice-smooth} を参照するか、直接論じればよい。
$C = \Lambda[y_1, \ldots, y_n]/(f_1, \ldots, f_m)$ と書く。
このとき $C \to R$ は方程式系 $f_1 = \ldots = f_m = 0$ の $R$ における解に
対応する。『環準同型の平滑化』の節
\ref{smoothing-section-approximation-G-rings} を参照せよ。したがって、上のような
すべての $\mathcal{X}$ に対して補題 \ref{artin-lemma-approximate-versal} の結論が
成り立つならば、$R$ に解をもつ方程式系は $\Lambda_{\mathfrak m}$ の Hensel 化に
解をもつ。言い換えると、$\Lambda_{\mathfrak m}^h$ に対して近似性質が成り立つ。
これは $\Lambda_{\mathfrak m}^h$ が G-環であることを意味し
（将来の参照をここに挿入する。『環準同型の平滑化』の節
\ref{smoothing-section-introduction} の議論も参照せよ）、さらに
$\Lambda_{\mathfrak m}$ が G-環であることを意味する。
\end{example}






\section{半普遍性の開性}
\label{artin-section-openness-versality}

\noindent
次に、半普遍性の開性を扱う。

\begin{definition}
\label{artin-definition-openness-versality}
$S$ を局所 Noether スキームとする。
\begin{enumerate}
\item $\mathcal{X}$ を $(\Sch/S)_{fppf}$ 上のグルーポイドをファイバーとする圏と
する。$S$ 上局所有限型なスキーム $U$、$U$ 上の $\mathcal{X}$ の対象 $x$、および
有限型点 $u_0 \in U$ が与えられ、$x$ が $u_0$ で半普遍的であるとき、開近傍
$u_0 \in U' \subset U$ で、$x$ が $U'$ のすべての有限型点において半普遍的で
あるものが存在するならば、$\mathcal{X}$ は {\it 半普遍性の開性} を満たすという。
\item $f : \mathcal{X} \to \mathcal{Y}$ を $(\Sch/S)_{fppf}$ 上の
グルーポイドをファイバーとする圏の $1$-射とする。$S$ 上局所有限型な
スキーム $U$ と、$U$ 上の $\mathcal{Y}$ の対象 $y$ が与えられたとき、
$(\Sch/U)_{fppf} \times_\mathcal{Y} \mathcal{X}$ に対して半普遍性の開性が
成り立つならば、$f$ は {\it 半普遍性の開性} を満たすという。
\end{enumerate}
\end{definition}

\noindent
半普遍性の開性は、しばしば検証が最も困難である。しかし、次の例はこの条件を
要求する必要があることを示す。

\begin{example}
\label{artin-example-versality}
$k$ を体とし、$\Lambda = k[s, t]$ とおく。規則
$$
F(A) =
\left\{
\begin{matrix}
* & \text{if there exist }f_1, \ldots, f_n \in A\text{ such that } \\
  & A = (s, t, f_1, \ldots, f_n)\text{ and } f_i s = 0\ \forall i \\
\emptyset & \text{else}
\end{matrix}
\right.
$$
で定まる関手 $F : \Lambda\text{-algebras} \longrightarrow \textit{Sets}$ を
考える。幾何学的には、$F(A) = *$ とは、$V(s, t) \subset \Spec(A)$ の
準コンパクト開近傍 $W$ で $s|_W = 0$ を満たすものが存在することを意味する。
$\mathcal{X} \subset (\Sch/\Spec(\Lambda))_{fppf}$ を、すべての $j$ に対して
$F(A_j) = *$ となるアフィン開被覆 $T = \bigcup \Spec(A_j)$ をもつスキーム
$T$ からなる充満部分圏とする。このとき $\mathcal{X}$ は [0], [1], [2], [3],
[4] を満たすが [5] を満たさない。実際、$U = \Spec(k[s, t]/(s))$ 上には、
$u_0 = (s, t)$ では半普遍的だが、他のどの点でも半普遍的でない対象 $x$ が
存在する。詳細は省略する。
\end{example}

\noindent
$S$ を局所 Noether スキームとする。$f : \mathcal{X} \to \mathcal{Y}$ を
$(\Sch/S)_{fppf}$ 上のグルーポイドをファイバーとする圏の $1$-射とする。
次の性質を考える。
\begin{equation}
\label{artin-equation-smooth}
\begin{matrix}
\text{for all fields }k\text{ of finite type over }S
\text{ and all }x_0 \in \Ob(\mathcal{X}_{\Spec(k)})\text{ the}\\
\text{map }
\mathcal{F}_{\mathcal{X}, k, x_0} \to \mathcal{F}_{\mathcal{Y}, k, f(x_0)}
\text{ of predeformation categories is smooth}
\end{matrix}
\end{equation}
この概念に関する補題をいくつか定式化する。まず、これを半普遍性（の開性）と
結びつける。

\begin{lemma}
\label{artin-lemma-versal-smooth}
$S$ を局所 Noether スキームとする。$\mathcal{X}$ を
$(\Sch/S)_{fppf}$ 上のグルーポイドをファイバーとする圏とする。
$U$ を $S$ 上局所有限型なスキームとし、$x$ を $U$ 上の $\mathcal{X}$ の対象とする。
$x$ が $U$ のすべての有限型点で半普遍的であり、$\mathcal{X}$ が (RS) を満たすと
仮定する。このとき $x : (\Sch/U)_{fppf} \to \mathcal{X}$ は
(\ref{artin-equation-smooth}) を満たす。
\end{lemma}

\begin{proof}
$l$ が $S$ 上有限型であるような射 $\Spec(l) \to U$ をとる。このとき像
$u_0 \in U$ は $U$ の有限型点で、$l/\kappa(u_0)$ は有限拡大である。
『射』の節 \ref{morphisms-section-points-finite-type} の議論を参照せよ。
したがって補題 \ref{artin-lemma-versal-change-of-field} により、
$\mathcal{F}_{(\Sch/U)_{fppf}, l, u_{l, 0}} \to
\mathcal{F}_{\mathcal{X}, l, x_{l, 0}}$ は滑らかである。
\end{proof}

\begin{lemma}
\label{artin-lemma-composition-smooth}
$S$ を局所 Noether スキームとする。$f : \mathcal{X} \to \mathcal{Y}$ と
$g : \mathcal{Y} \to \mathcal{Z}$ を、$(\Sch/S)_{fppf}$ 上の
グルーポイドをファイバーとする圏の合成可能な $1$-射とする。
$f$ と $g$ が (\ref{artin-equation-smooth}) を満たすならば、$g \circ f$ も満たす。
\end{lemma}

\begin{proof}
これは『形式的変形理論』の補題
\ref{formal-defos-lemma-smooth-properties} から形式的に従う。
\end{proof}

\begin{lemma}
\label{artin-lemma-base-change-smooth}
$S$ を局所 Noether スキームとする。$f : \mathcal{X} \to \mathcal{Y}$ と
$\mathcal{Z} \to \mathcal{Y}$ を、$(\Sch/S)_{fppf}$ 上の
グルーポイドをファイバーとする圏の $1$-射とする。
$f$ が (\ref{artin-equation-smooth}) を満たすならば、射影
$\mathcal{X} \times_\mathcal{Y} \mathcal{Z} \to \mathcal{Z}$ も満たす。
\end{lemma}

\begin{proof}
補題 \ref{artin-lemma-fibre-product-deformation-categories} と『形式的変形理論』の補題
\ref{formal-defos-lemma-smooth-properties} から直ちに従う。
\end{proof}

\begin{lemma}
\label{artin-lemma-smooth-smooth}
$S$ を局所 Noether スキームとする。$f : \mathcal{X} \to \mathcal{Y}$ を
$(\Sch/S)_{fppf}$ 上のグルーポイドをファイバーとする圏の $1$-射とする。
$f$ が対象について形式的に滑らかならば、$f$ は (\ref{artin-equation-smooth}) を満たす。
$f$ が代数空間により表現可能かつ滑らかならば、$f$ は
(\ref{artin-equation-smooth}) を満たす。
\end{lemma}

\begin{proof}
補題 \ref{artin-lemma-formally-smooth-on-deformation-categories} の言い換えである。
\end{proof}

\begin{lemma}
\label{artin-lemma-implies-smooth}
$S$ を局所 Noether スキームとする。$f : \mathcal{X} \to \mathcal{Y}$ を
$(\Sch/S)_{fppf}$ 上のグルーポイドをファイバーとする圏の $1$-射とする。
次を仮定する。
\begin{enumerate}
\item $f$ は代数空間により表現可能である。
\item $f$ は (\ref{artin-equation-smooth}) を満たす。
\item $\mathcal{X} \to (\Sch/S)_{fppf}$ は対象上で極限を保つ。
\item $\mathcal{Y}$ は極限を保つ。
\end{enumerate}
このとき $f$ は滑らかである。
\end{lemma}

\begin{proof}
証明の主要な材料は『射についてさらに』の補題
\ref{more-morphisms-lemma-lifting-along-artinian-at-point} である。これは、
$S$ 上有限型なスキームの射で (\ref{artin-equation-smooth}) を満たすものは滑らかで
あると（ほぼ）述べている。証明の他の議論は本質的に記帳である。

\medskip\noindent
$V$ を $S$ 上のスキームとし、$y$ を $V$ 上の $\mathcal{Y}$ の対象とする。
$2$-ファイバー積
$\mathcal{Z} = \mathcal{X} \times_{f, \mathcal{X}, y} (\Sch/V)_{fppf}$ を
表現する代数空間を $Z$ とする。射影 $Z \to V$ が滑らかであることを示さなければ
ならない。『代数スタック』の定義
\ref{algebraic-definition-relative-representable-property} を参照せよ。実際、
$V$ が $S$ 上局所有限表示なアフィンスキームである場合にこれを示せば十分である。
『表現可能性の判定条件』の補題
\ref{criteria-lemma-check-property-limit-preserving} を参照せよ。このとき補題
\ref{artin-lemma-limit-preserving-algebraic-space} により $(\Sch/V)_{fppf}$ は極限を保つ。
したがって補題 \ref{artin-lemma-fibre-product-limit-preserving} と
\ref{artin-lemma-limit-preserving-algebraic-space} により、$Z \to S$ は局所有限表示である。
スキーム $W$ と全射エタール射 $W \to Z$ を選ぶ。このとき $W$ は $S$ 上
局所有限表示である。

\medskip\noindent
$f$ は (\ref{artin-equation-smooth}) を満たすので、補題
\ref{artin-lemma-base-change-smooth} により $\mathcal{Z} \to (\Sch/V)_{fppf}$ も満たす。
次に、$(\Sch/W)_{fppf} \to \mathcal{Z}$ は (\ref{artin-equation-smooth}) を満たす。
これは補題 \ref{artin-lemma-smooth-smooth} による。したがって、合成
$$
(\Sch/W)_{fppf} \to \mathcal{Z} \to (\Sch/V)_{fppf}
$$
は (\ref{artin-equation-smooth}) を満たす。これは補題
\ref{artin-lemma-composition-smooth} による。『射についてさらに』の補題
\ref{more-morphisms-lemma-lifting-along-artinian-at-point} により、合成
$W \to Z \to V$ は $W$ のすべての有限型点 $w_0$ で滑らかである。滑らかな軌跡は
開なので、『射』の補題 \ref{morphisms-lemma-enough-finite-type-points} により、
$W \to V$ はスキームの滑らかな射である。したがって定義により、$Z \to V$ は
代数空間の滑らかな射である。
\end{proof}

\noindent
次の補題は、半普遍性の開性を実際に用いる方法を述べる。

\begin{lemma}
\label{artin-lemma-get-smooth}
$S$ を局所 Noether スキームとする。
$p : \mathcal{X} \to (\Sch/S)_{fppf}$ をグルーポイドをファイバーとする圏とする。
$k$ を $S$ 上の有限型な体とし、$x_0$ を $\Spec(k)$ 上の $\mathcal{X}$ の対象で、
像が $s \in S$ であるものとする。次を仮定する。
\begin{enumerate}
\item $\Delta : \mathcal{X} \to \mathcal{X} \times \mathcal{X}$ は代数空間により
表現可能である。
\item $\mathcal{X}$ は公理 [1], [2], [3] を満たす（節
\ref{artin-section-axioms} を参照せよ）。
\item $\mathcal{X}$ のすべての形式対象は有効である。
\item $\mathcal{X}$ に対して半普遍性の開性が成り立つ。
\item $\mathcal{O}_{S, s}$ は G-環である。
\end{enumerate}
このとき、有限型射 $U \to S$ と、$U$ 上の $\mathcal{X}$ の対象 $x$ で、
$$
x : (\Sch/U)_{fppf} \longrightarrow \mathcal{X}
$$
が滑らかであり、かつ剰余体が $k$ で $x|_{u_0} \cong x_0$ を満たす有限型点
$u_0 \in U$ が存在するようなものが存在する。
\end{lemma}

\begin{proof}
公理 [2]、補題 \ref{artin-lemma-deformation-category}、および注意
\ref{artin-remark-deformation-category-implies} により、
$\mathcal{F}_{\mathcal{X}, k, x_0}$ は (S1) と (S2) を満たす。
さらに公理 [3] により接空間は有限次元なので、『形式的変形理論』の補題
\ref{formal-defos-lemma-versal-object-existence} から、
$\mathcal{F}_{\mathcal{X}, k, x_0}$ は半普遍的形式対象 $\xi$ をもつ。
仮定 (3) は $\xi$ が有効であることをいう。公理 [1] と補題
\ref{artin-lemma-approximate-versal} により、有限型射 $U \to S$、$U$ 上の
$\mathcal{X}$ の対象 $x$、および剰余体が $k$ である $U$ の有限型点 $u_0$ で、
$x$ が $u_0$ で半普遍的かつ $x|_{\Spec(k)} \cong x_0$ となるものが存在する。
半普遍性の開性により $U$ を縮小し、$x$ が $U$ のすべての有限型点で半普遍的で
あると仮定してよい。ここで
$$
x : (\Sch/U)_{fppf} \longrightarrow \mathcal{X}
$$
が滑らかであると主張する。これにより補題が従う。実際、補題
\ref{artin-lemma-versal-smooth} により $x$ は (\ref{artin-equation-smooth}) を満たし、
補題 \ref{artin-lemma-implies-smooth} により証明が完了する。
\end{proof}





\section{公理}
\label{artin-section-axioms}

\noindent
$S$ を局所 Noether スキームとする。
$p : \mathcal{X} \to (\Sch/S)_{fppf}$ をグルーポイドをファイバーとする圏とする。
$\mathcal{X}$ に課す公理は次のとおりである。
\begin{enumerate}
\item[{[-1]}] 集合論的な問題に関心のない読者は無視してよい集合論的条件
\footnote{条件は次のとおりである。$k$ が $S$ 上の有限型な体を動くときの
$|\Ob(\mathcal{X}_{\Spec(k)})/\cong|$ と
$|\text{Arrows}(\mathcal{X}_{\Spec(k)})|$ のすべての基数の上限が、
$(\Sch/S)_{fppf}$ のある対象の大きさに対して $\leq$ である。}、
\item[{[0]}] $\mathcal{X}$ はエタール位相に関するグルーポイドのスタックである、
\item[{[1]}] $\mathcal{X}$ は極限を保つ、
\item[{[2]}] $\mathcal{X}$ は Rim--Schlessinger 条件 (RS) を満たす、
\item[{[3]}] すべての $k$ と $x_0$ に対して、空間
$T\mathcal{F}_{\mathcal{X}, k, x_0}$ と
$\text{Inf}(\mathcal{F}_{\mathcal{X}, k, x_0})$ は有限次元である
（\ref{artin-equation-tangent-space} および
\ref{artin-equation-infinitesimal-automorphisms} を参照せよ）、
\item[{[4]}] 関手 (\ref{artin-equation-approximation}) は同値である、
\item[{[5]}] $\mathcal{X}$ と
$\Delta : \mathcal{X} \to \mathcal{X} \times \mathcal{X}$ は半普遍性の開性を
満たす。
\end{enumerate}










\section{関手に対する公理}
\label{artin-section-axioms-functors}

\noindent
$S$ をスキームとする。$F : (\Sch/S)_{fppf}^{opp} \to \textit{Sets}$ を関手とする。
$F$ に付随する集合をファイバーとする圏を
$\mathcal{X} = \mathcal{S}_F$ と表す。『代数スタック』の節
\ref{algebraic-section-split} を参照せよ。本節では、$\mathcal{X}$ に適用される
上の内容と、$F$ に関する主張との間の翻訳を与える。

\medskip\noindent
$S$ を局所 Noether スキームとする。
$F : (\Sch/S)_{fppf}^{opp} \to \textit{Sets}$ を関手とする。
$k$ を $S$ 上有限型な体とし、$x_0 \in F(\Spec(k))$ とする。
付随する前変形圏 (\ref{artin-equation-predeformation-category}) は関手
$$
F_{k, x_0} : \mathcal{C}_\Lambda \longrightarrow \textit{Sets},
\quad
A \longmapsto \{x \in F(\Spec(A)) \mid x|_{\Spec(k)} = x_0 \}.
$$
に対応する。$\mathcal{C}_\Lambda$ 上の集合をファイバーとする余ファイバー圏と
関手 $\mathcal{C}_\Lambda \to \textit{Sets}$ を区別しないことを思い出そう。
『形式的変形理論』の注意
\ref{formal-defos-remarks-cofibered-groupoids}
（\ref{formal-defos-item-convention-cofibered-sets}）を参照せよ。
関手の変換 $a : F \to G$ が与えられたとき、$y_0 = a(x_0)$ とおけば射
$$
F_{k, x_0} \longrightarrow G_{k, y_0}
$$
を得る。(\ref{artin-equation-functoriality}) を参照せよ。補題
\ref{artin-lemma-formally-smooth-on-deformation-categories} によれば、
$a : F \to G$ が形式的に滑らかならば（『空間の射についてさらに』の定義
\ref{spaces-more-morphisms-definition-formally-smooth-etale-unramified} の意味で）、
$F_{k, x_0} \longrightarrow G_{k, y_0}$ は『形式的変形理論』の注意
\ref{formal-defos-remark-compare-smooth-schlessinger} の意味で滑らかである。

\medskip\noindent
補題 \ref{artin-lemma-pushout} によれば、$S$ 上のスキームの圏において
$Y' = Y \amalg_X X'$ であり、$X \to X'$ が無限小厚化、$X \to Y$ が
アフィンであるとき、$F$ が代数空間ならば写像
$$
F(Y \amalg_X X') \to F(Y) \times_{F(X)} F(X')
$$
は全単射である。一般の関手について、$Y, X, X'$ が $S$ 上有限型な
Artin 局所環のスペクトルである任意の押し出し $Y' = Y \amalg_X X'$ に対して
$$
F(Y \amalg_X X') \to F(Y) \times_{F(X)} F(X')
$$
が全単射であるとき、$F$ は {\it Rim--Schlessinger 条件} を満たす、または
$F$ は {\it (RS) を満たす} という。したがって、すべての代数空間は (RS) を
満たす。

\medskip\noindent
補題 \ref{artin-lemma-deformation-category} によれば、(RS) を満たす関手 $F$ が
与えられると、すべての $F_{k, x_0}$ は『形式的変形理論』の定義
\ref{formal-defos-definition-deformation-category} における変形関手、すなわち
『形式的変形理論』の注意
\ref{formal-defos-remark-compare-schlessinger-H4} における (RS) を満たす。
特に、接空間
$$
TF_{k, x_0} = \{x \in F(\Spec(k[\epsilon])) \mid x|_{\Spec(k)} = x_0\}
$$
は『形式的変形理論』の補題
\ref{formal-defos-lemma-tangent-space-vector-space} により $k$-ベクトル空間の
構造をもつ。

\medskip\noindent
補題 \ref{artin-lemma-finite-dimension} によれば、$S$ 上局所有限型な代数空間 $F$ は、
有限次元の接空間 $TF_{k, x_0}$ をもつ変形関手 $F_{k, x_0}$ を生じる。

\medskip\noindent
$F$ の {\it 形式対象}\footnote{これは Artin が形式変形と呼ぶものである。}
$\xi = (R, \xi_n)$ とは、剰余体が $S$ 上有限型である Noether 完備局所
$S$-代数 $R$ と、各添字について
$\xi_{n + 1}|_{\Spec(R/\mathfrak m^n)} = \xi_n$ を満たす元
$\xi_n \in F(\Spec(R/\mathfrak m^n))$ からなる。
形式対象 $\xi$ は $F_{R/\mathfrak m, \xi_1}$ の形式対象 $\xi$ を定める。
これが『形式的変形理論』の定義 \ref{formal-defos-definition-versal} の意味で
半普遍的であるとき、かつそのときに限り、$\xi$ は {\it 半普遍的} であるという。
形式対象 $\xi = (R, \xi_n)$ が {\it 有効} であるとは、すべての
$n \geq 1$ に対して $\xi_n = x|_{\Spec(R/\mathfrak m^n)}$ を満たす
$x \in F(\Spec(R))$ が存在することをいう。補題 \ref{artin-lemma-effective} によれば、
$F$ が代数空間ならば、すべての形式対象は有効である。

\medskip\noindent
$U$ を $S$ 上局所有限型なスキームとし、$x \in F(U)$ とする。
$u_0 \in U$ を有限型点とする。
$\xi = (\mathcal{O}_{U, u_0}^\wedge,
x|_{\Spec(\mathcal{O}_{U, u_0}/\mathfrak m_{u_0}^n)})$
が上で述べた意味で半普遍的形式対象であるとき、かつそのときに限り、$x$ は
$u_0$ で半普遍的であるという。

\medskip\noindent
$S$ を局所 Noether スキームとする。
$F : (\Sch/S)_{fppf}^{opp} \to \textit{Sets}$ を関手とする。
$F$ に課す公理は次のとおりである。
\begin{enumerate}
\item[{[-1]}] 集合論的な問題に関心のない読者は無視してよい集合論的条件
\footnote{条件は次のとおりである。$k$ が $S$ 上の有限型な体を動くときの
$|F(\Spec(k))|$ のすべての基数の上限が、$(\Sch/S)_{fppf}$ のある対象の大きさ
に対して $\leq$ である。}、
\item[{[0]}] $F$ はエタール位相に関する層である、
\item[{[1]}] $F$ は極限を保つ、
\item[{[2]}] $F$ は Rim--Schlessinger 条件 (RS) を満たす、
\item[{[3]}] すべての接空間 $TF_{k, x_0}$ は有限次元である、
\item[{[4]}] すべての形式対象は有効である、
\item[{[5]}] $F$ は半普遍性の開性を満たす。
\end{enumerate}
ここで {\it 極限を保つ} とは『空間の極限』の定義
\ref{spaces-limits-definition-locally-finite-presentation} で定義された概念であり、
{\it 半普遍性の開性} とは次をいう。$S$ 上局所有限型なスキーム $U$、
$x \in F(U)$、および $x$ が $u_0$ で半普遍的である有限型点
$u_0 \in U$ が与えられたとき、開近傍 $u_0 \in U' \subset U$ で、$x$ が
$U'$ のすべての有限型点で半普遍的であるものが存在する。






\section{代数空間}
\label{artin-section-algebraic-spaces}

\noindent
次が代数空間に関する最初の主要結果である。

\begin{proposition}
\label{artin-proposition-spaces-diagonal-representable}
$S$ を局所 Noether スキームとする。
$F : (\Sch/S)_{fppf}^{opp} \to \textit{Sets}$ を関手とする。次を仮定する。
\begin{enumerate}
\item $\Delta : F \to F \times F$ は代数空間により表現可能である。
\item $F$ は公理 [-1], [0], [1], [2], [3], [4], [5] を満たす
（節 \ref{artin-section-axioms-functors} を参照せよ）。
\item $S$ のすべての有限型点 $s$ に対して $\mathcal{O}_{S, s}$ は G-環である。
\end{enumerate}
このとき $F$ は代数空間である。
\end{proposition}

\begin{proof}
補題 \ref{artin-lemma-get-smooth} を $F$ に適用する。これを用いて、$S$ 上のすべての
有限型な体 $k$ と $x_0 \in F(\Spec(k))$ に対して、$S$ 上有限型なアフィンスキーム
$U_{k, x_0}$ と滑らかな射 $U_{k, x_0} \to F$ を、剰余体が $k$ である有限型点
$u_{k, x_0} \in U_{k, x_0}$ が存在し、$x_0$ が $u_{k, x_0}$ の像となるように
選ぶ。このとき
$$
U = \coprod\nolimits_{k, x_0} U_{k, x_0} \longrightarrow F
$$
は滑らかである\footnote{集合論的注意：公理 [-1] により添字集合に上界があるので、
この余積は $(\Sch/S)_{fppf}$ の対象と同型である。『集合』の補題
\ref{sets-lemma-what-is-in-it} を参照せよ。}。証明を完了するには、この写像が
全射であることを示せば十分である。『ブートストラップ』の補題
\ref{bootstrap-lemma-spaces-etale-smooth-cover} を参照せよ（ここで公理 [0] を用いる）。
『表現可能性の判定条件』の補題
\ref{criteria-lemma-check-property-limit-preserving} により、$V$ が $S$ 上
局所有限表示なアフィンスキームであるような $V \to F$ に対して、
$U \times_F V \to V$ が全射であることを示せば十分である。
$U \times_F V \to V$ は滑らかなので、その像は開である。したがって、その像が
$U \times_F V \to V$ の像として $V$ のすべての有限型点を含むことを示せば
十分である。『射』の補題
\ref{morphisms-lemma-enough-finite-type-points} を参照せよ。
$v_0 \in V$ を有限型点とする。このとき $k = \kappa(v_0)$ は $S$ 上の有限型な
体である。合成 $\Spec(k) \xrightarrow{v_0} V \to F$ を $x_0$ と表す。
このとき $(u_{k, x_0}, v_0) : \Spec(k) \to U \times_F V$ は $v_0$ へ写る点であり、
主張が従う。
\end{proof}

\begin{lemma}
\label{artin-lemma-monomorphism}
$S$ を局所 Noether スキームとする。$a : F \to G$ を関手
$(\Sch/S)_{fppf}^{opp} \to \textit{Sets}$ の変換とする。次を仮定する。
\begin{enumerate}
\item $a$ は単射である。
\item $F$ は公理 [0], [1], [2], [4], [5] を満たす。
\item $S$ のすべての有限型点 $s$ に対して $\mathcal{O}_{S, s}$ は G-環である。
\item $G$ は $S$ 上局所有限型な代数空間である。
\end{enumerate}
このとき $F$ は代数空間である。
\end{lemma}

\begin{proof}
補題 \ref{artin-lemma-finite-dimension} により、関手 $G$ は [3] を満たす。
$F \to G$ は単射なので、$F$ も [3] を満たす。また、$F \to G$ が単射なので、
スキーム $U$, $V$ と射 $U \to F$, $V \to F$ が与えられたとき、
$U \times_F V = U \times_G V$ である。したがって、仮定により $G$ について
成り立つことから、$\Delta : F \to F \times F$ は（スキームにより）表現可能で
ある。ゆえに命題 \ref{artin-proposition-spaces-diagonal-representable} を適用できる
\footnote{集合論的条件 [-1] は $G$ について成り立つので、$F$ についても
成り立つ。詳細は省略する。}。
\end{proof}











\section{代数スタック}
\label{artin-section-algebraic-stacks}

\noindent
命題 \ref{artin-proposition-second-diagonal-representable} が代数スタックに関する最初の
主要結果である。

\begin{lemma}
\label{artin-lemma-diagonal-representable}
$S$ を局所 Noether スキームとする。
$p : \mathcal{X} \to (\Sch/S)_{fppf}$ をグルーポイドをファイバーとする圏とする。
次を仮定する。
\begin{enumerate}
\item $\Delta : \mathcal{X} \to \mathcal{X} \times \mathcal{X}$ は代数空間により
表現可能である。
\item $\mathcal{X}$ は公理 [-1], [0], [1], [2], [3] を満たす
（節 \ref{artin-section-axioms} を参照せよ）。
\item $\mathcal{X}$ のすべての形式対象は有効である。
\item $\mathcal{X}$ は半普遍性の開性を満たす。
\item $S$ のすべての有限型点 $s$ に対して $\mathcal{O}_{S, s}$ は G-環である。
\end{enumerate}
このとき $\mathcal{X}$ は代数スタックである。
\end{lemma}

\begin{proof}
補題 \ref{artin-lemma-get-smooth} を $\mathcal{X}$ に適用する。これを用いて、$S$ 上の
すべての有限型な体 $k$ と、対象の各同型類
$x_0 \in \Ob(\mathcal{X}_{\Spec(k)})$ に対して、$S$ 上有限型なアフィンスキーム
$U_{k, x_0}$ と滑らかな射 $(\Sch/U_{k, x_0})_{fppf} \to \mathcal{X}$ を、
剰余体が $k$ である有限型点 $u_{k, x_0} \in U_{k, x_0}$ が存在し、$x_0$ が
$u_{k, x_0}$ の像となるように選ぶ。このとき
$$
(\Sch/U)_{fppf} \to \mathcal{X},
\quad\text{with}\quad
U = \coprod\nolimits_{k, x_0} U_{k, x_0}
$$
は滑らかである\footnote{集合論的注意：公理 [-1] により添字集合に上界があるので、
この余積は $(\Sch/S)_{fppf}$ の対象と同型である。『集合』の補題
\ref{sets-lemma-what-is-in-it} を参照せよ。}。証明を完了するには、この写像が
全射であることを示せば十分である。『表現可能性の判定条件』の補題
\ref{criteria-lemma-stacks-etale} を参照せよ（ここで公理 [0] を用いる）。
『表現可能性の判定条件』の補題
\ref{criteria-lemma-check-property-limit-preserving} により、$V$ が $S$ 上
局所有限表示なアフィンスキームであるような
$y : (\Sch/V)_{fppf} \to \mathcal{X}$ に対して、
$(\Sch/U)_{fppf} \times_\mathcal{X} (\Sch/V)_{fppf} \to (\Sch/V)_{fppf}$ が
全射であることを示せば十分である。仮定 (1) により、ファイバー積
$(\Sch/U)_{fppf} \times_\mathcal{X} (\Sch/V)_{fppf}$ は代数空間 $W$ により
表現される。このとき $W \to V$ は滑らかなので、その像は開である。したがって、
$W \to V$ の像が $V$ のすべての有限型点を含むことを示せば十分である。『射』の補題
\ref{morphisms-lemma-enough-finite-type-points} を参照せよ。
$v_0 \in V$ を有限型点とする。このとき $k = \kappa(v_0)$ は $S$ 上の有限型な
体である。$v_0$ による $y$ の引き戻しを $x_0 = y|_{\Spec(k)}$ と表す。
すると $(u_{k, x_0}, v_0)$ は射 $\Spec(k) \to W$ を与え、その $W \to V$ との
合成は $v_0$ となるので、主張が従う。
\end{proof}

\begin{proposition}
\label{artin-proposition-second-diagonal-representable}
$S$ を局所 Noether スキームとする。
$p : \mathcal{X} \to (\Sch/S)_{fppf}$ をグルーポイドをファイバーとする圏とする。
次を仮定する。
\begin{enumerate}
\item $\Delta_\Delta : \mathcal{X} \to
\mathcal{X} \times_{\mathcal{X} \times \mathcal{X}} \mathcal{X}$ は代数空間により
表現可能である。
\item $\mathcal{X}$ は公理 [-1], [0], [1], [2], [3], [4], [5] を満たす
（節 \ref{artin-section-axioms} を参照せよ）。
\item $S$ のすべての有限型点 $s$ に対して $\mathcal{O}_{S, s}$ は G-環である。
\end{enumerate}
このとき $\mathcal{X}$ は代数スタックである。
\end{proposition}

\begin{proof}
まず $\Delta : \mathcal{X} \to \mathcal{X} \times \mathcal{X}$ が代数空間により
表現可能であることを示す。そのためには、$V$ が $S$ 上局所有限表示な任意の
アフィンスキームで、$y$ が $V$ 上の $\mathcal{X} \times \mathcal{X}$ の対象で
あるとき、
$$
\mathcal{Y} =
\mathcal{X} \times_{\Delta, \mathcal{X} \times \mathcal{X}, y} (\Sch/V)_{fppf}
$$
が代数空間により表現可能であることを示せば十分である。『表現可能性の判定条件』の
補題 \ref{criteria-lemma-check-representable-limit-preserving}\footnote{『表現可能性の
判定条件』の補題 \ref{criteria-lemma-check-representable-limit-preserving} における
集合論的条件は成り立つ。$\mathcal{Y}$ を表現する代数空間 $Y$ の大きさには適切な
上界がある。実際、$Y \to S$ は局所有限型となり、$Y$ は公理 [-1] を満たす。
詳細は省略する。} を参照せよ。
$\mathcal{Y}$ は集合亜群をファイバーとすることに注意せよ（『スタック』の補題
\ref{stacks-lemma-isom-as-2-fibre-product}）。同型類の関手を
$Y : (\Sch/S)_{fppf}^{opp} \to \textit{Sets}$,
$T \mapsto \Ob(\mathcal{Y}_T)/\cong$ とする。命題
\ref{artin-proposition-spaces-diagonal-representable} を適用して、$Y$ が代数空間であることを
示す。

\medskip\noindent
条件 (1) と『表現可能性の判定条件』の補題
\ref{criteria-lemma-second-diagonal} により、
$\Delta_\mathcal{Y} : \mathcal{Y} \to \mathcal{Y} \times \mathcal{Y}$
（したがって $Y \to Y \times Y$ も）は代数空間により表現可能である。
『スタック』の補題 \ref{stacks-lemma-stack-in-setoids-characterize} および
\ref{stacks-lemma-2-fibre-product-gives-stack-in-setoids} により、$Y$ はエタール位相に
関する層、すなわち公理 [0] を満たすことに注意せよ。また、補題
\ref{artin-lemma-fibre-product-limit-preserving} により $Y$ は極限を保つ、すなわち
[1] を満たす。補題 \ref{artin-lemma-algebraic-stack-RS} と
\ref{artin-lemma-fibre-product-RS} により、$Y$ は (RS) を満たす、すなわち公理 [2] を
満たす。$Y$ に対する公理 [3] は補題 \ref{artin-lemma-finite-dimension} と
\ref{artin-lemma-fibre-product-tangent-spaces} から従う。公理 [4] は補題
\ref{artin-lemma-effective} と \ref{artin-lemma-fibre-product-effective} から従う。
$Y$ に対する公理 [5] は、$\mathcal{X}$ に対する公理 [5] の一部である
$\Delta_\mathcal{X}$ の半普遍性の開性から直接従う。したがって命題
\ref{artin-proposition-spaces-diagonal-representable} のすべての仮定が満たされ、$Y$ は
代数空間である。

\medskip\noindent
ここまでくると、補題 \ref{artin-lemma-diagonal-representable} により
$\mathcal{X}$ は代数スタックである。
\end{proof}





\section{強い Rim--Schlessinger 条件}
\label{artin-section-RS-star}

\noindent
本章の残りでは、Rim--Schlessinger 条件の次の真に強い版が重要な役割を果たす。

\begin{definition}
\label{artin-definition-RS-star}
$S$ をスキームとする。$\mathcal{X}$ を $(\Sch/S)_{fppf}$ 上の
グルーポイドをファイバーとする圏とする。$S$-代数のファイバー積図式
$$
\xymatrix{
B' \ar[r] & B \\
A' = A \times_B B' \ar[u] \ar[r] & A \ar[u]
}
$$
が与えられ、$B' \to B$ が平方零な核をもつ全射であるとき、ファイバー圏の関手
$$
\mathcal{X}_{\Spec(A')}
\longrightarrow
\mathcal{X}_{\Spec(A)} \times_{\mathcal{X}_{\Spec(B)}} \mathcal{X}_{\Spec(B')}
$$
が圏同値であるならば、$\mathcal{X}$ は {\it 条件 (RS*)} を満たすという。
\end{definition}

\noindent
いくつか注意する。定義 \ref{artin-definition-RS-star} のような
$A \to B \leftarrow B'$ に対して、
\begin{enumerate}
\item スキームの圏において
$\Spec(A') = \Spec(A) \amalg_{\Spec(B)} \Spec(B')$ である
（『射についてさらに』の補題
\ref{more-morphisms-lemma-pushout-along-thickening} を参照せよ）、また
\item $\mathcal{X}$ が代数スタックならば、補題
\ref{artin-lemma-algebraic-stack-RS-star} により $\mathcal{X}$ は (RS*) を満たす。
\end{enumerate}
$S$ が局所 Noether ならば、さらに
\begin{enumerate}
\item[(3)] $A$, $B$, $B'$ が $S$ 上有限型で $B$ が $A$ 上有限ならば、$A'$ は
$S$ 上有限型である\footnote{$\Spec(A)$ が $S$ のアフィン開集合へ写る場合、これは
『代数についてさらに』の補題
\ref{more-algebra-lemma-fibre-product-finite-type} から従う。一般の場合は
『代数についてさらに』の補題
\ref{more-algebra-lemma-diagram-localize} を用いて従う。}、また
\item[(4)] $\mathcal{X}$ が (RS*) を満たすならば、$\mathcal{X}$ は (RS) を満たす。
実際、(RS) は $A$, $B$, $B'$ が Artin 局所である (RS*) の場合をちょうど
覆っている。
\end{enumerate}

\begin{lemma}
\label{artin-lemma-algebraic-stack-RS-star}
$\mathcal{X}$ を基礎 $S$ 上の代数スタックとする。このとき $\mathcal{X}$ は
(RS*) を満たす。
\end{lemma}

\begin{proof}
これは補題 \ref{artin-lemma-pushout} から従う。定義
\ref{artin-definition-RS-star} に続く注意を参照せよ。
\end{proof}

\begin{lemma}
\label{artin-lemma-fibre-product-RS-star}
$S$ をスキームとする。$p : \mathcal{X} \to \mathcal{Y}$ と
$q : \mathcal{Z} \to \mathcal{Y}$ を、$(\Sch/S)_{fppf}$ 上の
グルーポイドをファイバーとする圏の $1$-射とする。
$\mathcal{X}$, $\mathcal{Y}$, $\mathcal{Z}$ が (RS*) を満たすならば、
$\mathcal{X} \times_\mathcal{Y} \mathcal{Z}$ も満たす。
\end{lemma}

\begin{proof}
証明は補題 \ref{artin-lemma-fibre-product-RS} の証明とまったく同じである。
\end{proof}








\section{半普遍性と一般化}
\label{artin-section-generalize-versality}

\noindent
(RS*) を満たし、極限を保つスタックについて、半普遍性が一般化のもとで保たれる
ことを示す。初読では本節を飛ばすことを勧める。

\begin{lemma}
\label{artin-lemma-single-point}
$S$ を局所 Noether スキームとする。$\mathcal{X}$ を
$(\Sch/S)_{fppf}$ 上のグルーポイドをファイバーとする圏で、(RS*) をもつものとする。
$x$ を $S$ 上有限型なアフィンスキーム $U$ 上の $\mathcal{X}$ の対象とする。
$u \in U$ を有限型点とし、$x$ は $u$ で半普遍的でないとする。このとき、
$U \to T$ 上にある
$\mathcal{X}$ の射 $x \to y$ で、次を満たすものが存在する。
\begin{enumerate}
\item 射 $U \to T$ は一階無限小厚化である。
\item 短完全列
$$
0 \to \kappa(u) \to \mathcal{O}_T \to \mathcal{O}_U \to 0
$$
がある。
\item $u$ の開近傍 $W \subset T$ と射 $\beta : y|_W \to x$ からなる組
$(W, \alpha)$ で、合成
$$
x|_{U \cap W} \xrightarrow{\text{restriction of }x \to y}
y|_W \xrightarrow{\beta} x
$$
が標準射 $x|_{U \cap W} \to x$ となるものは存在 {\bf しない}。
\end{enumerate}
\end{lemma}

\begin{proof}
$R = \mathcal{O}_{U, u}^\wedge$ とおき、$k = \kappa(u)$ を $R$ の剰余体とする。
$\xi$ を $x$ に付随する $R$ 上の $\mathcal{X}$ の形式対象とする。$x$ は $u$ で
半普遍的でないので、補題 \ref{artin-lemma-versality-matches} により $\xi$ は半普遍的で
ない。定義 \ref{artin-definition-versal-formal-object} に続く議論によれば、これは、
閉埋め込み $\Spec(k) \to \Spec(A) \to \Spec(B)$ 上にある $\mathcal{X}$ の射
$\xi_1 \to x_A \to x_B$ を見つけられることを意味する。ここで $A, B$ は剰余体
$k$ をもつ Artin 局所環である。さらに $n \geq 1$ と可換図式
$$
\vcenter{
\xymatrix{
& x_A \ar[ld] \\
\xi_n & \xi_1 \ar[u] \ar[l]
}
}
\quad\text{lying over}\quad
\vcenter{
\xymatrix{
& \Spec(A) \ar[ld] \\
\Spec(R/\mathfrak m^n) & \Spec(k) \ar[u] \ar[l]
}
}
$$
で、$m \geq n$ と可換図式
$$
\vcenter{
\xymatrix{
& & x_B \ar[lldd] \\
& & x_A \ar[ld] \ar[u] \\
\xi_m & \xi_n \ar[l] & \xi_1 \ar[u] \ar[l]
}
}
\:\text{lying over}\:
\vcenter{
\xymatrix{
& & \Spec(B) \ar[lldd] \\
& & \Spec(A) \ar[ld] \ar[u] \\
\Spec(R/\mathfrak m^m) &
\Spec(R/\mathfrak m^n) \ar[l] &
\Spec(k) \ar[u] \ar[l]
}
}
$$
が存在 {\bf しない} ようなものがある。さらに、$B \to A$ が小拡大、すなわち
全射 $B \to A$ の核 $I$ が $A$-加群として $k$ と同型であると仮定してよい。
これは『形式的変形理論』の注意
\ref{formal-defos-remark-versal-object} から従う。ここで単に
$$
T = U \amalg_{\Spec(A)} \Spec(B)
$$
と定義する。性質 (RS*) により、$T$ 上の $y$ で、その $\Spec(B)$ への制限が
$x_B$、$U$ への制限が $x$ となるものを得る（これが $U \to T$ 上の矢印
$x \to y$ を与える）。証明を完了するため、条件 (1), (2), (3) を検証する。

\medskip\noindent
押し出しの構成により、行が完全な可換図式
$$
\xymatrix{
0 \ar[r] &
I \ar[r] &
B \ar[r] &
A \ar[r] &
0 \\
0 \ar[r] &
I \ar[r] \ar[u] &
\Gamma(T, \mathcal{O}_T) \ar[r] \ar[u] &
\Gamma(U, \mathcal{O}_U) \ar[r] \ar[u] &
0
}
$$
を得る。これは直ちに (1) と (2) を証明する。証明を完了するため、背理法で
論じる。(3) のような組 $(W, \beta)$ があると仮定する。
$\Spec(B) \to T$ は $W$ を経由するので、射
$$
x_B \to y|_W \xrightarrow{\beta} x
$$
を得る。$B$ は剰余体 $k = \kappa(u)$ をもつ Artin 局所環なので、$x_B \to x$ は
ある $m \geq n$ に対して
$\Spec(\mathcal{O}_{U, u}/\mathfrak m_u^m)$ を経由する射
$\Spec(B) \to U$ 上にある。言い換えると、$x_B \to x$ は $\xi_m$ を経由し、
写像 $x_B \to \xi_m$ を与える。条件 (3) の射 $\alpha$ に対する両立条件は、
図式
$$
\xymatrix{
x_B \ar[d] & x_A \ar[d] \ar[l] \\
\xi_m & \xi_n \ar[l]
}
$$
が可換であるという条件になる。これは求めていた矛盾を与える。
\end{proof}

\begin{lemma}
\label{artin-lemma-generalization-versality}
$S$ を局所 Noether スキームとする。$\mathcal{X}$ を
$(\Sch/S)_{fppf}$ 上のグルーポイドをファイバーとする圏とする。次を仮定する。
\begin{enumerate}
\item $\Delta : \mathcal{X} \to \mathcal{X} \times \mathcal{X}$ は代数空間により
表現可能である。
\item $\mathcal{X}$ は (RS*) をもつ。
\item $\mathcal{X}$ は極限を保つ。
\end{enumerate}
$x$ を $S$ 上有限型なスキーム $U$ 上の $\mathcal{X}$ の対象とする。
$u \leadsto u_0$ を $U$ の有限型点の特殊化とし、$x$ は $u_0$ で半普遍的で
あるとする。このとき $x$ は $u$ で半普遍的である。
\end{lemma}

\begin{proof}
$U$ を縮小して、$U$ がアフィンで、$U$ が $S$ のアフィン開集合
$\Spec(\Lambda)$ へ写ると仮定してよい。$x$ が $u$ で半普遍的でないならば、
補題 \ref{artin-lemma-single-point} のような $U \to T$ 上の $x \to y$ を選べる。
$U = \Spec(R_0)$ および $T = \Spec(R)$ と書く。射 $U \to T$ は、核が平方零な
イデアルである全射環準同型 $R \to R_0$ に対応する。仮定 (3) により、ある有限型
$\Lambda$-部分代数 $R' \subset R$ に対する $U' = \Spec(R')$ 上の対象 $x'$ から
$y$ が得られる。$R'$ を増やして、$R' \to R_0$ が全射である、したがって
$U \subset U'$ が一階無限小厚化であると仮定してよく、実際そのようにする。
こうして
$$
x \to y \to x'
\text{ lying over }
U \to T \to U'
$$
を得る。仮定 (1) により、$S$ 上の代数空間 $Z$ が
$$
(\Sch/U)_{fppf} \times_{x, \mathcal{X}, x'} (\Sch/U')_{fppf}
$$
を表現する。『代数スタック』の補題
\ref{algebraic-lemma-representable-diagonal} を参照せよ。
$2$-ファイバー積の構成により、$Z$ の $V$-値点は、射
$a : V \to U$, $a' : V \to U'$ と射 $\alpha : a^*x \to (a')^*x'$ からなる三つ組
$(a, a', \alpha)$ に対応する。可換図式
$$
\xymatrix{
U \ar[rd] \ar[rdd] \ar[rrd] \\
& Z \ar[r]_{p'} \ar[d]^p & U' \ar[d] \\
& U \ar[r] & S
}
$$
を得る。射 $i : U \to Z$ は同型 $x \to x'|_U$ から得られる。
$z_0 = i(u_0) \in Z$ とおく。補題 \ref{artin-lemma-base-change-versal} により
$Z \to U'$ は $z_0$ で滑らかである。$U$ を $u_0$ のアフィン開近傍に、$U'$ を
対応する開集合に、$Z$ を $p$ と $p'$ によるこれらの開集合の逆像の共通部分に
置き換えると、$Z \to U'$ が $i(U)$ に沿って滑らかな状況になる。
$u \leadsto u_0$ なので、点 $u$ はこの開集合に属する。補題
\ref{artin-lemma-single-point} の条件 (3) は $U$ の縮小で明らかに保たれる
（スキーム $U$, $T$, $U'$ はすべて同じ台位相空間をもつ）。
$U \to U'$ はアフィンスキームの一階無限小厚化なので、
$p' \circ i' = \text{id}_{U'}$ を満たし、$U$ への制限が $i$ となる射
$i' : U' \to Z$ を選べる（『空間の射についてさらに』の補題
\ref{spaces-more-morphisms-lemma-smooth-formally-smooth}）。
普遍射 $p^*x \to (p')^*x'$ を $i'$ で引き戻して、$p \circ i' : U' \to U$ 上の射
$$
x' \to x
$$
で、その合成
$$
x \to x' \to x
$$
が恒等射であるものを得る。射 $T \to U'$ 上の $y \to x'$ があることを
思い出そう。合成して射 $y \to x$ を得るが、その存在は補題
\ref{artin-lemma-single-point} の条件 (3) に矛盾する。この矛盾により証明が完了する。
\end{proof}






\section{強い形式的有効性}
\label{artin-section-strong-formal-effectiveness}

\noindent
本節では、形式対象の有効性の強い版が半普遍性の開性を導く方法を示す。
\cite[Theorem 1.1]{Bhatt-Algebraize} の証明は、準コンパクトかつ準分離な代数空間が
注意 \ref{artin-remark-strong-effectiveness} で論じる強い形式的有効性を満たすことを示す。
さらに、ここで展開する理論には実例がある。後にこれを用いて、連接層のスタックと
複体のモジュライに対する半普遍性の開性を示す。『Quot』の定理
\ref{quot-theorem-coherent-algebraic-general} および
\ref{quot-theorem-complexes-algebraic} を参照せよ。

\begin{lemma}
\label{artin-lemma-infinite-sequence-pre}
$S$ を局所 Noether スキームとする。$\mathcal{X}$ を
$(\Sch/S)_{fppf}$ 上のグルーポイドをファイバーとする圏で、(RS*) をもつものとする。
$x$ を $S$ 上有限型なアフィンスキーム $U$ 上の $\mathcal{X}$ の対象とする。
$u_n \in U$, $n \geq 1$ を有限型点で、(a) $n \not = m$ のとき特殊化
$u_n \leadsto u_m$ がなく、(b) すべての $n$ に対して $x$ が $u_n$ で半普遍的で
ないものとする。このとき $S$ 上の射
$$
x \to x_1 \to x_2 \to \ldots
\quad\text{in }\mathcal{X}\text{ lying over }\quad
U \to U_1 \to U_2 \to \ldots
$$
で、次を満たすものが存在する。
\begin{enumerate}
\item 各 $n$ に対して射 $U \to U_n$ は一階無限小厚化である。
\item 各 $n$ に対して短完全列
$$
0 \to \kappa(u_n) \to \mathcal{O}_{U_n} \to \mathcal{O}_{U_{n - 1}} \to 0
$$
がある。$n = 1$ に対して $U_0 = U$ とする。
\item 各 $n$ に対して、$u_n$ の開近傍 $W \subset U_n$ と射
$\alpha : x_n|_W \to x$ からなる組 $(W, \alpha)$ で、合成
$$
x|_{U \cap W} \xrightarrow{\text{restriction of }x \to x_n}
x_n|_W \xrightarrow{\alpha} x
$$
が標準射 $x|_{U \cap W} \to x$ となるものは存在 {\bf しない}。
\end{enumerate}
\end{lemma}

\begin{proof}
点 $u_n$ の間に特殊化がなく（特に $u_n$ は相異なる）、すべての $n$ に対して、
$u_n \in U'$ かつ $i = 1, \ldots, n - 1$ に対して $u_i \not \in U'$ を満たす
開集合 $U' \subset U$ をとれる。補題 \ref{artin-lemma-single-point} により、各
$n \geq 1$ に対して、
$$
x \to y_n
\quad\text{in }\mathcal{X}\text{ lying over }\quad
U \to T_n
$$
で、次を満たすものをとれる。
\begin{enumerate}
\item 射 $U \to T_n$ は一階無限小厚化である。
\item 短完全列
$$
0 \to \kappa(u_n) \to \mathcal{O}_{T_n} \to \mathcal{O}_U \to 0
$$
がある。
\item $u_n$ の開近傍 $W \subset T_n$ と射 $\beta : y_n|_W \to x$ からなる組
$(W, \alpha)$ で、合成
$$
x|_{U \cap W} \xrightarrow{\text{restriction of }x \to y_n}
y_n|_W \xrightarrow{\beta} x
$$
が標準射 $x|_{U \cap W} \to x$ となるものは存在 {\bf しない}。
\end{enumerate}
したがって、帰納的に
$$
U_1 = T_1, \quad
U_{n + 1} = U_n \amalg_U T_{n + 1}
$$
と定義できる。$x_1 = y_1$ とおき、(RS*) を用いると、$U_n$ 上では $x_n$、
$T_{n + 1}$ 上では $y_{n + 1}$ に制限される $U_{n + 1}$ 上の $x_{n + 1}$ を
帰納的に得る。$U \to U_n$ に対する性質 (1) は、『射についてさらに』の補題
\ref{more-morphisms-lemma-pushout-along-thickening} における押し出しの構成から従う。
$U_n$ に対する性質 (2) も同様に、押し出しの構成により $T_n$ に対する性質 (2)
から従う。上で論じたように $u_n$ の開近傍 $U'$ へ縮小すると、$T_n$ と $U_n$ の
対応する開部分スキームが同型であることだけから、$(U_n, x_n)$ に対する性質 (3) は
$(T_n, y_n)$ に対する性質 (3) から従う。いくつかの詳細は省略する。
\end{proof}

\begin{remark}[強い有効性]
\label{artin-remark-strong-effectiveness}
$S$ を局所 Noether スキームとする。$\mathcal{X}$ を
$(\Sch/S)_{fppf}$ 上のグルーポイドをファイバーとする圏とする。次が与えられていると
仮定する。
\begin{enumerate}
\item アフィン開集合 $\Spec(\Lambda) \subset S$、
\item 遷移写像が全射で、その核が局所冪零である $\Lambda$-代数の逆系 $(R_n)$、
\item 系 $(\Spec(R_n))$ 上にある $\mathcal{X}$ の対象の系 $(\xi_n)$。
\end{enumerate}
この状況で $R = \lim R_n$ とおく。$\Spec(R)$ 上の $\mathcal{X}$ の対象 $\xi$ で、
$\Spec(R_n)$ への制限が系 $(\xi_n)$ を与えるものが存在するとき、$(\xi_n)$ は
{\it 有効} であるという。
\end{remark}

\noindent
$S$ 上のすべての代数スタック $\mathcal{X}$ が、注意
\ref{artin-remark-strong-effectiveness} のすべての系 $(\xi_n)$ は有効である、という形の
強い有効性公理を満たすわけではない。例は『例』の節
\ref{examples-section-non-formal-effectiveness} にある。

\begin{lemma}
\label{artin-lemma-SGE-implies-openness-versality}
$S$ を局所 Noether スキームとする。$\mathcal{X}$ を
$(\Sch/S)_{fppf}$ 上のグルーポイドをファイバーとする圏とする。次を仮定する。
\begin{enumerate}
\item $\Delta : \mathcal{X} \to \mathcal{X} \times \mathcal{X}$ は代数空間により
表現可能である。
\item $\mathcal{X}$ は (RS*) をもつ。
\item $\mathcal{X}$ は極限を保つ。
\item 注意 \ref{artin-remark-strong-effectiveness} のような系 $(\xi_n)$ で、すべての
$m \geq n$ に対して $\Ker(R_m \to R_n)$ が平方零イデアルであるものは有効である。
\end{enumerate}
このとき $\mathcal{X}$ は半普遍性の開性を満たす。
\end{lemma}

\begin{proof}
$S$ 上局所有限型なスキーム $U$、$U$ の有限型点 $u_0$、および $U$ 上の
$\mathcal{X}$ の対象 $x$ で、$x$ が $u_0$ で半普遍的であるものを選ぶ。
$U$ を縮小して、$U$ がアフィンで、$U$ が $S$ のアフィン開集合
$\Spec(\Lambda)$ へ写ると仮定してよい。$E \subset U$ を、有限型点 $u$ で
$x$ が $u$ で半普遍的でないものの集合とする。補題
\ref{artin-lemma-generalization-versality} により、
$u \in E$ ならば $u_0$ は $u$ の特殊化ではない。半普遍性の開性が成り立たない
ならば、$u_0$ は $E$ の閉包 $\overline{E}$ に属する。『性質』の補題
\ref{properties-lemma-countable-dense-subset} により、$E$ と同じ閉包をもつ可算部分集合
$E' \subset E$ を選べる。『性質』の補題
\ref{properties-lemma-maximal-points} により、$E'$ の点の間には特殊化がないと仮定して
よい。上で指摘したように $u_0$ は $E'$ のどの点の特殊化でもないので、$E'$ は
（可算）無限でなければならない。したがって
$E' = \{u_1, u_2, u_3, \ldots\}$ と書け、$u_i$ の間には特殊化がなく、$u_0$ は
$E'$ の閉包に属する。

\medskip\noindent
補題 \ref{artin-lemma-infinite-sequence-pre} のように、
$U \to U_1 \to U_2 \to \ldots$ 上にある
$x \to x_1 \to x_2 \to \ldots$ を選ぶ。
$U_n = \Spec(R_n)$ および $U = \Spec(R_0)$ と書き、$R = \lim R_n$ とおく。
$R \to R_0$ は全射で、その核は平方零イデアルであることに注意せよ。
仮定 (4) により、$\Spec(R)$ 上の $\xi$ で、その $R_n$ への基底変換が $x_n$ となる
ものを得る。仮定 (3) により、ある有限型 $\Lambda$-部分代数
$R' \subset R$ に対する $U' = \Spec(R')$ 上の対象 $\xi'$ から $\xi$ が得られる。
$R'$ を増やして、$R' \to R_0$ が全射である、したがって
$U \subset U'$ が一階無限小厚化であると仮定してよく、実際そのようにする。
こうして
$$
x \to x_1 \to x_2 \to \ldots \to \xi'
\text{ lying over }
U \to U_1 \to U_2 \to \ldots \to U'
$$
を得る。仮定 (1) により、$S$ 上の代数空間 $Z$ が
$$
(\Sch/U)_{fppf} \times_{x, \mathcal{X}, \xi'} (\Sch/U')_{fppf}
$$
を表現する。『代数スタック』の補題
\ref{algebraic-lemma-representable-diagonal} を参照せよ。
$2$-ファイバー積の構成により、$Z$ の $T$-値点は、射
$a : T \to U$, $a' : T \to U'$ と射
$\alpha : a^*x \to (a')^*\xi'$ からなる三つ組 $(a, a', \alpha)$ に対応する。
可換図式
$$
\xymatrix{
U \ar[rd] \ar[rdd] \ar[rrd] \\
& Z \ar[r]_{p'} \ar[d]^p & U' \ar[d] \\
& U \ar[r] & S
}
$$
を得る。射 $i : U \to Z$ は同型 $x \to \xi'|_U$ から得られる。
$z_0 = i(u_0) \in Z$ とおく。補題 \ref{artin-lemma-base-change-versal} により
$Z \to U'$ は $z_0$ で滑らかである。$U$ を $u_0$ のアフィン開近傍に、$U'$ を
対応する開集合に、$Z$ を $p$ と $p'$ によるこれらの開集合の逆像の共通部分に
置き換えると、$Z \to U'$ が $i(U)$ に沿って滑らかな状況になる。これはまた、
$u_n$ を部分列、すなわちその開集合に属する $u_n$ の添字だけからなる部分列へ
置き換えることも伴う。さらに、補題 \ref{artin-lemma-infinite-sequence-pre} の条件 (3) は
$U$ の縮小で明らかに保たれる（スキーム $U$, $U_n$, $U'$ はすべて同じ台位相空間を
もつ）。$U \to U'$ はアフィンスキームの一階無限小厚化なので、
$p' \circ i' = \text{id}_{U'}$ を満たし、$U$ への制限が $i$ となる射
$i' : U' \to Z$ を選べる（『空間の射についてさらに』の補題
\ref{spaces-more-morphisms-lemma-smooth-formally-smooth}）。普遍射
$p^*x \to (p')^*\xi'$ を $i'$ で引き戻すと、$p \circ i' : U' \to U$ 上の射
$$
\xi' \to x
$$
で、その合成
$$
x \to \xi' \to x
$$
が恒等射であるものを得る。射 $U_1 \to U'$ 上の $x_1 \to \xi'$ があることを
思い出そう。合成して射 $x_1 \to x$ を得るが、その存在は補題
\ref{artin-lemma-infinite-sequence-pre} の条件 (3) に矛盾する。この矛盾により証明が
完了する。
\end{proof}

\begin{remark}
\label{artin-remark-trade-openness-versality-diagonal-with-strong-effectiveness}
グルーポイドをファイバーとする圏の対角射の半普遍性の開性を、強い形式的有効性
公理から導く方法がある。$S$ を局所 Noether スキームとし、$\mathcal{X}$ を
$(\Sch/S)_{fppf}$ 上のグルーポイドをファイバーとする圏とする。次を仮定する。
\begin{enumerate}
\item $\Delta_\Delta : \mathcal{X} \to
\mathcal{X} \times_{\mathcal{X} \times \mathcal{X}} \mathcal{X}$ は代数空間により
表現可能である。
\item $\mathcal{X}$ は (RS*) をもつ。
\item $\mathcal{X}$ は極限を保つ。
\item 注意 \ref{artin-remark-strong-effectiveness} のような $S$-代数の逆系 $(R_n)$ で、
すべての $m \geq n$ に対して $\Ker(R_m \to R_n)$ が平方零イデアルであるものが
与えられたとき、関手
$$
\mathcal{X}_{\Spec(\lim R_n)} \longrightarrow
\lim_n \mathcal{X}_{\Spec(R_n)}
$$
は充満忠実である。
\end{enumerate}
このとき $\Delta : \mathcal{X} \to \mathcal{X} \times \mathcal{X}$ は半普遍性の
開性を満たす。これは、$V$ が $S$ 上局所有限表示な任意のアフィンスキームで、
$y$ が $V$ 上の $\mathcal{X} \times \mathcal{X}$ の対象であるとき、補題
\ref{artin-lemma-SGE-implies-openness-versality} を
$\mathcal{X} \times_{\Delta, \mathcal{X} \times \mathcal{X}, y}
(\Sch/V)_{fppf}$ の形のファイバー積へ適用することから従う。
これが必要になれば、この注意を補題に変え、詳細な証明を与える。
\end{remark}






\section{無限小変形}
\label{artin-section-inf}

\noindent
本節では、節 \ref{artin-section-tangent-spaces} で導入した接空間の概念の一般化を論じる。
これを適切に行うため、『形式的変形理論』の節
\ref{formal-defos-section-tangent-spaces-functors},
\ref{formal-defos-section-lifts},
\ref{formal-defos-section-infinitesimal-automorphisms} から記法を借りる。

\medskip\noindent
$S$ をスキームとする。$\mathcal{X}$ を $(\Sch/S)_{fppf}$ 上の
グルーポイドをファイバーとする圏とする。$S$-代数の準同型 $A' \to A$ と
$\Spec(A)$ 上の $\mathcal{X}$ の対象 $x$ が与えられたとき、$x$ の
$\Spec(A')$ への持ち上げの圏を $\textit{Lift}(x, A')$ と書く。
$\textit{Lift}(x, A')$ の対象とは、$\Spec(A) \to \Spec(A')$ 上にある
$\mathcal{X}$ の射 $x \to x'$ であり、$\textit{Lift}(x, A')$ の射は可換図式として
定義される。$\textit{Lift}(x, A')$ の同型類の集合を $\text{Lift}(x, A')$ と表す。
『形式的変形理論』の定義 \ref{formal-defos-definition-lifts} および注意
\ref{formal-defos-remark-omit-arrow} を参照せよ。
$A' \to A$ が局所冪零な核をもつ全射であるとき、$\text{Lift}(x, A')$ の元
$x'$ を $x$ の {\it （無限小）変形} という。この場合、
{\it $x$ 上の $x'$ の無限小自己同型群} は核
$$
\text{Inf}(x'/x) =
\Ker\left(
\text{Aut}_{\mathcal{X}_{\Spec(A')}}(x') \to
\text{Aut}_{\mathcal{X}_{\Spec(A)}}(x)\right)
$$
である。$\text{Inf}(x'/x)$ の元は、$\mathit{Aut}_\mathcal{X}(x')$ に付随する
集合をファイバーとする圏に対する、$\Spec(A')$ 上の $\text{id}_x$ の持ち上げと
同じものである。『形式的変形理論』の定義
\ref{formal-defos-definition-relative-infinitesimal-auts} および注意
\ref{formal-defos-remark-infaut-lifting-equalities} と比較せよ。

\medskip\noindent
$M$ が $A$-加群ならば、台となる $A$-加群が $A \oplus M$ で、積が
$(a, m) \cdot (a', m') = (aa', am' + a'm)$ で与えられる $A$-代数を $A[M]$ と
表す。$M = A$ のとき、これは $A$ 上の双対数環であり、慣例どおり
$A[\epsilon]$ と表す。$A$-代数準同型 $A[M] \to A$ がある。
$x$ の $\Spec(A[M])$ への引き戻しを、$x$ の $\Spec(A[M])$ への
{\it 自明な変形} という。

\begin{lemma}
\label{artin-lemma-functoriality}
$S$ をスキームとする。$f : \mathcal{X} \to \mathcal{Y}$ を
$(\Sch/S)_{fppf}$ 上のグルーポイドをファイバーとする圏の $1$-射とする。
$S$-代数の可換図式
$$
\xymatrix{
B' \ar[r] & B \\
A' \ar[u] \ar[r] & A \ar[u]
}
$$
をとる。$x$ を $\Spec(A)$ 上の $\mathcal{X}$ の対象、$y$ を $\Spec(B)$ 上の
$\mathcal{Y}$ の対象とし、$\phi : f(x)|_{\Spec(B)} \to y$ を $\Spec(B)$ 上の
$\mathcal{Y}$ の射とする。このとき、$f$ と $\phi$ により誘導される持ち上げの圏の
標準関手
$$
\textit{Lift}(x, A') \longrightarrow \textit{Lift}(y, B')
$$
がある。この構成は、明らかな仕方でグルーポイドをファイバーとする圏の $1$-射の
合成と両立する。
\end{lemma}

\begin{proof}
この補題は自明である。
\end{proof}

\noindent
$S$ を基礎スキームとする。$\mathcal{X}$ を $(\Sch/S)_{fppf}$ 上の
グルーポイドをファイバーとする圏とする。対象が組 $(x, A' \to A)$ で、
\begin{enumerate}
\item $A' \to A$ は、核が平方零イデアルである $S$-代数の全射であり、
\item $x$ は $\Spec(A)$ 上にある $\mathcal{X}$ の対象である
\end{enumerate}
ような圏を定義する。射 $(y, B' \to B) \to (x, A' \to A)$ は、$S$-代数の可換図式
$$
\xymatrix{
B' \ar[r] & B \\
A' \ar[u] \ar[r] & A \ar[u]
}
$$
と、$\Spec(B)$ 上の射 $x|_{\Spec(B)} \to y$ により与えられる。これを
{\it 変形状況} の圏と呼ぶ。

\begin{lemma}
\label{artin-lemma-properties-lift-RS-star}
$S$ をスキームとする。$\mathcal{X}$ を $(\Sch/S)_{fppf}$ 上の
グルーポイドをファイバーとする圏とし、$\mathcal{X}$ は条件 (RS*) を満たすと
仮定する。$A$ を $S$-代数とし、$x$ を $\Spec(A)$ 上の $\mathcal{X}$ の対象とする。
\begin{enumerate}
\item $A$-線形関手 $\text{Inf}_x : \text{Mod}_A \to \text{Mod}_A$ が存在し、
変形状況 $(x, A' \to A)$ と持ち上げ $x'$ が与えられたとき、
$I = \Ker(A' \to A)$ に対して同型
$\text{Inf}_x(I) \to \text{Inf}(x'/x)$ がある。
\item $A$-線形関手 $T_x : \text{Mod}_A \to \text{Mod}_A$ が存在し、次を満たす。
\begin{enumerate}
\item $\text{Mod}_A$ の $M$ が与えられたとき、全単射
$T_x(M) \to \text{Lift}(x, A[M])$ がある。
\item 変形状況 $(x, A' \to A)$ が与えられたとき、作用
$$
T_x(I) \times \text{Lift}(x, A') \to \text{Lift}(x, A')
$$
がある。ここで $I = \Ker(A' \to A)$ である。
$\text{Lift}(x, A') \not = \emptyset$ ならば、この作用は単純推移的である。
\end{enumerate}
\end{enumerate}
\end{lemma}

\begin{proof}
$\text{Inf}_x$ を関手
$$
\text{Mod}_A \longrightarrow \textit{Sets},\quad
M \longrightarrow
\text{Inf}(x'_M/x) = \text{Lift}(\text{id}_x, A[M])
$$
として定義する。これは $M$ を、$x$ の $\Spec(A[M])$ への自明な変形 $x'_M$ の
無限小自己同型群、または同じことだが $\mathit{Aut}_\mathcal{X}(x'_M)$ における
$\text{id}_x$ の持ち上げの群へ写す。$T_x$ を、$x$ の $\Spec(A[M])$ への無限小
変形の同型類の関手
$$
\text{Mod}_A \longrightarrow \textit{Sets},\quad
M \longrightarrow \text{Lift}(x, A[M])
$$
として定義する。『形式的変形理論』の補題
\ref{formal-defos-lemma-linear-functor} を $\text{Inf}_x$ と $T_x$ に適用する。
(RS*) により、圏として
$$
\textit{Lift}(x, A[M \times N]) =
\textit{Lift}(x, A[M]) \times \textit{Lift}(x, A[N])
$$
であり、自明な変形も対応するので、この補題を適用できる。

\medskip\noindent
$(x, A' \to A)$ を変形状況とする。規則
$g(a_1, a_2) = \overline{a_1} \oplus a_2 - a_1$ で定まる環準同型
$g : A' \times_A A' \to A[I]$ を考える。
$(a_1, a_2) \mapsto (a_1, g(a_1, a_2))$ で与えられる同型
$$
A' \times_A A' \longrightarrow A' \times_A A[I]
$$
がある。この同型は第一因子上の $A'$ への射影、したがって $A$ への射影と可換する。
よって (RS*) を二度適用すると、圏同値
\begin{align*}
\textit{Lift}(x, A') \times \textit{Lift}(x, A')
& =
\textit{Lift}(x, A' \times_A A') \\
& =
\textit{Lift}(x, A' \times_A A[I]) \\
& =
\textit{Lift}(x, A') \times \textit{Lift}(x, A[I])
\end{align*}
を得る。これらの写像と最後の積の最終因子への射影を用いると、「差写像」
$$
\text{Inf}(x'/x) \times  \text{Inf}(x'/x)
\longrightarrow
\text{Inf}_x(I)
\quad\text{and}\quad
\text{Lift}(x, A') \times \text{Lift}(x, A')
\longrightarrow
T_x(I)
$$
を得る。これらの差写像は、図式
$$
\xymatrix{
A' \times_A A' \times_A A'
\ar[rrrrr]_-{(a_1, a_2, a_3) \mapsto (g(a_1, a_2), g(a_2, a_3))}
\ar[rrrrrd]_{(a_1, a_2, a_3) \mapsto g(a_1, a_3)} & & & & &
A[I] \times_A A[I] = A[I \times I] \ar[d]^{+} \\
& & & & & A[I]
}
$$
が可換であることから、推移律
「$(x'_1 - x'_2) + (x'_2 - x'_3) = x'_1 - x'_3$」を満たす。
上の圏同値の列を逆にすると、$\text{Inf}(x'/x)$、それぞれ
$\text{Lift}(x, A')$ が空でなければ自由かつ推移的な作用を得る。
$\text{Inf}(x'/x)$ は群なので常に空でないことに注意せよ。
\end{proof}

\begin{remark}[関手性]
\label{artin-remark-functoriality}
補題 \ref{artin-lemma-properties-lift-RS-star} と同じ仮定と記法を用いる。
$A \to B$ を環準同型とし、$y = x|_{\Spec(B)}$ とする。
$M \in \text{Mod}_A$, $N \in \text{Mod}_B$ とし、$M \to N$ を
$A$-線形写像とする。このとき、環準同型 $A[M] \to B[N]$ から得られる引き戻し
関手
$$
\textit{Lift}(x, A[M]) \to \textit{Lift}(y, B[N])
$$
があるので、標準写像 $\text{Inf}_x(M) \to \text{Inf}_y(N)$ および
$T_x(M) \to T_y(N)$ がある。同様に、変形状況の射
$(y, B' \to B) \to (x, A' \to A)$ が与えられると、引き戻し関手
$\textit{Lift}(x, A') \to \textit{Lift}(y, B')$ を得る。
$\text{Inf}_x$ と $T_x$ 上の作用、加法、スカラー倍の構成は持ち上げの圏における
射だけを用いるので（『形式的変形理論』の補題
\ref{formal-defos-lemma-linear-functor} の証明を参照せよ）、上の構成は関手的である。
言い換えると、$A$-線形写像
$$
\text{Inf}_x(M) \to \text{Inf}_y(N)
\quad\text{and}\quad
T_x(M) \to T_y(N)
$$
で、図式
$$
\vcenter{
\xymatrix{
\text{Inf}_y(J) \ar[r] & \text{Inf}(y'/y) \\
\text{Inf}_x(I) \ar[r] \ar[u] & \text{Inf}(x'/x) \ar[u]
}
}
\quad\text{and}\quad
\vcenter{
\xymatrix{
T_y(J) \times \text{Lift}(y, B') \ar[r] & \text{Lift}(y, B') \\
T_x(I) \times \text{Lift}(x, A') \ar[r] \ar[u] & \text{Lift}(x, A') \ar[u]
}
}
$$
を可換にするものを得る。ここで $I = \Ker(A' \to A)$,
$J = \Ker(B' \to B)$ であり、$x'$ は $x$ の $A'$ への持ち上げ
（常に存在するとは限らない）、$y' = x'|_{\Spec(B')}$ である。
\end{remark}

\begin{remark}[自己同型]
\label{artin-remark-automorphisms}
補題 \ref{artin-lemma-properties-lift-RS-star} と同じ仮定と記法を用いる。
$x', x''$ を $x$ の $A'$ への持ち上げとする。このとき合成写像
$$
\text{Inf}(x'/x) \times
\Mor_{\textit{Lift}(x, A')}(x', x'') \times \text{Inf}(x''/x)
\longrightarrow
\Mor_{\textit{Lift}(x, A')}(x', x'').
$$
がある。$\textit{Lift}(x, A')$ はグルーポイドなので、
$\Mor_{\textit{Lift}(x, A')}(x', x'')$ が空でなければ、これは
$\Mor_{\textit{Lift}(x, A')}(x', x'')$ 上の $\text{Inf}(x'/x)$ の単純推移的な
左作用と、$\text{Inf}(x''/x)$ による単純推移的な右作用を定める。
さて補題によれば
$\text{Inf}(x'/x) = \text{Inf}_x(I) = \text{Inf}(x''/x)$ である。
これらの同一視を介して、上で述べた二つの作用は一致すると主張する。
実際、$x' \not \cong x''$ ならば主張は明らかであり、または
$x' \cong x''$ ならば $x'' = x'$ と仮定してよく、この場合は
$\text{Inf}(x'/x)$ が可換であることから結果が従う。特に、良定義な作用
$$
\text{Inf}_x(I) \times \Mor_{\textit{Lift}(x, A')}(x', x'')
\longrightarrow
\Mor_{\textit{Lift}(x, A')}(x', x'')
$$
を得る。この作用は $\Mor_{\textit{Lift}(x, A')}(x', x'')$ が空でなければ直ちに
単純推移的である。
\end{remark}

\begin{remark}
\label{artin-remark-short-exact-sequence-thickenings}
$S$ をスキームとする。$\mathcal{X}$ を $(\Sch/S)_{fppf}$ 上の
グルーポイドをファイバーとする圏とする。$A$ を $S$-代数とする。
変形状況の {\it 短完全列}
$$
(x, A_1' \to A) \to (x, A_2' \to A) \to (x, A_3' \to A)
$$
という概念がある。対応する核 $I_i = \Ker(A_i' \to A)$ の間の写像が
$A$-加群の短完全列
$$
0 \to I_3 \to I_2 \to I_1 \to 0
$$
を与えることを要求する。この場合、写像 $A_3' \to A_1'$ は $A$ を経由するので、
標準的な同型 $A_1' = A[I_1]$ があることに注意せよ。
\end{remark}

\begin{lemma}
\label{artin-lemma-ses-inf-and-T}
$S$ をスキームとする。$p : \mathcal{X} \to \mathcal{Y}$ と
$q : \mathcal{Z} \to \mathcal{Y}$ を、$(\Sch/S)_{fppf}$ 上の
グルーポイドをファイバーとする圏の $1$-射とする。
$\mathcal{X}$, $\mathcal{Y}$, $\mathcal{Z}$ が (RS*) を満たすと仮定する。
$A$ を $S$-代数とし、$w$ を $A$ 上の
$\mathcal{W} = \mathcal{X} \times_\mathcal{Y} \mathcal{Z}$ の対象とする。
$w$ から得られる $\mathcal{X}, \mathcal{Y}, \mathcal{Z}$ の対象を
$x, y, z$ と表す。任意の $A$-加群 $M$ に対して、$A$-加群の $6$-項完全列
$$
\xymatrix{
0 \ar[r] &
\text{Inf}_w(M) \ar[r] &
\text{Inf}_x(M) \oplus \text{Inf}_z(M) \ar[r] &
\text{Inf}_y(M) \ar[lld] \\
 &
T_w(M) \ar[r] &
T_x(M) \oplus T_z(M) \ar[r] &
T_y(M)
}
$$
がある。
\end{lemma}

\begin{proof}
補題 \ref{artin-lemma-fibre-product-RS-star} により $\mathcal{W}$ は (RS*) を満たすので、
$T_w(M)$ と $\text{Inf}_w(M)$ は定義される。水平矢印は補題
\ref{artin-lemma-functoriality} の関手性を用いて定義される。

\medskip\noindent
「境界」写像 $\delta : \text{Inf}_y(M) \to T_w(M)$ を定義する。
『圏』の補題 \ref{categories-lemma-2-product-fibred-categories}、より正確には
『圏』の補題 \ref{categories-lemma-2-product-categories-over-C} における
$2$-ファイバー積の記述で
$w = (x, z, p(x) \to y \to q(z))$ となるような同型
$p(x) \to y$ と $y \to q(z)$ を選ぶ。
$x, y, z, w$ の $A[M]$ 上の自明な変形をそれぞれ $x', y', z', w'$ と表す。
引き戻しにより同型 $y' \to p(x')$ と $q(z') \to y'$ を得る。
$\alpha \in \text{Inf}_y(M)$ の元は、$A[M]$ 上の自己同型
$\alpha : y' \to y'$ で、$A$ 上の $y$ への制限が恒等射となるものと同じである。
したがって
$$
\delta(\alpha) =
(x', z', p(x') \to y' \xrightarrow{\alpha} y' \to q(z'))
$$
とおくことにより、$T_w(M)$ の対象を得る。これは『形式的変形理論』の補題
\ref{formal-defos-lemma-morphism-linear-functors} により $A$-加群の写像である。

\medskip\noindent
証明の残りは『形式的変形理論』の補題
\ref{formal-defos-lemma-deformation-categories-fiber-product-morphisms} の証明と
まったく同じである。
\end{proof}

\begin{remark}[従来の接空間との両立性]
\label{artin-remark-compare-deformation-spaces}
$S$ を局所 Noether スキームとする。$\mathcal{X}$ を
$(\Sch/S)_{fppf}$ 上のグルーポイドをファイバーとする圏とし、$\mathcal{X}$ は
(RS*) をもつと仮定する。$k$ を $S$ 上有限型な体とし、$x_0$ を
$\Spec(k)$ 上の $\mathcal{X}$ の対象とする。このとき $k$-ベクトル空間の等式
$$
T\mathcal{F}_{\mathcal{X}, k, x_0} = T_{x_0}(k)
\quad\text{and}\quad
\text{Inf}(\mathcal{F}_{\mathcal{X}, k, x_0}) =
\text{Inf}_{x_0}(k)
$$
がある。等号の左辺の空間は (\ref{artin-equation-tangent-space}) と
(\ref{artin-equation-infinitesimal-automorphisms}) で与えられ、右辺の空間は補題
\ref{artin-lemma-properties-lift-RS-star} で与えられる。
\end{remark}

\begin{remark}[標準元]
\label{artin-remark-canonical-element}
補題 \ref{artin-lemma-properties-lift-RS-star} と同じ仮定と記法を用いる。
$\Spec(A) \to S$ が環準同型 $\Lambda \to A$ に対応するようなアフィン開集合
$\Spec(\Lambda) \subset S$ を選ぶ。環準同型
$$
A \longrightarrow A[\Omega_{A/\Lambda}],
\quad
a \longmapsto (a, \text{d}_{A/\Lambda}(a))
$$
を考える。対応する射 $\Spec(A[\Omega_{A/\Lambda}]) \to \Spec(A)$ に沿って
$x$ を引き戻し、$A[\Omega_{A/\Lambda}]$ 上の $x$ の変形 $x_{can}$ を得る。
これを {\it 標準元}
$$
x_{can} \in T_x(\Omega_{A/\Lambda}) = \text{Lift}(x, A[\Omega_{A/\Lambda}]).
$$
と呼ぶ。次に、$\Lambda$ は Noether で $\Lambda \to A$ は有限型であると仮定する。
$k = \kappa(\mathfrak p)$ を $U = \Spec(A)$ の有限型点 $u_0$ における剰余体とする。
$x_0 = x|_{u_0}$ とおく。(RS*) と $A[k] = A \times_k k[k]$ であることにより、
空間 $T_x(k)$ は変形関手 $\mathcal{F}_{\mathcal{X}, k, x_0}$ の接空間である。
$$
T\mathcal{F}_{U, k, u_0} =
\text{Der}_\Lambda(A, k) = \Hom_A(\Omega_{A/\Lambda}, k)
$$
による同一視（『形式的変形理論』の例
\ref{formal-defos-example-tangent-space-prorepresentable-functor} を参照せよ）と
$T_x$ の関手性を介して、標準元は $U$ 上の対象 $x$ により誘導される接空間上の
写像を生じる。すなわち $\theta \in T\mathcal{F}_{U, k, u_0}$ は
$T_x(k) = T\mathcal{F}_{\mathcal{X}, k, x_0}$ における
$T_x(\theta)(x_{can})$ へ写る。
\end{remark}

\begin{remark}[標準自己同型]
\label{artin-remark-canonical-isomorphism}
$S$ を局所 Noether スキームとする。$\mathcal{X}$ を
$(\Sch/S)_{fppf}$ 上のグルーポイドをファイバーとする圏とし、$\mathcal{X}$ は
条件 (RS*) を満たすと仮定する。$A$ を $S$-代数で、$\Spec(A) \to S$ が
アフィン開集合へ写るものとし、$x, y$ を $\Spec(A)$ 上の $\mathcal{X}$ の対象とする。
さらに $A \to B$ を環準同型とし、
$\alpha : x|_{\Spec(B)} \to y|_{\Spec(B)}$ を $\Spec(B)$ 上の $\mathcal{X}$ の射とする。
環準同型
$$
B \longrightarrow B[\Omega_{B/A}],
\quad
b \longmapsto (b, \text{d}_{B/A}(b))
$$
を考える。対応する射 $\Spec(B[\Omega_{B/A}]) \to \Spec(B)$ に沿って
$\alpha$ を引き戻し、$B[\Omega_{B/A}]$ 上の $x$ と $y$ の引き戻しの間の射
$\alpha_{can}$ を得る。一方、$B$ の $B[\Omega_{B/A}]$ の第一直和因子への埋め込みに
対応する射 $\Spec(B[\Omega_{B/A}]) \to \Spec(B)$ によって $\alpha$ を引き戻すことも
できる。注意 \ref{artin-remark-automorphisms} の議論により、差
$$
\varphi(x, y, \alpha) = \alpha_{can} - \alpha|_{\Spec(B[\Omega_{B/A}])} \in
\text{Inf}_{x|_{\Spec(B)}}(\Omega_{B/A}).
$$
をとれる。これを {\it 標準自己同型} と呼ぶ。これはすべてのデータ
$A$, $x$, $y$, $A \to B$, $\alpha$ に依存する。
\end{remark}

\section{障害理論}
\label{artin-section-obstruction-theory}

\noindent
本節では、障害理論とは何かを述べる。無限小変形および無限小自己同型の空間とは
異なり、障害理論は付加的なデータである。この定式化は、補題
\ref{artin-lemma-properties-lift-RS-star} と注意 \ref{artin-remark-functoriality} の結果に
動機づけられている。

\begin{definition}
\label{artin-definition-obstruction-theory}
$S$ を局所 Noether 的基礎とする。$\mathcal{X}$ を $(\Sch/S)_{fppf}$ 上の
グルーポイドをファイバーとする圏とする。{\it 障害理論}とは、次のデータからなる：
\begin{enumerate}
\item $\Spec(A) \to S$ があるアフィン開集合へ写るような任意の $S$-代数 $A$ と、
$\Spec(A)$ 上の $\mathcal{X}$ の任意の対象 $x$ に対し、{\it 障害加群}を与える
$A$-線形関手
$$
\mathcal{O}_x : \text{Mod}_A \to \text{Mod}_A
$$
\item (1) における $(x, A)$、環準同型 $A \to B$、
$M \in \text{Mod}_A$、$N \in \text{Mod}_B$、および $A$-線形写像
$M \to N$ に対し、誘導される $A$-線形写像
$\mathcal{O}_x(M) \to \mathcal{O}_y(N)$。ここで $y = x|_{\Spec(B)}$ とする；さらに
\item 任意の変形状況 $(x, A' \to A)$ に対し、{\it 障害}元
$o_x(A') \in \mathcal{O}_x(I)$。ここで $I = \Ker(A' \to A)$ とする。
\end{enumerate}
これらのデータには、次の条件を課す：
\begin{enumerate}
\item[(i)] 関手性を与える写像により、障害加群は三つ組 $(x, A, M)$ の圏から
集合の圏への関手をなす；
\item[(ii)] 変形状況の任意の射
$(y, B' \to B) \to (x, A' \to A)$ に対して、元 $o_x(A')$ は
$o_y(B')$ へ写る；
\item[(iii)] 任意の変形状況 $(x, A' \to A)$ に対して
$$
\text{Lift}(x, A') \not = \emptyset
\Leftrightarrow
o_x(A') = 0
$$
が成り立つ。
\end{enumerate}
\end{definition}

\noindent
最後の条件がこの用語を説明している。加群 $\mathcal{O}_x(A')$ を
{\it 障害加群}と呼び、元 $o_x(A')$ を {\it 障害}と呼ぶ。
多くの障害理論はさらに別の性質をもち、それらを有用なものとするには追加の条件が
必要である。さらに、これは一つの例示的な定義にすぎない。たとえば定義において、
$S$ 上有限型な変形状況のみを考えることもできる。

\medskip\noindent
障害理論を導入する主な理由の一つは、半普遍性の開性を確認することである。
この種の結果の一例が、下の補題 \ref{artin-lemma-get-openness-obstruction-theory} である。
これを行う最初の着想は Artin によるものであり、序論で挙げた Artin の論文を参照せよ。
その後、たとえば Flenner \cite{Flenner}、Hall \cite{Hall-coherent}、
Hall と Rydh \cite{rydh_axioms}、Olsson \cite{olsson_deformation}、
Olsson と Starr \cite{olsson-starr}、および Lieblich
\cite{lieblich-complexes} の研究に受け継がれている（参照文献は無作為な順序である）。
さらに、特定のグルーポイドをファイバーとする圏については、しばしば著者が
当面の問題に適応させた理論を多少展開する。われわれも後でこの理論を展開する
（ここに将来の参照を挿入する）。

\begin{lemma}
\label{artin-lemma-get-openness-obstruction-theory}
\begin{reference}
これは \cite[Theorem 4.4]{Hall-coherent} である。
\end{reference}
$S$ を局所 Noether スキームとする。$\mathcal{X}$ を $(\Sch/S)_{fppf}$ 上の
グルーポイドをファイバーとする圏とする。次を仮定する：
\begin{enumerate}
\item $\Delta : \mathcal{X} \to \mathcal{X} \times \mathcal{X}$ は
代数空間で表現可能である；
\item $\mathcal{X}$ は (RS*) をもつ；
\item $\mathcal{X}$ は極限を保つ；
\item 障害理論が存在する\footnote{証明を調べると、実際には定義
\ref{artin-definition-obstruction-theory} における障害類の関手性 (ii) を、
$B = A$ かつ $y = x$ である写像
$(y, B' \to B) \to (x, A' \to A)$ に対して確認すれば十分であるとわかる。}；
\item $\Spec(A)$ 上の $\mathcal{X}$ の対象 $x$ と $A$-加群 $M_n$、$n \geq 1$ に対して
\begin{enumerate}
\item $T_x(\prod M_n) = \prod T_x(M_n)$；
\item $\mathcal{O}_x(\prod M_n) \to \prod \mathcal{O}_x(M_n)$ は単射である。
\end{enumerate}
\end{enumerate}
このとき $\mathcal{X}$ は半普遍性の開性を満たす。
\end{lemma}

\begin{proof}
補題 \ref{artin-lemma-SGE-implies-openness-versality} の条件 (4) を確認することにより証明する。
$(\xi_n)$ と $(R_n)$ を注意 \ref{artin-remark-strong-effectiveness} におけるものとし、
すべての $m \geq n$ に対して $\Ker(R_m \to R_n)$ が平方零イデアルであると仮定する。
$A = R_1$ および $x = \xi_1$ とおく。$M_n = \Ker(R_n \to R_1)$ と書く。
このとき $M_n$ は $A$-加群である。$R = \lim R_n$ とおく。
$$
\tilde R = \{(r_1, r_2, r_3 \ldots) \in \prod R_n
\text{ such that all have the same image in }A\}
$$
とおく。このとき $\tilde R \to A$ は全射で、その核は $M = \prod M_n$ である。
写像 $R \to \tilde R$ と写像
$\tilde R \to A[M]$、$(r_1, r_2, r_3, \ldots) \mapsto
(r_1, r_2 - r_1, r_3 - r_2, \ldots)$ がある。これらを合わせると、変形状況の短完全列
$$
(x, R \to A) \to (x, \tilde R \to A) \to (x, A[M])
$$
を得る。注意 \ref{artin-remark-short-exact-sequence-thickenings} を参照せよ。
対応する核の列 $0 \to \lim M_n \to M \to M \to 0$ は、加群系 $(M_n)$ の極限を
計算する標準列である。

\medskip\noindent
$o_x(\tilde R) \in \mathcal{O}_x(M)$ を障害元とする。持ち上げ $\xi_n$ があるので、
$o_x(\tilde R)$ は $\mathcal{O}_x(M_n)$ において零へ写る。仮定 (5)(b) により
$o_x(\tilde R) = 0$ である。$x$ の $\Spec(\tilde R)$ への持ち上げ
$\tilde \xi$ を一つ選ぶ。$\tilde \xi_n$ を $\tilde \xi$ の $\Spec(R_n)$ への制限とする。
補題 \ref{artin-lemma-properties-lift-RS-star} の (2)(b) により、
$t_n \cdot \tilde \xi_n = \xi_n$ となる元 $t_n \in T_x(M_n)$ が存在する。
仮定 (5)(a) により、$T_x(M_n)$ において $t_n$ へ写る $t \in T_x(M)$ をとれる。
$\tilde \xi$ を $t \cdot \tilde \xi$ で置き換えれば、すべての $n$ に対して
$\tilde \xi$ の $\Spec(R_n)$ への制限が $\xi_n$ となる。
特に、$\xi_{n + 1}$ の $\Spec(R_n)$ への制限は $\xi_n$ なので、
$\tilde \xi$ の $\Spec(A[M])$ への制限 $\overline{\xi}$ は、すべての $n$ に対して
$\Spec(A[M_n])$ 上で自明変形へ制限される。したがって仮定 (5)(a) により、
$\overline{\xi}$ は $x$ の自明変形である。$R = \tilde R \times_{A[M]} A$ に適用した
公理 (RS*) により、$\tilde \xi$ は $R$ 上の $x$ の変形 $\xi$ の引き戻しである。
これで証明が完了する。
\end{proof}

\begin{example}
\label{artin-example-global-sections}
$S = \Spec(\Lambda)$ とし、ここで $\Lambda$ は Noether 環とする。
$W \to S$ をスキームの射とする。$\mathcal{F}$ を $S$ 上平坦な準連接
$\mathcal{O}_W$-加群とする。関手
$$
F : (\Sch/S)_{fppf}^{opp} \longrightarrow \textit{Sets},
\quad
T/S \longrightarrow H^0(W_T, \mathcal{F}_T)
$$
を考える。ここで $W_T = T \times_S W$ は基底変換であり、$\mathcal{F}_T$ は
$\mathcal{F}$ の $W_T$ への引き戻しである。$T = \Spec(A)$ のときは
$W_T = W_A$ などと書く。$F$ に付随するグルーポイドをファイバーとする圏を
$\mathcal{X} \to (\Sch/S)_{fppf}$ とする。このとき $\mathcal{X}$ は障害理論をもつ。
実際、
\begin{enumerate}
\item $\Lambda$ 上の $A$ と $x \in H^0(W_A, \mathcal{F}_A)$ が与えられたとき、
$\mathcal{O}_x(M) = H^1(W_A, \mathcal{F}_A \otimes_A M)$ とおく；
\item 変形状況 $(x, A' \to A)$ が与えられたとき、$o_x(A') \in \mathcal{O}_x(A)$ を、
境界写像
$$
H^0(W_A, \mathcal{F}_A) \longrightarrow H^1(W_A, \mathcal{F}_A \otimes_A I)
$$
による $x$ の像とする。この境界写像は短完全列
$$
0 \to \mathcal{F}_A \otimes_A I \to
\mathcal{F}_{A'} \to \mathcal{F}_A \to 0.
$$
から生じる。
\end{enumerate}
いくつかの詳細は省略した。特に、上の短完全列の構成（この構成では $W_A$ と
$W_{A'}$ が同じ台位相空間をもつことを用いる）と、$\mathcal{F}$ が $S$ 上平坦で
あることから上の列が短完全になる理由の説明を省略した。
\end{example}

\begin{example}[鍵となる例]
\label{artin-example-key}
$S = \Spec(\Lambda)$ とし、ここで $\Lambda$ は Noether 環とする。
$X = \Spec(R)$ および $R = \Lambda[x_1, \ldots, x_n]/J$ として
$\mathcal{X} = (\Sch/X)_{fppf}$ とする。素朴な余接複体 $\NL_{R/\Lambda}$ は
（標準的に）
$$
J/J^2
\longrightarrow
\bigoplus\nolimits_{i = 1, \ldots, n} R\text{d}x_i,
$$
とホモトピー同値である。代数の補題 \ref{algebra-lemma-NL-homotopy} を参照せよ。
変形状況 $(x, A' \to A)$ を考え、$A' \to A$ の核を $I$ と書く。
対象 $x$ は、すべての $f \in J$ に対して $f(a_1, \ldots, a_n) = 0$ が $A$ において
成り立つような $a_i \in A$ からなる $(a_1, \ldots, a_n)$ に対応する。次のようにおく：
\begin{align*}
\mathcal{O}_x(A')
& =
\Hom_R(J/J^2, I)/\Hom_R(R^{\oplus n}, I) \\
& =
\Ext^1_R(\NL_{R/\Lambda}, I) \\
& =
\Ext^1_A(\NL_{R/\Lambda} \otimes_R A, I).
\end{align*}
$A$ における $a_i$ の持ち上げ $a_i' \in A'$ を選ぶ。このとき $o_x(A')$ は、
$f \in J$ を $f(a_1', \ldots, a'_n) \in I$ へ送る写像 $J/J^2 \to I$ の類である。
$o_x(A')$ が選択によらないことの確認は省略する。$o_x(A') = 0$ ならば写像が
持ち上がることは明らかである。最後に、関手性は直ちにわかる。
こうして障害理論を得る。$o_x(A')$ は、$D(A)$ における合成
$$
\NL_{R/\Lambda} \to \NL_{A/\Lambda} \to \NL_{A/A'} = I[1]
$$
として、もう少し標準的に記述できることに注意する。最後の同一視については、
代数の補題 \ref{algebra-lemma-NL-surjection} を参照せよ。
\end{example}

\section{素朴な障害理論}
\label{artin-section-naive-obstruction-theory}

\noindent
本節の表題は、本節で素朴な余接複体を用いることに由来する。局所 Noether スキーム
$S$ 上の、与えられたグルーポイドをファイバーとする圏について、
$(x, A' \to A)$ を変形状況とする。鍵となる例 \ref{artin-example-key} は、任意の障害理論が
$D(A)$ における $A$ の素朴な余接複体を終域とする射と密接に関係すべきことを示唆する。
これを詳しく調べると、補題 \ref{artin-lemma-characterize-versal} において半普遍性の判定条件を
得て、それは補題 \ref{artin-lemma-openness} における半普遍性の開性の判定条件を導く。
この概念をもう少し形式化するため、定義 \ref{artin-definition-naive-obstruction-theory} で
素朴な障害理論という概念を導入する。

\medskip\noindent
以下では、代数の節 \ref{algebra-section-netherlander} で定義された素朴な余接複体を用いる。
特に、$A' \to A$ が平方零な核 $I$ をもつ $\Lambda$-代数の全射ならば、写像
$$
\NL_{A'/\Lambda} \to \NL_{A/\Lambda} \to \NL_{A/A'}
$$
があり、その合成は零とホモトピー同値である（代数の注意
\ref{algebra-remark-composition-homotopy-equivalent-to-zero} を参照せよ）。
ここでは完全な余接複体でなく素朴な余接複体を用いているため、これは一般には
識別三角形をなさない。ホモトピー同値 $\NL_{A/A'} \to I[1]$ がある
（これは次数 $-1$ に置かれた $I$ からなる複体である。代数の補題
\ref{algebra-lemma-NL-surjection} を参照せよ）。最後に、標準写像
$\NL_{A/\Lambda} \to \Omega_{A/\Lambda}$ があることに注意する。

\begin{lemma}
\label{artin-lemma-compute-ext-into-field}
$A \to k$ を、$k$ が体であるような環準同型とする。$E \in D^-(A)$ とする。
このとき $\Ext^i_A(E, k) = \Hom_k(H^{-i}(E \otimes^\mathbf{L} k), k)$ である。
\end{lemma}

\begin{proof}
省略する。ヒント：$E$ を自由 $A$-加群からなる上に有界な複体で置き換え、
両辺を計算せよ。
\end{proof}

\begin{lemma}
\label{artin-lemma-construct-essential-surjection}
$\Lambda \to A \to k$ を Noether 環の有限型環準同型とし、$A$ のある素イデアル
$\mathfrak p$ に対して $k = \kappa(\mathfrak p)$ とする。
$\xi : E \to \NL_{A/\Lambda}$ を $D^{-}(A)$ の射とし、
$H^{-1}(\xi \otimes^{\mathbf{L}} k)$ は全射でないと仮定する。
このとき、次を満たす $\Lambda$-代数の全射 $A' \to A$ が存在する：
\begin{enumerate}
\item[(a)] $I = \Ker(A' \to A)$ は平方零であり、$A$-加群として $k$ と同型である；
\item[(b)] $\Omega_{A'/\Lambda} \otimes k = \Omega_{A/\Lambda} \otimes k$；
\item[(c)] $E \to \NL_{A/A'}$ は零である。
\end{enumerate}
\end{lemma}

\begin{proof}
$f \in A$、$f \not \in \mathfrak p$ とする。$A'' \to A_f$ が、誘導される写像
$E \otimes_A A_f \to \NL_{A_f/\Lambda}$ に対して (a)、(b)、(c) を満たすと仮定する。
代数続論の補題 \ref{algebra-lemma-localize-NL} を参照せよ。このとき
$A' = A'' \times_{A_f} A$ とおけば解を得る。実際、
$\Ker(A' \to A) = \Ker(A'' \to A) = I$ なので、$A' \to A$ が (a) を満たすことは
明らかである。$f$ を持ち上げる $f'' \in A''$ を選ぶ。このとき $A'$ の
$(f'', f)$ における局所化は $A''$ と同型である（たとえば代数続論の補題
\ref{more-algebra-lemma-diagram-localize} による）。したがって (b) と (c) も
$A'$ に対して明らかである。このようにして、$A$ を局所化 $A_f$ で
（有限回）置き換えてよいことがわかる。特に、このような置き換えの後、
$\mathfrak p$ は $A$ の極大イデアルであると仮定してよい。射の補題
\ref{morphisms-lemma-point-finite-type} を参照せよ。

\medskip\noindent
表示 $A = \Lambda[x_1, \ldots, x_n]/J$ を選ぶ。このとき
$\NL_{A/\Lambda}$ は（標準的に）
$$
J/J^2
\longrightarrow
\bigoplus\nolimits_{i = 1, \ldots, n} A\text{d}x_i,
$$
とホモトピー同値である。代数の補題 \ref{algebra-lemma-NL-homotopy} を参照せよ。
必要なら局所化することにより（Nakayama の補題を用いる）、$f_j \otimes 1$ が
$J/J^2 \otimes_A k$ の基底をなすような $J$ の生成元
$f_1, \ldots, f_m$ を選べる。さらに番号を付け替えることにより、
$\text{d}f_1, \ldots, \text{d}f_r$ の像が
$J/J^2 \otimes k \to \bigoplus k\text{d}x_i$ の像の基底をなし、
$\text{d}f_{r + 1}, \ldots, \text{d}f_m$ が $\bigoplus k\text{d}x_i$ において
零へ写ると仮定できる。これらの選択のもとで、空間
$$
H^{-1}(\NL_{A/\Lambda} \otimes^{\mathbf{L}}_A k) =
H^{-1}(\NL_{A/\Lambda} \otimes_A k)
$$
は $f_{r + 1} \otimes 1, \ldots, f_m \otimes 1$ を基底にもつ。
もう一度基底を取り替えることにより、$H^{-1}(\xi \otimes^{\mathbf{L}} k)$ の像は
$f_{r + 1} \otimes 1, \ldots, f_{m - 1} \otimes 1$ の $k$-線形包に含まれると
仮定してよい。次のようにおく：
$$
A' = \Lambda[x_1, \ldots, x_n]/(f_1, \ldots, f_{m - 1}, \mathfrak pf_m)
$$
構成により $A' \to A$ は (a) を満たす。$\text{d}f_m$ は
$\bigoplus k\text{d}x_i$ において零へ写るので、(b) が成り立つ。最後に構成により、
誘導される写像 $E \to \NL_{A/A'} = I[1]$ は零写像
$H^{-1}(E \otimes_A^\mathbf{L} k) \to I \otimes_A k$ を誘導する。
補題 \ref{artin-lemma-compute-ext-into-field} により、この合成は零である。
\end{proof}

\noindent
次の補題がわれわれの主要な技術的結果である。

\begin{lemma}
\label{artin-lemma-characterize-versal}
$S$ を局所 Noether スキームとする。$\mathcal{X}$ を $(\Sch/S)_{fppf}$ 上の
グルーポイドをファイバーとする圏で、(RS*) を満たすものとする。
$U = \Spec(A)$ を $S$ 上有限型なアフィンスキームで、あるアフィン開集合
$\Spec(\Lambda)$ へ写るものとする。$x$ を $U$ 上の $\mathcal{X}$ の対象とする。
$\xi : E \to \NL_{A/\Lambda}$ を $D^{-}(A)$ の射とする。次を仮定する：
\begin{enumerate}
\item[(i)] 任意の変形状況 $(x, A' \to A)$ に対して、$x$ が $\Spec(A')$ へ
持ち上がることと、$E \to \NL_{A/\Lambda} \to \NL_{A/A'}$ が零であることは同値である；
\item[(ii)] 関手の同型 $T_x(-) \to \Ext^0_A(E, -)$ が存在し、そのもとで
$E \to \NL_{A/\Lambda} \to \Omega^1_{A/\Lambda}$ は標準元に対応する
（注意 \ref{artin-remark-canonical-element} を参照せよ）。
\end{enumerate}
$u_0 \in U$ を、剰余体が $k = \kappa(u_0)$ である有限型点とする。
次の条件を考える：
\begin{enumerate}
\item $x$ は $u_0$ において半普遍的である；
\item $\xi : E \to \NL_{A/\Lambda}$ は全射
$H^{-1}(E \otimes_A^{\mathbf{L}} k) \to
H^{-1}(\NL_{A/\Lambda} \otimes_A^{\mathbf{L}} k)$ と単射
$H^0(E \otimes_A^{\mathbf{L}} k) \to
H^0(\NL_{A/\Lambda} \otimes_A^{\mathbf{L}} k)$ を誘導する。
\end{enumerate}
このとき常に (2) $\Rightarrow$ (1) が成り立ち、$u_0$ が閉点ならば
(1) $\Rightarrow$ (2) も成り立つ。
\end{lemma}

\begin{proof}
$\mathfrak p = \Ker(A \to k)$ を $u_0$ に対応する素イデアルとする。

\medskip\noindent
$x$ は $u_0$ において半普遍的であり、$u_0$ は $U$ の閉点であると仮定する。
$H^{-1}(\xi \otimes_A^{\mathbf{L}} k)$ が全射でないならば、
補題 \ref{artin-lemma-construct-essential-surjection} におけるような、核が $I$ である拡大
$A' \to A$ をとる。$u_0$ は閉点なので $I$ は有限 $A$-加群であり、したがって
$A'$ は有限型 $\Lambda$-代数である（$u_0$ が閉でなければこれは成り立たない）。
特に $A'$ は Noether である。$A'$ の性質 (c) と $\xi$ の (i) により、
$x$ は $A'$ 上の対象 $x'$ へ持ち上がる。$\mathfrak p' \subset A'$ を $k$ への
全射の核とする。Artin--Rees（代数の補題 \ref{algebra-lemma-Artin-Rees}）により、
$(\mathfrak p')^n \cap I = 0$ となる $n > 1$ が存在する。このとき
$$
B' = A'/(\mathfrak p')^n \longrightarrow A/\mathfrak p^n = B
$$
は Artin 局所環の小さい本質的拡大である。形式的変形理論の補題
\ref{formal-defos-lemma-essential-surjection} を参照せよ。一方、$x$ は $u_0$ において
半普遍的であり、$x'|_{\Spec(B')}$ は $x|_{\Spec(B)}$ の持ち上げなので、
合成 $A/\mathfrak p^m \to B' \to B$ が商写像となるような整数 $m \geq n$ と写像
$q : A/\mathfrak p^m \to B'$ が存在する。$B'$ の極大イデアルの $n$ 乗は零なので、
この $q$ は $B$ を経由する。これは $B' \to B$ が本質的全射であることに反する。
この矛盾により、$H^{-1}(\xi \otimes_A^{\mathbf{L}} k)$ は全射である。

\medskip\noindent
$x$ は $u_0$ において半普遍的であると仮定する。補題
\ref{artin-lemma-compute-ext-into-field} により、写像
$H^0(\xi \otimes_A^{\mathbf{L}} k)$ は写像
$\Ext^0_A(\NL_{A/\Lambda}, k) \to \text{Ext}^0_A(E, k)$ の双対である。ここで
$$
\Ext^0_A(\NL_{A/\Lambda}, k) = \text{Der}_\Lambda(A, k)
\quad\text{and}\quad
T_x(k) = \Ext^0_A(E, k)
$$
である。条件 (ii) により、写像
$\Ext^0_A(\NL_{A/\Lambda}, k) \to \text{Ext}^0_A(E, k)$ は、
$u_0$ における $U$ への接ベクトル $\theta$ を $x_0$ の対応する無限小変形へ送る。
注意 \ref{artin-remark-canonical-element} を参照せよ。したがって $x$ が半普遍的ならば、
形式的変形理論の補題 \ref{formal-defos-lemma-versal-criterion} によりこの写像は全射である。
ゆえに $H^0(\xi \otimes_A^{\mathbf{L}} k)$ は単射である。これにより、$u_0$ が
閉点の場合の (1) $\Rightarrow$ (2) の証明が完了する。

\medskip\noindent
証明の残りでは、$H^{-1}(E \otimes_A^\mathbf{L} k) \to
H^{-1}(\NL_{A/\Lambda} \otimes_A^\mathbf{L} k)$ は全射であり、
$H^0(E \otimes_A^\mathbf{L} k) \to
H^0(\NL_{A/\Lambda} \otimes_A^\mathbf{L} k)$ は単射であると仮定する。
$R = A_\mathfrak p^\wedge$ とおき、$\eta$ を $x|_{\Spec(R)}$ に付随する $R$ 上の
形式対象とする。接空間上の写像 $d\underline{\eta}$ は全射である。実際、これは単射
$H^0(E \otimes_A^{\mathbf{L}} k) \to
H^0(\NL_{A/\Lambda} \otimes_A^{\mathbf{L}} k)$ の双対と同一視される
（前段落を参照せよ）。形式的変形理論の補題
\ref{formal-defos-lemma-versal-criterion} によれば、次を証明すれば十分である：
$C' \to C$ を、剰余体が $k$ である有限型 Artin 局所 $\Lambda$-代数の小さい拡大とする。
$R \to C$ を、剰余体の同一視と両立する $\Lambda$-代数準同型とする。
$y = x|_{\Spec(C)}$ とし、$y'$ を $y$ の $C'$ への持ち上げとする。
示すべきことは、$\Lambda$-代数準同型 $R \to C$ を $R \to C'$ へ持ち上げられることである。

\medskip\noindent
$\Lambda$-代数準同型 $A \to C$ を持ち上げれば十分であることに注意する。
$I = \Ker(C' \to C)$ とする。$I$ は $1$-次元 $k$-ベクトル空間である。
$A \to C$ を持ち上げることへの障害 $ob$ は
$\Ext^1_A(\NL_{A/\Lambda}, I)$ の元である。例 \ref{artin-example-key} を参照せよ。
補題 \ref{artin-lemma-compute-ext-into-field} と仮定により、$\xi$ は単射
$$
\Ext^1_A(\NL_{A/\Lambda}, I)
\longrightarrow
\Ext^1_A(E, I)
$$
を誘導する。$ob$ の構成と (i) により、$\Ext^1_A(E, I)$ における $ob$ の像は、
$x$ を $A \times_C C'$ へ持ち上げることへの障害である。(RS*) により、
$y/C$ が $y'/C'$ へ持ち上がるという事実から、$x$ は $A \times_C C'$ へ持ち上がる。
したがって $ob = 0$ であり、証明が完了する。
\end{proof}

\noindent
上の主要補題により、いくつかの場合に半普遍性の開性を得られる。

\begin{lemma}
\label{artin-lemma-openness}
$S$ を局所 Noether スキームとする。$\mathcal{X}$ を $(\Sch/S)_{fppf}$ 上の
グルーポイドをファイバーとする圏で、(RS*) を満たすものとする。
$U = \Spec(A)$ を $S$ 上有限型なアフィンスキームで、あるアフィン開集合
$\Spec(\Lambda)$ へ写るものとする。$x$ を $U$ 上の $\mathcal{X}$ の対象とする。
$\xi : E \to \NL_{A/\Lambda}$ を $D^{-}(A)$ の射とする。次を仮定する：
\begin{enumerate}
\item[(i)] 任意の変形状況 $(x, A' \to A)$ に対して、$x$ が $\Spec(A')$ へ
持ち上がることと、$E \to \NL_{A/\Lambda} \to \NL_{A/A'}$ が零であることは同値である；
\item[(ii)] 関手の同型 $T_x(-) \to \Ext^0_A(E, -)$ が存在し、そのもとで
$E \to \NL_{A/\Lambda} \to \Omega^1_{A/\Lambda}$ は標準元に対応する
（注意 \ref{artin-remark-canonical-element} を参照せよ）；
\item[(iii)] $E$ のコホモロジー群は有限 $A$-加群である。
\end{enumerate}
$x$ が閉点 $u_0 \in U$ において半普遍的ならば、$x$ が $U'$ のすべての有限型点で
半普遍的となるような開近傍 $u_0 \in U' \subset U$ が存在する。
\end{lemma}

\begin{proof}
$C$ を $\xi$ の錐とする。このとき $D^{-}(A)$ における識別三角形
$$
E \to \NL_{A/\Lambda} \to C \to E[1]
$$
がある。補題 \ref{artin-lemma-characterize-versal} により、$x$ が $u_0$ において
半普遍的であるという仮定から $H^{-1}(C \otimes^\mathbf{L} k) = 0$ が従う。
代数続論の補題 \ref{more-algebra-lemma-cut-complex-in-two} により、$u_0$ に対応する
素イデアルに含まれない $f \in A$ で、任意の $A_f$-加群 $M$ に対して
$H^{-1}(C \otimes^\mathbf{L}_A M) = 0$ となるものが存在する。
補題 \ref{artin-lemma-characterize-versal} をもう一度用いると、開集合 $D(f) \subset U$ の
すべての有限型点において半普遍性が成り立つことがわかる。
\end{proof}

\noindent
上の技術的補題は次の定義を示唆する。

\begin{definition}
\label{artin-definition-naive-obstruction-theory}
$S$ を局所 Noether 的基礎とする。$\mathcal{X}$ を $(\Sch/S)_{fppf}$ 上の
グルーポイドをファイバーとする圏とし、$\mathcal{X}$ は (RS*) を満たすと仮定する。
{\it 素朴な障害理論}とは、次のデータからなる：
\begin{enumerate}
\item
\label{artin-item-map}
$\Spec(A) \to S$ がアフィン開集合 $\Spec(\Lambda) \subset S$ へ写るような任意の
$S$-代数 $A$ と、$\Spec(A)$ 上の $\mathcal{X}$ の任意の対象 $x$ に対して、
対象 $E_x \in D^-(A)$ と写像 $\xi_x : E \to \NL_{A/\Lambda}$ が与えられる；
\item
\label{artin-item-inf}
(\ref{artin-item-map}) における $(x, A)$ に対して、関手の変換
$$
\text{Inf}_x( - ) \to \Ext^{-1}_A(E_x, -)
\quad\text{and}\quad
T_x(-) \to \Ext^0_A(E_x, -)
$$
がある；
\item
\label{artin-item-functoriality}
(\ref{artin-item-map}) における $(x, A)$ と環準同型 $A \to B$ に対して、
$y = x|_{\Spec(B)}$ とおくと、$D(A)$ における関手性写像 $E_x \to E_y$ がある。
\end{enumerate}
これらのデータには、次の条件を課す：
\begin{enumerate}
\item[(i)]
(\ref{artin-item-functoriality}) の状況において、図式
$$
\xymatrix{
E_y \ar[r]_{\xi_y} & \NL_{B/\Lambda} \\
E_x \ar[u] \ar[r]^{\xi_x} & \NL_{A/\Lambda} \ar[u]
}
$$
は $D(A)$ において可換である；
\item[(ii)]
(\ref{artin-item-map}) における $(x, A)$ と $A \to B \to C$ が与えられたとき、
$y = x|_{\Spec(B)}$ および $z = x|_{\Spec(C)}$ とおけば、関手性写像
$E_x \to E_y$ と $E_y \to E_z$ の合成は関手性写像 $E_x \to E_z$ である；
\item[(iii)]
(\ref{artin-item-inf}) の写像は同型であり、関手性写像および注意
\ref{artin-remark-functoriality} の写像と両立する；
\item[(iv)]
合成 $E_x \to \NL_{A/\Lambda} \to \Omega_{A/\Lambda}$ は
$T_x(\Omega_{A/\Lambda}) = \Ext^0(E_x, \Omega_{A/\Lambda})$ の標準元に対応する。
注意 \ref{artin-remark-canonical-element} を参照せよ；
\item[(v)]
$I = \Ker(A' \to A)$ である変形状況 $(x, A' \to A)$ が与えられたとき、
合成 $E_x \to \NL_{A/\Lambda} \to \NL_{A/A'}$ が
$$
\Hom_A(E_x, \NL_{A/\Lambda}) = \Ext^0_A(E_x, \NL_{A/A'}) =
\Ext^1_A(E_x, I)
$$
において零であることと、$x$ が $A'$ へ持ち上がることは同値である。
\end{enumerate}
\end{definition}

\noindent
したがって特に、$\mathcal{O}_x( - ) = \Ext^1_A(E_x, -)$ とおくことにより、
節 \ref{artin-section-obstruction-theory} におけるような障害理論を得る。

\begin{lemma}
\label{artin-lemma-naive-obstruction-theory-qis}
$S$ と $\mathcal{X}$ を定義 \ref{artin-definition-naive-obstruction-theory} におけるものとし、
$\mathcal{X}$ には素朴な障害理論が与えられているとする。
$A \to B$ と $y \to x$ を (\ref{artin-item-functoriality}) におけるものとする。
$k$ を、体であるような $B$-代数とする。このとき関手性写像 $E_x \to E_y$ は、
$i = 0, 1$ に対して全単射
$$
H^i(E_x \otimes_A^{\mathbf{L}} k) \to H^i(E_y \otimes_B^{\mathbf{L}} k)
$$
を誘導する。
\end{lemma}

\begin{proof}
$z = x|_{\Spec(k)}$ とおく。(RS*) と $A[k] = A \times_k k[k]$ および
$B[k] = B \times_k k[k]$ により
$$
\textit{Lift}(x, A[k]) = \textit{Lift}(z, k[k])
\quad\text{and}\quad
\textit{Lift}(y, B[k]) = \textit{Lift}(z, k[k])
$$
である。したがって素朴な障害理論の性質により、関手性写像 $E_x \to E_y$ は
$i = -1, 0$ に対して全単射
$\Ext^i_A(E_x, k) \to \text{Ext}^i_B(E_y, k)$ を誘導する。
補題 \ref{artin-lemma-compute-ext-into-field} により、写像
$H^i(E_x \otimes_A^{\mathbf{L}} k) \to H^i(E_y \otimes_B^{\mathbf{L}} k)$、
$i = 0, 1$ は双対ベクトル空間上で同型を誘導し、したがってそれ自身も同型である。
\end{proof}

\begin{lemma}
\label{artin-lemma-naive-obstruction-theory-gives-openness}
$S$ を局所 Noether スキームとする。
$p : \mathcal{X} \to (\Sch/S)_{fppf}^{opp}$ をグルーポイドをファイバーとする圏とする。
$\mathcal{X}$ は (RS*) を満たし、$\mathcal{X}$ は素朴な障害理論をもつと仮定する。
定義 \ref{artin-definition-naive-obstruction-theory} の複体 $E_x$ が、$A$ が $S$ 上有限型で
あるような組 $(A, x)$ に対して有限生成コホモロジー群をもつならば、
$\mathcal{X}$ に対して半普遍性の開性が成り立つ。
\end{lemma}

\begin{proof}
$U$ を $S$ 上局所有限型なスキームとし、$x$ を $U$ 上の $\mathcal{X}$ の対象とする。
$u_0$ を $U$ の有限型点で、$x$ が $u_0$ において半普遍的となるものとする。
まず $U$ を、$u_0$ が閉点で、かつ $U \to S$ がアフィン開集合
$\Spec(\Lambda)$ へ写るようなアフィンスキームへ縮小してよい。
$U = \Spec(A)$ と書く。$\xi_x : E_x \to \NL_{A/\Lambda}$ を障害写像とする。
ここで補題 \ref{artin-lemma-openness} を適用すれば結論を得る。
\end{proof}

\section{双対的な概念}
\label{artin-section-dual}

\noindent
$(x, A' \to A)$ を、局所 Noether スキーム $S$ 上の与えられたグルーポイドを
ファイバーとする圏 $\mathcal{X}$ に対する変形状況とする。$\mathcal{X}$ は障害理論を
もつと仮定する。定義 \ref{artin-definition-obstruction-theory} を参照せよ。
実際には、しばしば $A$-加群の複体 $K^\bullet$ と関手の同型
$$
\text{Inf}_x(-) \to H^0(K^\bullet \otimes_A^\mathbf{L} -),\quad
T_x(-) \to H^1(K^\bullet \otimes_A^\mathbf{L} -),\quad
\mathcal{O}_x(-) \to H^2(K^\bullet \otimes_A^\mathbf{L} -)
$$
がある。本節ではこれを少し形式化し、いくつかの場合にこれが半普遍性の開性の確認を
どのように導くかを示す。

\begin{example}
\label{artin-example-global-sections-dual}
$\Lambda, S, W, \mathcal{F}$ を例 \ref{artin-example-global-sections} におけるものとする。
$W \to S$ は固有で、$\mathcal{F}$ は連接であると仮定する。スキームのコホモロジーの
注意 \ref{coherent-remark-explain-perfect-direct-image} により、
$\mathcal{F}$ のコホモロジーを普遍的に計算する有限射影 $\Lambda$-加群の有限複体
$N^\bullet$ が存在する。特に、例 \ref{artin-example-global-sections} の障害空間は
$\mathcal{O}_x(M) = H^1(N^\bullet \otimes_\Lambda M)$ である。
したがって $K^\bullet = N^\bullet \otimes_\Lambda A[-1]$ とおけば、
$\mathcal{O}_x(M) = H^2(K^\bullet \otimes_A^\mathbf{L} M)$ である。
\end{example}

\begin{situation}
\label{artin-situation-dual}
$S$ を局所 Noether スキームとする。$\mathcal{X}$ を $(\Sch/S)_{fppf}$ 上の
グルーポイドをファイバーとする圏とする。関手 $T_x(-)$ について語れるように、
$\mathcal{X}$ は (RS*) をもつと仮定する。補題
\ref{artin-lemma-properties-lift-RS-star} を参照せよ。
$U = \Spec(A)$ を $S$ 上有限型なアフィンスキームで、アフィン開集合
$\Spec(\Lambda)$ へ写るものとする。$x$ を $U$ 上の $\mathcal{X}$ の対象とする。
次が与えられていると仮定する：
\begin{enumerate}
\item $A$-加群の複体 $K^\bullet$；
\item 関手の変換 $T_x(-) \to H^1(K^\bullet \otimes_A^\mathbf{L} -)$；
\item 核が $I = \Ker(A' \to A)$ である任意の変形状況 $(x, A' \to A)$ に対する元
$o_x(A') \in H^2(K^\bullet \otimes_A^\mathbf{L} I)$。
\end{enumerate}
これらは次の（最小限の）条件を満たすものとする：
\begin{enumerate}
\item[(i)] 変換 $T_x(-) \to H^1(K^\bullet \otimes_A^\mathbf{L} -)$ は同型である；
\item[(ii)] 変形状況の射 $(x, A'' \to A) \to (x, A' \to A)$ が与えられたとき、
元 $o_x(A')$ は写像
$H^2(K^\bullet \otimes_A^\mathbf{L} I) \to
H^2(K^\bullet \otimes_A^\mathbf{L} I')$ によって元 $o_x(A'')$ へ写る。
ここで $I' = \Ker(A'' \to A)$ とする；
\item[(iii)] $x$ が $\Spec(A')$ 上の対象へ持ち上がることと
$o_x(A') = 0$ は同値である。
\end{enumerate}
無限小自己同型も取り入れることは可能だが、可能な限り鋭い結果を得るため、
ここではそうしない。
\end{situation}

\noindent
状況 \ref{artin-situation-dual} では、$K^\bullet \otimes_A^\mathbf{L} \NL_{A/\Lambda}$ が
重要な役割を果たす。元
$\xi \in H^1(K^\bullet \otimes_A^\mathbf{L} \NL_{A/\Lambda})$ が与えられたとする。
このとき、(1) 平方零な核 $I$ をもつ任意の $\Lambda$-代数の全射 $A' \to A$ に対して、
標準写像 $\NL_{A/\Lambda} \to \NL_{A/A'} = I[1]$ は $\xi$ を元
$\xi_{A'} \in H^2(K^\bullet \otimes_A^\mathbf{L} I)$ へ送り、(2) 写像
$\NL_{A/\Lambda} \to \Omega_{A/\Lambda}$ は $\xi$ を
$H^1(K^\bullet \otimes_A^\mathbf{L} \Omega_{A/\Lambda})$ の元 $\xi_{can}$ へ送る。

\begin{lemma}
\label{artin-lemma-dual-obstruction}
状況 \ref{artin-situation-dual} において、さらに次を仮定する：
\begin{enumerate}
\item[(iv)] 注意 \ref{artin-remark-short-exact-sequence-thickenings} におけるような
変形状況の短完全列と持ち上げ $x'_2 \in \text{Lift}(x, A_2')$ が与えられたとき、
$o_x(A_3') \in H^2(K^\bullet \otimes_A^\mathbf{L} I_3)$ は
$\partial\theta$ に等しい。ここで
$\theta \in H^1(K^\bullet \otimes_A^\mathbf{L} I_1)$ は、
$A_1' = A[I_1]$ と与えられた写像
$T_x(-) \to H^1(K^\bullet \otimes_A^\mathbf{L} -)$ を介して
$x'_2|_{\Spec(A_1')}$ に対応する元である。
\end{enumerate}
この場合、次を満たす元
$\xi \in H^1(K^\bullet \otimes_A^\mathbf{L} \NL_{A/\Lambda})$ が存在する：
\begin{enumerate}
\item 任意の変形状況 $(x, A' \to A)$ に対して $\xi_{A'} = o_x(A')$ である；
\item $\xi_{can}$ は、与えられた変換
$T_x(-) \to H^1(K^\bullet \otimes_A^\mathbf{L} -)$ を介して、注意
\ref{artin-remark-canonical-element} の標準元と一致する。
\end{enumerate}
\end{lemma}

\begin{proof}
核が $J$ である $\alpha : \Lambda[x_1, \ldots, x_n] \to A$ を選ぶ。
$P = \Lambda[x_1, \ldots, x_n]$ と書く。証明の残りでは
$$
\NL(\alpha) = (J/J^2 \longrightarrow \bigoplus A \text{d}x_i)
$$
を用いる。これは代数の補題 \ref{algebra-lemma-NL-homotopy} および代数続論の補題
\ref{more-algebra-lemma-derived-tor-homotopy} により許される。
元 $o_x(P/J^2) \in H^2(K^\bullet \otimes_A^\mathbf{L} J/J^2)$ と商
$$
C = (P/J^2 \times \bigoplus A \text{d}x_i)/(J/J^2)
$$
を考える。ここで $J/J^2$ は対角的に埋め込まれる。$C \to A$ は核が
$\bigoplus A\text{d}x_i$ である全射であることに注意する。さらに、押し出しにおいて
$f \in P$ の類を $(f, \text{d}f)$ の類へ送ることにより与えられる、$C \to A$ の
切断 $A \to C$ がある。後で用いるため、対応する射
$\Spec(C) \to \Spec(A)$ に沿った $x$ の引き戻しを $x_C$ と書く。
したがって $o_x(C) = 0$ である。ゆえに $o_x(P/J^2)$ は
$H^2(K^\bullet \otimes_A^\mathbf{L} \bigoplus A\text{d}x_i)$ において零へ写る。
したがって $o_x(P/J^2)$ へ写る元
$\xi \in H^1(K^\bullet \otimes_A^\mathbf{L} \NL(\alpha))$ が存在する。

\medskip\noindent
任意の変形状況 $(x, A' \to A)$ に対して、$A$ への増大写像と両立する
$\Lambda$-代数準同型 $P/J^2 \to A'$ が存在することに注意する。
したがって構成と状況 \ref{artin-situation-dual} の性質 (ii) により、元 $\xi$ は
補題の第一の性質を満たす。

\medskip\noindent
$\xi$ の選択は、$H^1(K^\bullet \otimes_A^\mathbf{L} \bigoplus A\text{d}x_i)$ の
元を選ぶ不定性を除いて定まっていたことに注意する。前述の群の元によって $\xi$ を
変更した後、補題の第二の主張が成り立つようにできることを示す。
$C' \subset C$ を、$\Lambda$-代数準同型 $P/J^2 \to C$（第一因子の包含）による
$P/J^2$ の像とする。
$\Ker(C' \to A) = \Im(J/J^2 \to \bigoplus A\text{d}x_i)$ であることに注意する。
$\overline{C} = A[\Omega_{A/\Lambda}]$ とおく。写像
$P/J^2 \times \bigoplus A \text{d}x_i \to \overline{C}$、
$(f, \sum f_i \text{d}x_i) \mapsto (f \bmod J, \sum f_i \text{d}x_i)$ は
全射 $C \to \overline{C}$ を経由する。このとき
$$
(x, \overline{C} \to A) \to (x, C \to A) \to (x, C' \to A)
$$
は変形状況の短完全列である。付随する分裂
$\overline{C} = A[\Omega_{A/\Lambda}]$（注意
\ref{artin-remark-short-exact-sequence-thickenings} による）は、上で与えた分裂に等しい。
さらに、切断 $A \to C$ と写像 $C \to \overline{C}$ の合成は、注意
\ref{artin-remark-canonical-element} の写像
$(1, \text{d}) : A \to A[\Omega_{A/\Lambda}]$ である。
したがって $x_C$ は
$T_x(\Omega_{A/\Lambda}) = \text{Lift}(x, A[\Omega_{A/\Lambda}])$ の標準元
$x_{can}$ へ制限される。条件 (iv) により、$o_x(P/J^2)$ は
$$
H^1(K^\bullet \otimes_A^\mathbf{L} \Im(J/J^2 \to \bigoplus A\text{d}x_i))
$$
において $\partial x_{can}$ へ写る。構成により $\xi$ は $o_x(P/J^2)$ へ写る。
したがって $x_{can}$ と $\xi_{can}$ は上の群の同じ元へ写る。これは
（コホモロジー長完全列により）両者が
$H^1(K^\bullet \otimes_A^\mathbf{L} \bigoplus A\text{d}x_i)$ の元だけ異なることを意味し、
望む結論を得る。
\end{proof}

\begin{lemma}
\label{artin-lemma-dual-openness}
状況 \ref{artin-situation-dual} において、補題 \ref{artin-lemma-dual-obstruction} の (iv) が
成り立ち、$K^\bullet$ は $D(A)$ の完全対象であると仮定する。この場合、
$x$ が閉点 $u_0 \in U$ において半普遍的ならば、$x$ が $U'$ のすべての有限型点で
半普遍的となるような開近傍 $u_0 \in U' \subset U$ が存在する。
\end{lemma}

\begin{proof}
$K^\bullet$ は有限射影 $A$-加群の有限複体であると仮定してよい。
したがって $K^\bullet$ との導来テンソル積は、単に $K^\bullet$ とテンソルをとることに
等しい。$E^\bullet$ を $K^\bullet$ の双対完全複体とする。代数続論の補題
\ref{more-algebra-lemma-dual-perfect-complex} を参照せよ。
（したがって $E^n = \Hom_A(K^{-n}, A)$ であり、その微分は $K^\bullet$ の微分の
転置である。）$E \in D^{-}(A)$ を複体 $E^\bullet[-1]$ により表現される対象とする。
$\xi \in H^1(\text{Tot}(K^\bullet \otimes_A \NL_{A/\Lambda}))$ を補題
\ref{artin-lemma-dual-obstruction} で構成された元とし、対応する写像を
$\xi : E = E^\bullet[-1] \to \NL_{A/\Lambda}$ と書く（前掲箇所）。
組 $(E, \xi)$ は補題 \ref{artin-lemma-openness} のすべての仮定を満たすと主張する。
これにより証明が完了する。

\medskip\noindent
実際、補題 \ref{artin-lemma-openness} の仮定 (i) は、補題
\ref{artin-lemma-dual-obstruction} の結論 (1) と、前掲箇所による
$H^2(K^\bullet \otimes_A^\mathbf{L} -) = \Ext^1(E, -)$ から従う。
補題 \ref{artin-lemma-openness} の仮定 (ii) は、補題
\ref{artin-lemma-dual-obstruction} の結論 (2) と、前掲箇所による
$H^1(K^\bullet \otimes_A^\mathbf{L} -) = \Ext^0(E, -)$ から従う。
補題 \ref{artin-lemma-openness} の仮定 (iii) は明らかである。
\end{proof}

\section{Noether スキーム上の極限を保つ関手}
\label{artin-section-noetherian}

\noindent
（局所）Noether スキームの充満部分圏上だけで定義された関手やスタックを考えると
便利なことがある。本節では、代数空間の場合にこれを論じる。

\medskip\noindent
$S$ を局所 Noether スキームとする。圏を正確に揃えるため、いくらか細かく述べよう。
集合論的問題を無視する場合には、以下で選ぶスキームの集合を全スキームの集まりで
置き換えてよい。位相の注意 \ref{topologies-remark-choice-sites} と同様に、$S$ を含む
スキームの圏 $\Sch_\alpha$ を、同じ台圏 $\Sch_\alpha/S$ をもつ大サイト
$(\Sch/S)_{Zar}$、$(\Sch/S)_\etale$、$(\Sch/S)_{smooth}$、
$(\Sch/S)_{syntomic}$、$(\Sch/S)_{fppf}$ が得られるように選ぶ。
局所 Noether スキームからなる充満部分圏を
$$
\textit{Noetherian}_\alpha \subset \Sch_\alpha
$$
と書く。これにより充満部分圏
$$
\textit{Noetherian}_\alpha/S \subset \Sch_\alpha/S
$$
が定まる。$\tau \in \{Zariski, \etale, smooth, syntomic, fppf\}$ に対して、次が成り立つ：
\begin{enumerate}
\item $f : X \to Y$ が $\Sch_\alpha/S$ の射で、
$Y$ が $\textit{Noetherian}_\alpha/S$ に属し、$f$ が局所有限型ならば、
$X$ は $\textit{Noetherian}_\alpha/S$ に属する；
\item $f : X \to Y$ と $g : Z \to Y$ が $\textit{Noetherian}_\alpha/S$ の射で、
$f$ が局所有限型ならば、$\textit{Noetherian}_\alpha/S$ におけるファイバー積
$X \times_Y Z$ が存在し、$\Sch_\alpha/S$ におけるファイバー積と一致する；
\item $\{X_i \to X\}_{i \in I}$ が $(\Sch/S)_\tau$ の被覆で、
$X$ が $\textit{Noetherian}_\alpha/S$ に属するならば、各 $X_i$ も
$\textit{Noetherian}_\alpha/S$ に属する；
\item $\textit{Noetherian}_\alpha/S$ に、対象が
$\textit{Noetherian}_\alpha/S$ に属する $(\Sch/S)_\tau$ の被覆の集合を与えた圏は
サイトであり、これを $(\textit{Noetherian}/S)_\tau$ と書く；
\item 包含関手 $(\textit{Noetherian}/S)_\tau \to (\Sch/S)_\tau$ は
充満忠実、連続、かつ余連続である。
\end{enumerate}
サイトの補題 \ref{sites-lemma-cocontinuous-morphism-topoi} および
\ref{sites-lemma-when-shriek} により、トポスの射
$$
g_\tau : \Sh((\textit{Noetherian}/S)_\tau) \longrightarrow \Sh((\Sch/S)_\tau)
$$
を得る。その引き戻し関手は、包含関手
$(\textit{Noetherian}/S)_\tau \to (\Sch/S)_\tau$ に沿った層の制限である。

\begin{remark}[警告]
\label{artin-remark-no-fibre-products}
サイト $(\textit{Noetherian}/S)_\tau$ はファイバー積をもたない。
したがって層を扱う際には注意が必要である。たとえば連続な包含関手
$(\textit{Noetherian}/S)_\tau \to (\Sch/S)_\tau$ はサイトの射を定めない。
$\tau = fppf$ の場合の例については、例の節
\ref{examples-section-sheaves-locally-Noetherian} を参照せよ。
\end{remark}

\noindent
$F : (\textit{Noetherian}/S)_\tau^{opp} \to \textit{Sets}$ を関手とする。
$(\textit{Noetherian}/S)_\tau$ のアフィンスキームの任意の有向極限
$X = \lim X_i$ に対して $F(X) = \colim F(X_i)$ となるとき、$F$ は
{\it 極限を保つ}という。

\begin{lemma}
\label{artin-lemma-canonical-extension}
$\tau \in \{Zariski, \etale, smooth, syntomic, fppf\}$ とする。包含関手
$(\textit{Noetherian}/S)_\tau \to (\Sch/S)_\tau$ に沿った制限は、次の圏の間の圏同値を
定める：
\begin{enumerate}
\item $(\Sch/S)_\tau$ 上の極限を保つ層の圏；
\item $(\textit{Noetherian}/S)_\tau$ 上の極限を保つ層の圏。
\end{enumerate}
\end{lemma}

\begin{proof}
$F : (\textit{Noetherian}/S)_\tau^{opp} \to \textit{Sets}$ を、極限を保ち、かつ層である
関手とする。位相の補題 \ref{topologies-lemma-extend} および
\ref{topologies-lemma-extend-sheaf-general} により、極限を保ち、層であり、$F$ に制限される
一意的な関手 $F' : (\Sch/S)_\tau^{opp} \to \textit{Sets}$ が存在する。
実際、位相の補題 \ref{topologies-lemma-extend} における $F'$ の構成は $F$ に関して
関手的であり、この構成は制限の擬逆である。いくつかの詳細は省略する。
\end{proof}

\begin{lemma}
\label{artin-lemma-representable-limit-preserving}
$X$ を $(\textit{Noetherian}/S)_\tau$ の対象とする。
\par\noindent
点の関手
$h_X : (\textit{Noetherian}/S)_\tau^{opp} \to \textit{Sets}$ が極限を保つならば、
$X$ は $S$ 上局所有限表示である。
\end{lemma}

\begin{proof}
$V \subset X$ をアフィン開部分スキームで、あるアフィン開集合 $U \subset S$ へ
写るものとする。代数の補題 \ref{algebra-lemma-ring-colimit-fp} により、
$U$ 上有限表示なアフィンスキーム $V_i$ の有向極限として $V = \lim V_i$ と書ける。
仮定により、射 $V \to X$ はある $i$ に対して $V \to V_i \to X$ と分解する。
$i$ を大きくすることにより、$V_i$ における $V \subset X$ の逆像は最終的に $V_i$ に
等しくなるので、$V_i \to X$ は $V$ を経由すると仮定してよい。極限の補題
\ref{limits-lemma-descend-opens} を参照せよ。このとき $U$ 上のスキームの圏において、
恒等射 $V \to V$ はある $i$ に対して $V_i$ を経由する。したがって $V \to U$ は
有限表示である。対応する代数的事実は、$B$ が $A$-代数で、
$\text{id} : B \to B$ が有限表示 $A$-代数を経由するならば、$B$ は $A$ 上有限表示で
あるというものである（よい演習である）。ゆえに $X$ は $S$ 上局所有限表示である。
\end{proof}

\noindent
次の補題には、代数空間で表現可能な変換に対する変種がある。

\begin{lemma}
\label{artin-lemma-representable}
$\tau \in \{Zariski, \etale, smooth, syntomic, fppf\}$ とする。
$F', G' : (\Sch/S)_\tau^{opp} \to \textit{Sets}$ を、極限を保つ層とする。
$a' : F' \to G'$ を関手の変換とする。$a : F \to G$ を、
$a' : F' \to G'$ の $(\textit{Noetherian}/S)_\tau$ への制限と書く。
次は同値である：
\begin{enumerate}
\item $a'$ は表現可能である（関手の変換としての意味である。圏の定義
\ref{categories-definition-representable-morphism} を参照せよ）；
\item $(\textit{Noetherian}/S)_\tau$ の任意の対象 $V$ と任意の写像 $V \to G$ に対して、
ファイバー積 $F \times_G V : (\textit{Noetherian}/S)_\tau^{opp} \to \textit{Sets}$ は
表現可能な関手である；
\item (2) と同じであるが、$S$ 上アフィン有限型で $S$ のアフィン開集合へ写る
$V$ だけを考える。
\end{enumerate}
\end{lemma}

\begin{proof}
(1) を仮定する。空間の極限の補題
\ref{spaces-limits-lemma-locally-finite-presentation-permanence} により、変換 $a'$ は
極限を保つ\footnote{$\tau \not = fppf$ であってもこれは意味をもつ。実際、
$(\Sch/S)_\tau$ の台圏は $(\Sch/S)_{fppf}$ の台圏と等しく、この主張は位相に
言及しない。}。(2) におけるような $\xi : V \to G$ をとる。
$V' = V$ と書くが、これは $(\Sch/S)_\tau$ の対象とみなしたものである。
$G$ は $G'$ の $(\textit{Noetherian}/S)_\tau$ への制限なので、
$\xi \in G(V)$ は $\xi' \in G'(V')$ に対応する。仮定により
$V' \times_{\xi', G'} F'$ はスキーム $U'$ で表現される。射影
$V' \times_{\xi', G'} F' \to V'$ に対応するスキームの射 $U' \to V'$ は、
空間の極限の補題
\ref{spaces-limits-lemma-base-change-locally-finite-presentation} および極限の命題
\ref{limits-proposition-characterize-locally-finite-presentation} により局所有限表示である。
したがって $U'$ は局所 Noether スキームであり、それゆえ $U'$ は
$(\textit{Noetherian}/S)_\tau$ のある対象 $U$ と同型である。このとき望みどおり、
$U$ は $F \times_G V$ を表現する。

\medskip\noindent
(2) $\Rightarrow$ (3) は直ちにわかる。(3) を仮定し、(1) を証明する。
$T$ を $(\Sch/S)_\tau$ の対象とし、射 $T \to G'$ が与えられているとする。
関手 $F' \times_{G'} T$ が $T$ 上のスキーム $X$ で表現されることを示さなければ
ならない。$\mathcal{B}$ を、$S$ のアフィン開集合へ写る $T$ のアフィン開集合の
集合とする。これは $T$ の位相の基底である。以下で、$W \in \mathcal{B}$ に対して
ファイバー積 $F' \times_{G'} W$ が $W$ 上のスキーム $X_W$ で表現されることを示す。
$\mathcal{B}$ において $W_1 \subset W_2$ ならば、$X_{W_1}$ と
$X_{W_2} \times_{W_2} W_1$ はともに関手 $F' \times_{G'} W_1$ を表現するので、
同型 $X_{W_1} \to X_{W_2} \times_{W_2} W_1$ を得る。これらの同型は標準的であり、
構成の補題 \ref{constructions-lemma-relative-glueing} に述べられたコサイクル条件を満たす。
したがってスキーム $X_W$ を貼り合わせて $T$ 上のスキーム $X$ を得る。
貼り合わせ写像と写像 $X_W \to F'$ の両立性により、写像 $X \to F'$ を得る。
その結果得られる写像 $X \to F' \times_{G'} T$ は同型である。実際、これは $T$ 上で
局所的に確認できる（この射の始域と終域は Zariski 位相について層だからである）。

\medskip\noindent
$W$ をアフィンスキームで、あるアフィン開集合 $U \subset S$ へ写るものとする。
写像 $W \to G'$ が与えられているとする。引き続き (3) を仮定して、
$F' \times_{G'} W$ がスキームで表現されることを示さなければならない。
代数の補題 \ref{algebra-lemma-ring-colimit-fp} により、$U$ 上有限表示な
アフィンスキーム $V'_i$ の有向極限として $W = \lim V'_i$ と書ける。
$V'_i$ は Noether スキーム上有限型なので、$V'_i$ は Noether スキームである。
$V_i = V'_i$ と書くが、これは $(\textit{Noetherian}/S)_\tau$ の対象とみなしたものである。
$G'$ は極限を保つので、$W \to G'$ が合成 $W \to V'_i \to G'$ となるような
$i$ と写像 $V'_i \to G'$ を選べる。$G$ は $G'$ の
$(\textit{Noetherian}/S)_\tau$ への制限なので、射 $V'_i \to G'$ は射
$V_i \to G$ と同じものである（上を参照せよ）。仮定 (3) により、関手
$F \times_G V_i$ は $(\textit{Noetherian}/S)_\tau$ の対象 $X_i$ で表現される。
関手 $F \times_G V_i$ は極限を保つ。実際、これは $F' \times_{G'} V'_i$ の制限であり、
後者は、空間の極限の補題
\ref{spaces-limits-lemma-fibre-product-locally-finite-presentation}、$F'$ と $G'$ が
極限を保つという仮定、および $V'_i$ の点の関手が極限を保つことを述べる極限の注意
\ref{limits-remark-limit-preserving} により、極限を保つ。
補題 \ref{artin-lemma-representable-limit-preserving} により、$X_i$ は $S$ 上局所有限表示である。
$X'_i = X_i$ と書くが、これは $(\Sch/S)_\tau$ の対象とみなしたものである。
すると $F' \times_{G'} V'_i$ と点の関手 $h_{X'_i}$ はともに
$h_{X_i} : (\textit{Noetherian}/S)_\tau^{opp} \to \textit{Sets}$ を
$(\Sch/S)_\tau$ 上の極限を保つ層へ拡張したものである。補題
\ref{artin-lemma-canonical-extension} の圏同値により、$X'_i$ は
$F' \times_{G'} V'_i$ を表現すると結論する。最後に
$$
F' \times_{G'} W = F' \times_{G'} V'_i \times_{V'_i} W =
X'_i \times_{V'_i} W
$$
は望みどおり表現可能である。
\end{proof}

\section{Noether 的状況における代数空間}
\label{artin-section-algebraic-spaces-noetherian}

\noindent
$S$ を局所 Noether スキームとする。節 \ref{artin-section-noetherian} で調べたサイトを次で表す：
\par
\centerline{$(\textit{Noetherian}/S)_\etale \subset (\Sch/S)_\etale$}
\par\noindent
$F : (\textit{Noetherian}/S)_\etale^{opp} \to \textit{Sets}$ を関手とする。
すなわち $F$ は $(\textit{Noetherian}/S)_\etale$ 上の前層である。この状況では、節
\ref{artin-section-axioms-functors} の公理 [-1]、[0]、[1]、[2]、[3]、[4]、[5] はすべて
意味をもつ。各々を順に見直し、ここで何を意味するかを正確に述べる。

\medskip\noindent
公理 [-1]。これは集合論的問題に関心のない読者が無視してよい集合論的条件である。
これは $S$ 上有限型な体のスペクトルにおける $F$ の値に関する条件であり、そのような
スペクトルは $(\textit{Noetherian}/S)_\etale$ の対象なので、$F$ に対して意味をもつ。

\medskip\noindent
公理 [0]。これは $F$ が $(\textit{Noetherian}/S)_\etale^{opp}$ 上の層である、
すなわちエタール被覆に対する層条件を満たすという公理である。

\medskip\noindent
公理 [1]。これは、節 \ref{artin-section-noetherian} で定義した意味で $F$ が極限を保つという
公理である。すなわち、$(\textit{Noetherian}/S)_\etale$ のアフィンスキームの任意の
有向極限 $X = \lim X_i$ に対して $F(X) = \colim F(X_i)$ が成り立つ。

\medskip\noindent
公理 [2]。これは $F$ が Rim--Schlessinger 条件 (RS) を満たすという公理である。
定義 \ref{artin-definition-RS} の条件 (RS) の定義と節 \ref{artin-section-axioms-functors} の議論を
見ると、これは次を意味する：$Y, X, X'$ が Artin 局所環のスペクトルであるような、
$S$ 上有限型なスキームの任意の押し出し $Y' = Y \amalg_X X'$ に対して、写像
$$
F(Y \amalg_X X') \to F(Y) \times_{F(X)} F(X')
$$
は全単射である。この条件は意味をもつ。実際、スキーム $X$、$X'$、$Y$、$Y'$ は
$S$ 上有限型なので $(\text{Noetherian}/S)_\etale$ に属する。

\medskip\noindent
公理 [3]。これはすべての接空間 $TF_{k, x_0}$ が有限次元であるという公理である。
接空間 $TF_{k, x_0}$ は、$S$ 上有限型な体 $k$ に対する $\Spec(k)$ と
$\Spec(k[\epsilon])$ における $F$ の値から構成され、したがって圏
$(\textit{Noetherian}/S)_\etale$ の対象における値から得られるので、これは意味をもつ。

\medskip\noindent
公理 [4]。これはすべての形式対象が有効であるという公理である。節
\ref{artin-section-formal-objects} および \ref{artin-section-axioms-functors} の議論を見ると、
これは Noether スキームにおける関手の値だけを用いるので、この条件は $F$ に対して
意味をもつ。

\medskip\noindent
公理 [5]。これは $F$ が半普遍性の開性を満たすという公理である。これは次を
意味することを思い出そう：$U$ を $S$ 上局所有限型なスキームとし、
$x \in F(U)$ とし、$u_0 \in U$ を $x$ が $u_0$ において半普遍的となる有限型点とする。
このとき、$x$ が $U'$ のすべての有限型点において半普遍的となるような開近傍
$u_0 \in U' \subset U$ が存在する。前と同様に、これを確認するには Noether スキームに
おける関手の値だけが用いられる。

\begin{proposition}
\label{artin-proposition-spaces-diagonal-representable-noetherian}
$S$ を局所 Noether スキームとする。
$F : (\textit{Noetherian}/S)_\etale^{opp} \to \textit{Sets}$ を関手とする。
次を仮定する：
\begin{enumerate}
\item $\Delta : F \to F \times F$ は表現可能である（関手の変換としての意味である。
圏の定義 \ref{categories-definition-representable-morphism} を参照せよ）；
\item $F$ は公理 [-1]、[0]、[1]、[2]、[3]、[4]、[5] を満たす（上を参照せよ）；
\item $S$ のすべての有限型点 $s$ に対して $\mathcal{O}_{S, s}$ は G-環である。
\end{enumerate}
このとき、$(\textit{Noetherian}/S)_\etale$ への制限が $F$ となる一意的な代数空間
$F' : (\Sch/S)_{fppf}^{opp} \to \textit{Sets}$ が存在する
（説明については証明を参照せよ）。
\end{proposition}

\begin{proof}
サイト $(\Sch/S)_{fppf}$ と $(\Sch/S)_\etale$ は同じ台圏をもつことを思い出そう。
節 \ref{artin-section-noetherian} の議論を参照せよ。同様に、サイト
$(\textit{Noetherian}/S)_\etale$ と $(\textit{Noetherian}/S)_{fppf}$ も同じ台圏をもつ。
公理 [0] と [1] により、関手 $F$ は層であり、極限を保つ。
$F' : (\Sch/S)_\etale^{opp} \to \textit{Sets}$ を、エタール位相について層であり、
極限を保つような $F$ の一意的な拡張とする。補題 \ref{artin-lemma-canonical-extension} を
参照せよ。このとき $F'$ は、節 \ref{artin-section-axioms-functors} で与えられた
公理 [0] と [1] を満たす。補題 \ref{artin-lemma-representable} により、
$\Delta' : F' \to F' \times F'$ は（スキームで）表現可能である。
一方、$F'$ が節 \ref{artin-section-axioms-functors} の公理 [-1]、[2]、[3]、[4]、[5] を
満たすことは直ちに明らかである。実際、これらはいずれも
$(\textit{Noetherian}/S)_\etale$ の対象における $F'$ の値だけを用い、対応する条件を
$F$ に対して仮定している。したがって命題
\ref{artin-proposition-spaces-diagonal-representable} により、$F'$ は代数空間である。
\end{proof}

\section{収縮に関する Artin の定理}
\label{artin-section-contractions}

\noindent
本節では形式代数空間の言葉を自由に用いる。形式空間の節
\ref{formal-spaces-section-introduction} を参照せよ。収縮に関する Artin の定理は、
Artin の論文 \cite{ArtinII} の二つの主定理の一つである。もう一つは膨張に関する
彼の定理であり、これは形式空間の代数化の節 \ref{restricted-section-dilatations} で
述べ、証明した。

\begin{situation}
\label{artin-situation-contractions}
$S$ を局所 Noether スキームとする。$X'$ を $S$ 上局所有限型な代数空間とする。
$T' \subset |X'|$ を閉部分集合とする。$U' \subset X'$ を
$|U'| = |X'| \setminus T'$ である開部分空間とする。$W$ を $S$ 上の局所 Noether
形式代数空間で、$W_{red}$ が $S$ 上局所有限型であるものとする。最後に
$$
g : X'_{/T'} \longrightarrow W
$$
を形式的修正とする。形式空間の代数化の定義
\ref{restricted-definition-formal-modification} を参照せよ。
$X'_{/T'}$ は $T'$ に沿った $X'$ の形式完備化を表すことを思い出そう。
形式空間の節 \ref{formal-spaces-section-completion} を参照せよ。
\end{situation}

\noindent
上の状況における目標は、$S$ 上の代数空間の固有射 $f : X' \to X$、
閉部分集合 $T \subset |X|$、および形式代数空間の同型 $a : X_{/T} \to W$ で、
次を満たすものが存在することを証明することである：
\begin{enumerate}
\item $T'$ は $|f| : |X'| \to |X|$ による $T$ の逆像である；
\item $f : X' \to X$ は $U'$ を $X$ の開部分空間 $U$ へ同型に写す；
\item $g = a \circ f_{/T}$ である。ここで
$f_{/T} : X'_{/T'} \to X_{/T}$ は誘導される射である。
\end{enumerate}
$(f : X' \to X, T, a)$ を {\it 解}と呼ぶことにする。

\medskip\noindent
Artin の方針に従い、$S$ 上の局所 Noether スキームの圏上の関手 $F$ を構成し、
命題 \ref{artin-proposition-spaces-diagonal-representable-noetherian} を用いて $F$ が
代数空間であることを示し、$X = F$ とおけばよいことを証明する。

\begin{remark}
\label{artin-remark-G-rings}
特に、望む結果がすべての状況 \ref{artin-situation-contractions} に対して成り立つことは
証明できない。これは $S$ の局所環が G-環であると仮定する必要があるからである。
一般の場合に結果を証明できる場合、または反例がある場合には、
\href{mailto:stacks.project@gmail.com}{stacks.project@gmail.com} まで知らせてほしい。
\end{remark}

\noindent
状況 \ref{artin-situation-contractions} において、$V$ を $S$ 上の局所 Noether スキームとする。
$V$ における関手 $F$ の値は、次の条件を満たすすべての三つ組
$$
(Z, u' : V \setminus Z \to U', \hat x : V_{/Z} \to W)
$$
とする：
\begin{enumerate}
\item $Z \subset V$ は閉部分集合である；
\item $u' : V \setminus Z \to U'$ は $S$ 上の射である；
\item $\hat x : V_{/Z} \to W$ は $S$ 上の形式代数空間の進射である；
\item $u'$ と $\hat x$ は両立する（下を参照せよ）。
\end{enumerate}
両立条件は次のものである。形式的修正 $g$ を引き戻すことにより、形式的修正
$$
X'_{/T'} \times_{g, W, \hat x} V_{/Z} \longrightarrow V_{/Z}
$$
を得る。形式空間の代数化の補題
\ref{restricted-lemma-base-change-formal-modification} を参照せよ。
膨張に関する主定理（形式空間の代数化の定理
\ref{restricted-theorem-dilatations}）により、$V \setminus Z$ 上で同型となる代数空間の
一意的な固有射 $V' \to V$ で、$V'_{/Z} \to V_{/Z}$ が上に表示した射と同型となるものが
存在する。言い換えると、ある射 $\hat x' : V'_{/Z} \to X'_{/T'}$ に対して、
形式代数空間の Cartesian 図式
$$
\xymatrix{
V'_{/Z} \ar[r] \ar[d]_{\hat x'} & V_{/Z} \ar[d]^{\hat x} \\
X'_{/T'} \ar[r]^g & W
}
$$
がある。以下、特に断らず $V \setminus Z$ を $V'$ の開部分空間とみなす。
両立条件とは、$V \setminus Z \subset V'$ および $V'_{/Z}$ 上でそれぞれ $u'$ と
$\hat x$ に制限される射 $x' : V' \to X'$ が存在することである。言い換えると、図式
$$
\xymatrix{
V \setminus Z \ar[r] \ar[d]_{u'} &
V' \ar[d]^{x'} &
V'_{/Z} \ar[l] \ar[d]^{\hat x'} \ar[r] &
V_{/Z} \ar[d]^{\hat x} \\
U' \ar[r] &
X' &
X'_{/T'} \ar[r]^g \ar[l] &
W
}
$$
が可換であることである。形式空間の代数化の補題
\ref{restricted-lemma-faithful} により、射 $x'$ は存在すれば一意的であることに注意する。
この状況を「{\it $V' \to V$、$\hat x'$、および $x'$ が $u'$ と $\hat x$ の間の
両立性を証示する}」と表す。

\begin{remark}
\label{artin-remark-how-to-think-compatibility}
状況 \ref{artin-situation-contractions} において、$V$ を $S$ 上の局所 Noether スキームとする。
$(Z, u', \hat x)$ を上の (1)、(2)、(3) を満たす三つ組とする。両立条件 (4) を
どのように考えればよいかを説明する。ここでは $S$ 上の解析空間からなる架空の圏
$\textit{An}_S$ と、架空の解析化関手
$$
\left\{
\begin{matrix}
\text{locally Noetherian formal} \\
\text{algebraic spaces over }S
\end{matrix}
\right\}
\longrightarrow
\textit{An}_S,
\quad\quad
Y \longmapsto Y^{an}
$$
を用いるので、この説明は数学的に厳密なものではない。たとえば
$S = \Spec(k)$ 上の $Y = \text{Spf}(k[[t]])$ ならば、$Y^{an}$ は開単位円板と
考えるべきである。$Y = \Spec(k)$ ならば、$Y^{an}$ は一点である。
圏 $\textit{An}_S$ は開埋入と閉埋入をもち、閉部分の開補集合をとれるべきである。
$Y$ が与えられたとき、射 $Y_{red} \to Y$ は閉埋入
$Y_{red}^{an} \to Y^{an}$ を誘導すべきである。その開補集合を
$Y^{rig} = Y^{an} \setminus Y_{red}^{an}$ とおく。$Y$ が代数空間で
$Z \subset Y$ が閉ならば、射 $Y_{/Z} \to Y$ は開埋入
$Y_{/Z}^{an} \to Y^{an}$ を誘導し、これはさらに開埋入
$$
can : (Y_{/Z})^{rig} \longrightarrow (Y \setminus Z)^{an}
$$
を誘導すべきである。また、$g : Y' \to Y$ が局所 Noether 形式代数空間の形式的修正ならば、
誘導される射 $g^{rig} : (Y')^{rig} \to Y^{rig}$ は同型であるべきである。
$\text{An}_S$ と解析化関手が与えられれば、図式
$$
\xymatrix{
(V_{/Z})^{rig} \ar[rr]_{can} \ar[d]_{(g^{rig})^{-1} \circ \hat x^{an}} & &
(V \setminus Z)^{an} \ar[d]^{(u')^{an}} \\
(X'_{/T'})^{rig} \ar[rr]^{can} & & (X' \setminus T')^{an}
}
$$
が可換であるという要請を考えられる。これは
$g^{rig} : (X'_{T'})^{rig} \to W^{rig}$ が同型であり、
$U' = X' \setminus T'$ なので意味をもつ。最後に、解析化関手にいくつかの忠実性を
仮定すれば、この要請は上で定式化した両立条件と同値になる。
この説明によって、$u'$ と $\hat x$ の両立性を、ある可換図式の存在を要求することでは
なく、いくつかの写像が等しいという要請として読者が考えるようになることを期待する。
\end{remark}

\begin{lemma}
\label{artin-lemma-functor}
状況 \ref{artin-situation-contractions} において、$S$ 上の局所 Noether スキーム $V$ を
両立条件を満たす三つ組 $(Z, u', \hat x)$ の集合へ送り、$S$ 上の局所 Noether
スキームの射 $\varphi : V_2 \to V_1$ を写像
$$
F(\varphi) : F(V_1) \longrightarrow F(V_2)
$$
へ送る規則 $F$ を考える。この写像は $F(V_1)$ の元
$(Z_1, u'_1, \hat x_1)$ を、次で与えられる $F(V_2)$ の
$(Z_2, u'_2, \hat x_2)$ へ送る：
\begin{enumerate}
\item $Z_2 \subset V_2$ は $\varphi$ による $Z_1$ の逆像である；
\item $u'_2$ は $u'_1$ と
$\varphi|_{V_2 \setminus Z_2} : V_2 \setminus Z_2 \to V_1 \setminus Z_1$ の合成である；
\item $\hat x_2$ は $\hat x_1$ と
$\varphi_{/Z_2} : V_{2, /Z_2} \to V_{1, /Z_1}$ の合成である。
\end{enumerate}
この規則は反変関手である。
\end{lemma}

\begin{proof}
$u'_2$ と $\hat x_2$ の間の両立条件を見るため、$V'_1 \to V_1$、$\hat x'_1$、$x'_1$ が
$u'_1$ と $\hat x_1$ の間の両立性を証示するとする。
$V'_2 = V_2 \times_{V_1} V'_1$ とおき、$\hat x'_2$ を $\hat x'_1$ と
$V'_{2, /Z_2} \to V'_{1, /Z_1}$ の合成とし、$x'_2$ を $x'_1$ と
$V'_2 \to V'_1$ の合成とする。このとき $V'_2 \to V_2$、$\hat x'_2$、$x'_2$ が
$u'_2$ と $\hat x_2$ の間の両立性を証示する。詳細な確認は省略する。
\end{proof}

\begin{lemma}
\label{artin-lemma-solution}
状況 \ref{artin-situation-contractions} において解 $(f : X' \to X, T, a)$ が存在するならば、
$S$ 上の局所 Noether スキーム $V$ の圏上で関手的な全単射
$F(V) = \Mor_S(V, X)$ がある。
\end{lemma}

\begin{proof}
$V$ を $S$ 上の局所 Noether スキームとし、$x : V \to X$ を $S$ 上の射とする。
このとき、次のように $F(V)$ の元 $(Z, u', \hat x)$ を得る：
\begin{enumerate}
\item $Z \subset V$ は $x$ による $T$ の逆像である；
\item $u' : V \setminus Z \to U' = U$ は $x$ の $V \setminus Z$ への制限である；
\item $\hat x : V_{/Z} \to W$ は $x_{/Z} : V_{/Z} \to X_{/T}$ と同型
$a : X_{/T} \to W$ の合成である。
\end{enumerate}
この三つ組は両立条件を満たす。実際、$V' = V \times_{x, X} X'$ ととれ、
$\hat x'$ を射影 $x' : V' \to X'$ の完備化ととれる。

\medskip\noindent
逆に、$F(V)$ の元 $(Z, u', \hat x)$ が与えられているとする。$u'$ と $\hat x$ と
両立する一意的な射 $x : V \to X$ があると主張する。実際、$V' \to V$、
$\hat x'$、$x'$ が $u'$ と $\hat x$ の間の両立性を証示するとする。このとき形式空間の
代数化の命題 \ref{restricted-proposition-glue-modification} は、$V_{/Z}$ 上で
$\hat x$ と一致し、$V'$ 上で $x'$ と一致する（したがって特に $V \setminus Z$ 上で
$u'$ と一致する）一意的な射 $x : V \to X$ を得るためにまさに必要な結果である。

\medskip\noindent
上の二つの構成がそれぞれの定義域の間の互いに逆な全単射を定めることの確認は省略する。
\end{proof}

\begin{lemma}
\label{artin-lemma-functor-is-solution}
状況 \ref{artin-situation-contractions} において、$S$ 上局所有限型な代数空間 $X$ と、
$S$ 上の局所 Noether スキーム $V$ の圏上の関手的な全単射
$F(V) = \Mor_S(V, X)$ が存在するならば、$X$ は解である。
\end{lemma}

\begin{proof}
状況 \ref{artin-situation-contractions} の直後に挙げた性質 (1)、(2)、(3) をもつ固有射
$f : X' \to X$、閉部分集合 $T \subset |X|$、および同型 $a : X_{/T} \to W$ を
構成しなければならない。

\medskip\noindent
台圏を注意深く一致させるため、本証明の議論はいくらか細かい。本段落では、
大胆な読者がより臆することなく進む方法を説明する。すなわち、関手 $F$ の定義を
$S$ 上のすべての局所 Noether 代数空間へ拡張してよい。そうすれば、これらの
代数空間の圏上の関手として $F$ と $X$ が一致する、すなわち $X$ が $F$ を
表現すると結論してよい。次に $F(X)$ の普遍対象 $(T, u', \hat x)$ を考える。
すると読者は、$u'$ と $\hat x$ の間の両立性を証示する三つ組
$X'' \to X$、$\hat x'$、$x'$ に対して、射 $x' : X'' \to X'$ が同型であることを
見いだす。この $x'$ を逆にすることで $f : X' \to X$ が得られる。
最後に、既に $T \subset |X|$ があり、$\hat x$ が同型であることを示せる。
これを最後の構成要素 $a$ として用いられる。

\medskip\noindent
$h_X(-) = \Mor_S(-, X)$ を、節 \ref{artin-section-noetherian} の圏
$(\textit{Noetherian}/S)_\etale$ に制限した $X$ の点の関手と書く。
空間の極限の注意 \ref{spaces-limits-remark-limit-preserving} により、代数空間
$X$ と $X'$ は極限を保つ。したがってその制限 $h_X$ と $h_{X'}$ も極限を保つ。
ゆえに $f$ を構成するには、変換 $h_{X'} \to h_X = F$ を構成すれば十分である。
補題 \ref{artin-lemma-canonical-extension} を参照せよ。そこで $V \to S$ を
$(\textit{Noetherian}/S)_\etale$ の対象とし、
$\tilde x : V \to X'$ を $h_{X'}(V)$ の元とする。このとき、次のように
$F(V)$ の元 $(Z, u', \hat x)$ を得る：
\begin{enumerate}
\item $Z \subset V$ は $\tilde x$ による $T'$ の逆像である；
\item $u' : V \setminus Z \to U'$ は $\tilde x$ の $V \setminus Z$ への制限である；
\item $\hat x : V_{/Z} \to W$ は $x_{/Z} : V_{/Z} \to X'_{/T'}$ と
$g : X'_{/T'} \to W$ の合成である。
\end{enumerate}
この三つ組は両立条件を満たす。まず $V' \to V$ と
$\hat x' : V'_{/Z'} \to X'_{/T'}$ は常に得られる（補題 \ref{artin-lemma-functor} に先立つ
議論を参照せよ）。次に $x' : V' \to X'$ を、$V' \to V$ と射
$\tilde x : V \to X'$ の合成として定義すればよい。これでうまくいくことの確認は省略する。

\medskip\noindent
$\xi : V \to X$ がエタール射で $V$ がスキームならば、
$\xi = (Z, u', \hat x) \in F(V) = h_X(V) = X(V)$ を得る。もちろん、
$\varphi : V' \to V$ がさらにスキームのエタール射ならば、$F(V')$ へ引き戻した
$(Z, u', \hat x)$ は $\xi \circ \varphi$ に対応する。閉部分集合 $T \subset |X|$ は、
$\xi = (Z, u', \hat x)$ に対する $\xi : V \to X$ が $T$ を $Z$ へ引き戻すような
閉部分集合として定義する。

\medskip\noindent
$S$ 上の Noether スキーム $V$ と、上のような $(Z, u', \hat x)$ に対応する射
$\xi : V \to X$ を考える。このとき、$\xi(V)$ が集合論的に $T$ に含まれることと、
位相空間として $V = Z$ であることは同値である。したがって関手として
$X_{/T}$ は $W$ と一致することがわかる。これにより同型 $a : X_{/T} \to W$ を得る。
（ここでは小さな詳細を省略した。それは、局所 Noether 形式代数空間 $X_{/T}$ と
$W$ について、$S$ 上の局所 Noether スキームにおける値をとった後に同型を得ることを
確認すれば十分であるという点である。）

\medskip\noindent
条件 (1)、(2)、(3) の証明は省略する。
\end{proof}

\begin{remark}
\label{artin-remark-diagonal}
状況 \ref{artin-situation-contractions} において、$V$ を $S$ 上の局所 Noether スキームとする。
$i = 1, 2$ に対して $(Z_i, u'_i, \hat x_i) \in F(V)$ とする。
$i = 1, 2$ に対して、$V'_i \to V$、$\hat x'_i$、$x'_i$ が
$u'_i$ と $\hat x_i$ の間の両立性を証示するとする。

\medskip\noindent
$V' = V'_1 \times_V V'_2$ とおく。$E' \to V'$ を二つの射
$$
V' \to V'_1 \xrightarrow{x'_1} X'
\quad\text{and}\quad
V' \to V'_2 \xrightarrow{x'_2} X'
$$
の等化子とする。$Z = Z_1 \cap Z_2$ とおく。$E_W \to V_{/Z}$ を二つの射
$$
V_{/Z} \to V_{/Z_1} \xrightarrow{\hat x_1} W
\quad\text{and}\quad
V_{/Z} \to V_{/Z_2} \xrightarrow{\hat x_2} W
$$
の等化子とする。$E' \to V$ は分離かつ局所有限型であり、$E_W$ は $V$ 上分離な
局所 Noether 形式代数空間であることに注意する。関係するさまざまな射の両立性から、
次がわかる：
\begin{enumerate}
\item $\Im(E' \to V) \cap (Z_1 \cup Z_2)$ は $Z = Z_1 \cap Z_2$ に含まれる；
\item 射 $E' \times_V (V \setminus Z) \to V \setminus Z$ は単射であり、
$V \setminus (Z_1 \cup Z_2)$ への $u'_1$ と $u'_2$ の制限の等化子に等しい；
\item 射 $E'_{/Z} \to V_{/Z}$ は $E_W$ を経由し、図式
$$
\xymatrix{
E'_{/Z} \ar[r] \ar[d] & X'_{/T'} \ar[d]^g \\
E_W \ar[r] & W
}
$$
は Cartesian である。特に、射 $E'_{/Z} \to E_W$ は $g$ の基底変換として
形式的修正である；
\item $E'$、$(E' \to V)^{-1}Z$、および $E'_{/Z} \to E_W$ は、
局所 Noether スキーム $V$ を基礎スキームとする状況 \ref{artin-situation-contractions} における
三つ組である；
\item 局所 Noether スキームの射 $\varphi : A \to V$ が与えられたとき、次は同値である：
\begin{enumerate}
\item $(Z_1, u'_1, \hat x_1)$ と $(Z_2, u'_2, \hat x_2)$ は
$F(A)$ の同じ元へ制限される；
\item $A \setminus \varphi^{-1}(Z) \to V \setminus Z$ は
$E' \times_V (V \setminus Z)$ を経由し、
$A_{/\varphi^{-1}(Z)} \to V_{/Z}$ は $E_W$ を経由する。
\end{enumerate}
\end{enumerate}
補題 \ref{artin-lemma-solution} と \ref{artin-lemma-functor-is-solution} を用いると、(4) の三つ組に
対する解 $E \to V$ が存在するならば、$E$ は $V$ 上の局所 Noether スキームの圏上で
$F \times_{\Delta, F \times F} V$ を表現すると結論する。
\end{remark}

\begin{lemma}
\label{artin-lemma-closed-immersion}
状況 \ref{artin-situation-contractions} において、次を満たす閉部分集合 $Z \subset S$ が
与えられていると仮定する：
\begin{enumerate}
\item $X'$ における $Z$ の逆像は $T'$ である；
\item $U' \to S \setminus Z$ は閉埋入である；
\item $W \to S_{/Z}$ は閉埋入である。
\end{enumerate}
このとき解 $(f : X' \to X, T, a)$ が存在し、さらに $X \to S$ は閉埋入である。
\end{lemma}

\begin{proof}
$X \cap (S \setminus Z) = U'$ かつ $X_{/Z} = W$ となる閉部分スキーム
$X \subset S$ があると仮定する。このとき $X$ は関手 $F$ を表現し
（いくつかの詳細は省略する）、したがって解である。$X$ を見いだすことは明らかに
$S$ 上局所的な問題である。このようにして、次段落で論じる場合へ帰着する。

\medskip\noindent
$S = \Spec(A)$ はアフィンであると仮定する。$I \subset A$ を $Z$ を切り出す根基
イデアルとし、$I = (f_1, \ldots, f_r)$ と書く。仮定により次が与えられている：
\begin{enumerate}
\item 閉埋入 $U' \to S \setminus Z$ はイデアル $J_i \subset A[1/f_i]$ を定め、
$J_i$ と $J_j$ は $A[1/f_if_j]$ において同じイデアルを生成する；
\item 閉埋入 $W \to S_{/Z}$ は、$A$ の $I$-進完備化 $A^\wedge$ のあるイデアル
$J' \subset A^\wedge$ に対する写像
$\text{Spf}(A^\wedge/J') \to \text{Spf}(A^\wedge)$ である。
\end{enumerate}
証明を完了するには、$J_i = J[1/f_i]$ かつ $J' = JA^\wedge$ となるイデアル
$J \subset A$ を見いだす必要がある。代数続論の命題
\ref{more-algebra-proposition-equivalence} により、すべての $i$ に対して $J_i$ と $J'$ が
$A^\wedge[1/f_i]$ において同じイデアルを生成することを示せば十分である。

\medskip\noindent
$A' = H^0(X', \mathcal{O})$ は有限 $A$-代数であり、その形成は平坦な基底変換と
可換であることを思い出そう（空間のコホモロジーの補題
\ref{spaces-cohomology-lemma-proper-over-affine-cohomology-finite} および
\ref{spaces-cohomology-lemma-flat-base-change-cohomology}）。
$J'' = \Ker(A \to A')$ と書く\footnote{読者の予想に反して、イデアル $J$ と $J''$ は
一般には一致しない。}。$A[1/f_i]$ のスペクトルへの基底変換から、
$J_i = J''A[1/f_i]$ が従う。可換図式
$$
\xymatrix{
X' \ar[d] &
X'_{/T'} \times_{S_{/Z}} \text{Spf}(A^\wedge) \ar[l] \ar[d] &
X'_{/T'} \times_W \text{Spf}(A^\wedge/J') \ar@{=}[l] \ar[d] \\
\Spec(A) &
\text{Spf}(A^\wedge) \ar[l] &
\text{Spf}(A^\wedge/J') \ar[l]
}
$$
があることに注意する。中央の垂直射は、明らかな閉部分集合に沿った左の垂直射の
完備化である。形式関数定理により
$$
(A')^\wedge = \Gamma(X' \times_S \Spec(A^\wedge), \mathcal{O}) =
\lim H^0(X' \times_S \Spec(A/I^n), \mathcal{O})
$$
である。空間のコホモロジーの定理
\ref{spaces-cohomology-theorem-formal-functions} を参照せよ。図式から、$J'$ は
$(A')^\wedge$ において零へ写ると結論する。したがって
$J' \subset J'' A^\wedge$ である。射
$$
X'_{/T'} \to
\text{Spf}(A^\wedge/J''A^\wedge) \to
\text{Spf}(A^\wedge/J') = W
$$
を考える。合成 $g$ は形式的修正（特に rig-エタールかつ rig-全射）であり、第二の射は
閉埋入（特に進的単射）である。したがって
$X'_{/T'} \to \text{Spf}(A^\wedge/J''A^\wedge)$ は rig-全射かつ rig-エタールである。
形式空間の代数化の補題
\ref{restricted-lemma-rig-surjective-alternative-permanence} および
\ref{restricted-lemma-rig-etale-alternative-permanence} を参照せよ。
形式空間の代数化の補題 \ref{restricted-lemma-rig-etale-descent} および
\ref{restricted-lemma-permanence-rig-surjective} を適用すると、
$\text{Spf}(A^\wedge/J''A^\wedge) \to W$ は rig-エタールかつ rig-全射であると結論する。
形式空間の代数化の補題
\ref{restricted-lemma-closed-immersion-rig-smooth-rig-surjective} により、ある $n > 0$ に
対して $I^n J'' A^\wedge \subset J'$ である。したがって
$J'' A^\wedge[1/f_i] = J' A^\wedge[1/f_i]$ であり、望みどおりすべての $i$ に対して
$J_i A^\wedge[1/f_i] = J' A^\wedge[1/f_i]$ を得る。
\end{proof}

\begin{lemma}
\label{artin-lemma-diagonal-contractions}
状況 \ref{artin-situation-contractions} において、$X' \to S$ と $W \to S$ は分離であると
仮定する。このとき対角射 $\Delta : F \to F \times F$ は閉埋入で表現可能である。
\end{lemma}

\begin{proof}
補題 \ref{artin-lemma-closed-immersion} と注意 \ref{artin-remark-diagonal} の議論を組み合わせればよい。
\end{proof}

\begin{lemma}
\label{artin-lemma-sheaf}
状況 \ref{artin-situation-contractions} において、関手 $F$ は $S$ 上の局所 Noether スキームの
すべてのエタール被覆に対して層条件を満たす。
\end{lemma}

\begin{proof}
省略する。ヒント：射はエタール局所的に定義できる。
\end{proof}

\begin{lemma}
\label{artin-lemma-limit-preserving}
状況 \ref{artin-situation-contractions} において、関手 $F$ は極限を保つ。すなわち、
$S$ 上の Noether アフィンスキームの任意の有向極限 $V = \lim V_\lambda$ に対して
$F(V) = \colim F(V_\lambda)$ である。
\end{lemma}

\begin{proof}
これは途方もなく長い証明である。その大部分は、極限とエタール局所化に関する
標準的な議論からなる。何か興味深いことが起こる証明の最後の部分まで読み飛ばすよう、
読者に勧める。

\medskip\noindent
$V = \lim_{\lambda \in \Lambda} V_i$ を $S$ 上のスキームの有向極限で、$V$ と
$V_\lambda$ が Noether、遷移射がアフィンであるものとする。スキームの極限に関する
内容については、極限の節 \ref{limits-section-limits} を参照せよ。
$\colim F(V_\lambda) \to F(V)$ が全単射であることを証明する。

\medskip\noindent
単射性の証明：記法。
$\lambda \in \Lambda$ とし、$\xi_{\lambda, 1}, \xi_{\lambda, 2} \in F(V_\lambda)$ を、
$F(V)$ の同じ元へ制限される元とする。
$\xi_{\lambda, 1} = (Z_{\lambda, 1}, u'_{\lambda, 1},
\hat x_{\lambda, 1})$ および
$\xi_{\lambda, 2} = (Z_{\lambda, 2}, u'_{\lambda, 2}, \hat x_{\lambda, 2})$ と書く。

\medskip\noindent
単射性の証明：$Z_{\lambda, i}$ の一致。
$Z_{\lambda, 1}$ と $Z_{\lambda, 2}$ は $V$ の同じ閉部分集合へ制限されるので、
$i$ を大きくした後 $Z_{\lambda, 1} = Z_{\lambda, 2}$ と仮定してよい。
極限の補題 \ref{limits-lemma-inverse-limit-top} および位相の補題
\ref{topology-lemma-describe-limits} を参照せよ。共通の値を
$Z_\lambda \subset V_\lambda$ と書き、$\mu \geq \lambda$ に対して $V_\mu$ における
逆像を $Z_\mu \subset V_\mu$、$V$ における逆像を $Z$ と書く。
$Z$ と $Z_\mu$ を被約閉部分スキームとみなせば、スキームとして
$Z = \lim_{\mu \geq \lambda} Z_\mu$ であることを以下で用いる。

\medskip\noindent
単射性の証明：$u'_{\lambda, i}$ の一致。
$U'$ は $S$ 上局所有限型であり、$u'_{\lambda, 1}$ と $u'_{\lambda, 2}$ の
$V \setminus Z$ への制限は同じなので、$\lambda$ を大きくした後
$u'_{\lambda, 1} = u'_{\lambda, 2}$ と仮定してよい。極限の命題
\ref{limits-proposition-characterize-locally-finite-presentation} を参照せよ。
共通の値を $u'_\lambda$ と書き、その $V \setminus Z$ への制限を $u'$ と書く。

\medskip\noindent
単射性の証明：言い換え。
この時点で
$\xi_{\lambda, 1} = (Z_\lambda, u'_\lambda, \hat x_{\lambda, 1})$ および
$\xi_{\lambda, 2} = (Z_\lambda, u'_\lambda, \hat x_{\lambda, 2})$ である。
証明のこの部分で直面する主な問題は、$\lambda$ を大きくした後に射
$\hat x_{\lambda, 1}$ と $\hat x_{\lambda, 2}$ が同じになることを示すことである。

\medskip\noindent
単射性の証明：$\hat x_{\lambda, i}|_{Z_\lambda}$ の一致。
射
$\hat x_{\lambda, 1}|_{Z_\lambda}, \hat x_{\lambda, 2}|_{Z_\lambda} :
Z_\lambda \to W_{red}$ を考える。これらの射は同じ射 $Z \to W_{red}$ へ制限される。
$W_{red}$ は $S$ 上局所有限型なスキームなので、極限の命題
\ref{limits-proposition-characterize-locally-finite-presentation} を用いると、
$\lambda$ をより大きな添字で置き換えた後
$\hat x_{\lambda, 1}|_{Z_\lambda} = \hat x_{\lambda, 2}|_{Z_\lambda} :
Z_\lambda \to W_{red}$ と仮定してよい。

\medskip\noindent
単射性の証明：終結。
次に、注意 \ref{artin-remark-diagonal} の議論を $V_\lambda$ と二つの元
$\xi_{\lambda, 1}, \xi_{\lambda, 2} \in F(V_\lambda)$ に適用する。これにより次を得る：
\begin{enumerate}
\item $e_\lambda : E_\lambda' \to V_\lambda$ は分離かつ局所有限型である；
\item $e_\lambda^{-1}(V_\lambda \setminus Z_\lambda) \to
V_\lambda \setminus Z_\lambda$ は同型である；
\item $E_{W, \lambda} \to V_{\lambda, /Z_\lambda}$ は
$\hat x_{\lambda, 1}$ と $\hat x_{\lambda, 2}$ の等化子である単射である；
\item $E'_{\lambda, /Z_\lambda} \to E_{W, \lambda}$ は形式的修正である。
\end{enumerate}
(2) は、注意 \ref{artin-remark-diagonal} の主張 (2) と、上で
$u'_{\lambda, 1} = u'_{\lambda, 2} = u'_\lambda$ を得た準備から従う。
$\hat x_{\lambda, 1}|_{Z_\lambda} = \hat x_{\lambda, 2}|_{Z_\lambda}$ なので
$Z_\lambda = (V_{\lambda, /Z_\lambda})_{red}$ は $E_{W, \lambda}$ を経由する。
したがって形式空間の補題 \ref{formal-spaces-lemma-monomorphism-iso-over-red} により、
$E_{W, \lambda} \to V_{\lambda, /Z_\lambda}$ は閉埋入である。次に、注意
\ref{artin-remark-diagonal} の主張 (4) と、$V_\lambda$ 上の三つ組
$E_\lambda'$、$e_\lambda^{-1}(Z_\lambda)$、
$E'_{\lambda, /Z_\lambda} \to E_{W, \lambda}$ に適用した補題
\ref{artin-lemma-closed-immersion} から、この三つ組の解である閉埋入
$E_\lambda \to V_\lambda$ が存在することがわかる。次に注意
\ref{artin-remark-diagonal} の主張 (5) と補題 \ref{artin-lemma-solution} を用いると、
$E_\lambda$ は $\xi_{\lambda, 1}$ と $\xi_{\lambda, 2}$ の「等化子」である。
特に $V \to V_\lambda$ は $E_\lambda$ を経由する。極限の命題
\ref{limits-proposition-characterize-locally-finite-presentation} をもう一度用いると、
$V_\mu \to V_\lambda$ が $E_\lambda$ を経由するような $\mu \geq \lambda$ が見つかる。
したがって望みどおり、$\xi_{\lambda, 1}$ と $\xi_{\lambda, 2}$ の $V_\mu$ への
引き戻しは同じである。

\medskip\noindent
全射性の証明：主張。
$\xi = (Z, u', \hat x)$ を $F(V)$ の元とする。$\xi$ へ制限される元
$\xi_\lambda \in F(V_\lambda)$ と $\lambda \in \Lambda$ を見いださなければならない。

\medskip\noindent
全射性の証明：問題はエタール局所的である。
証明の前半で示した一意性と補題 \ref{artin-lemma-sheaf} における $F$ の層条件により、
問題はエタール位相について $V$ 上局所的である。より正確には $v \in V$ とする。
エタール射 $(\tilde V, \tilde v) \to (V, v)$ と、ある $\lambda$、あるエタール射
$\tilde V_\lambda \to V_\lambda$、およびある元
$\tilde \xi_\lambda \in F(\tilde V_\lambda)$ で、
$\tilde V = \tilde V_\lambda \times_{V_\lambda} V$ かつ
$\xi|_U = \tilde \xi_\lambda|_U$ となるものを見いだせば十分であると主張する。
この主張の詳しい証明は省略する\footnote{これを証明するには、射
$\tilde V \to V$ の集まりを有限エタール被覆へまとめ、対応する射
$\tilde V_\lambda \to V_\lambda$ も（$\lambda$ を大きくした後）エタール被覆をなすことを
示す。次に単射性を用いて、元 $\tilde \xi_\lambda$ が（$\lambda$ を大きくした後）
貼り合わさることを確かめ、$F$ の層条件を用いてこれらの元を $F(V_\lambda)$ の元へ
降下させる。}。

\medskip\noindent
全射性の証明：問題の言い換え。
$\tilde V$ がアフィンであるような任意のエタール射
$(\tilde V, \tilde v) \to (V, v)$ は、ある $\lambda$ に対するアフィンな
$\tilde V_\lambda$ をもつエタール射 $\tilde V_\lambda \to V_\lambda$ の基底変換である。
たとえば位相の補題 \ref{topologies-lemma-limit-fppf-topology} を参照せよ。
$\tilde V_\lambda$ が与えられれば
$\tilde V = \lim_{\mu \geq \lambda} \tilde V_\lambda \times_{V_\lambda} V_\mu$ である。
したがって、$\tilde V$ がアフィンであるエタール射
$(\tilde V, \tilde v) \to (V, v)$ が与えられたとき、$(V, v)$ を
$(\tilde V, \tilde v)$ で、$\xi$ を $\xi$ の $\tilde V$ への制限で置き換えてよい。

\medskip\noindent
全射性の証明：基礎がアフィンの場合への帰着。特に、$\tilde S \subset S$ を
$v \in V$ が $\tilde S$ の点へ写るようなアフィン開部分スキームとする。
前段落に従って $V$ を $\tilde V = \tilde S \times_S V$ で置き換えてよい。
もちろん、この置き換えを行うならば、$\tilde S$ 上の局所 Noether スキームの圏へ
制限した関手 $F$ について問題を解けば十分である。この関手は当然、状況全体を
$\tilde S$ へ基底変換して得られる関手である。したがって証明の残りでは、
$S = \Spec(R)$ は Noether アフィンスキームであると仮定する。

\medskip\noindent
全射性の証明：容易な場合。
$v \in V \setminus Z$ ならば、$\tilde V = V \setminus Z$ ととれる。
極限の補題 \ref{limits-lemma-descend-opens} により、これはある $\lambda$ に対して
開部分スキーム $\tilde V_\lambda \subset V_\lambda$ へ降下する。
次に $\lambda$ を大きくした後、$u'$ へ制限される射
$u'_\lambda : \tilde V_\lambda \to U'$ があると仮定してよい。
$\tilde \xi_\lambda = (\emptyset, u'_\lambda, \emptyset)$ ととれば、望む
$F(\tilde V_\lambda)$ の元を得る。

\medskip\noindent
全射性の証明：困難な場合と、$W$ がアフィンの場合への帰着。
最も困難な場合は $v \in Z \subset V$ を考えるときに生じる。これを $W$ が
アフィン形式スキームである場合へ帰着できると主張する。この議論は読み飛ばすよう
読者に勧める\footnote{Artin はこの問題を回避する方法でこの補題を証明しており、
その結果、単射性を先に証明する必要がない。すなわち Artin は一貫して、登場する
すべての空間の有限アフィン・エタール被覆を用い、証明中それらの間の写像を追跡する。
後から考えれば、この方法の方がここでの方法より望ましいかもしれない。}。
像の被約部分上で $\hat x : V_{/Z} \to W$ による $v$ の像が
$\tilde W \to W$ の像に入るようなエタール射 $\tilde W \to W$ を選べる。
ここで $\tilde W$ はアフィン形式代数空間である。このとき射
$$
p : \tilde W \times_{W, g} X'_{/T'} \longrightarrow X'_{/T'}
$$
および
$$
q : \tilde W \times_{W, \hat x} V_{/Z} \to V_{/Z}
$$
は局所 Noether 形式代数空間のエタール射である。形式空間の代数化の定理
\ref{restricted-theorem-dilatations-general} の（容易な場合）により、局所有限型で、
$U'$ 上で同型であり、$\tilde X'_{/T'} \to X'_{/T'}$ が $p$ と同型となるような
代数空間の射 $\tilde X' \to X'$ が存在する。形式空間の代数化の補題
\ref{restricted-lemma-output-etale} により、射 $\tilde X' \to X'$ はエタールである。
$\tilde T' \subset |\tilde X'|$ を $T'$ の逆像とし、$\tilde U' \subset \tilde X'$ を
その補開部分空間とする。$\tilde g' : \tilde X'_{/\tilde T'} \to \tilde W$ を、
$g$ の $\tilde W \to W$ による基底変換である形式的修正とする。このとき
$$
\tilde X',\ \tilde T',\ \tilde U',\ \tilde W,
\ \tilde g : \tilde X'_{/\tilde T'} \to \tilde W
$$
は状況 \ref{artin-situation-contractions} の別の例である。この三つ組から構成される関手を
$\tilde F$ と書く。射 $\tilde X' \to X'$ と $\tilde W \to W$ を明らかな方法で用いて
構成される関手の変換
$$
\tilde F \longrightarrow F
$$
がある。詳細は省略する。

\medskip\noindent
全射性の証明：困難な場合と、$W$ がアフィンの場合への帰着、その 2。
上で用いたものと同じ定理により、局所有限型で、$V \setminus Z$ 上で同型であり、
$\tilde V_{/Z} \to V_{/Z}$ が $q$ と同型となるような代数空間の射
$\tilde V \to V$ が存在する。$\tilde Z \subset \tilde V$ を $Z$ の逆像とする。
形式空間の代数化の補題 \ref{restricted-lemma-output-etale} および
\ref{restricted-lemma-output-separated} により、射 $\tilde V \to V$ はエタールかつ
分離である。特に $\tilde V$ は（局所 Noether）スキームである。たとえば空間の射の命題
\ref{spaces-morphisms-proposition-locally-quasi-finite-separated-over-scheme} を参照せよ。
射 $u'$ があり、これを射
$$
\tilde u' : \tilde V \setminus \tilde Z \longrightarrow \tilde U'
$$
とみなせる。ここで $\tilde U' \subset \tilde X'$ は $U'$ へ同型に写る開部分である。
また射
$$
\tilde {\hat x} :
\tilde V_{/\tilde Z} = \tilde W \times_{W, \hat x} V_{/Z}
\longrightarrow
\tilde W
$$
がある。ここでは単に射影を用いる。したがって三つ組
$$
\tilde \xi = (\tilde Z, \tilde u', \tilde {\hat x}) \in \tilde F(\tilde V)
$$
を得る。両立条件の証明は省略する。ヒント：$V' \to V$、$\hat x'$、$x'$ が
$u'$ と $\hat x$ の間の両立性を証示するならば、
$\tilde V' = V' \times_V \tilde V$ とおく。これは射 $\tilde{\hat x}'$ と
$\tilde x'$ を備え、これでうまくいくことを示せる。変換 $\tilde F \to F$ による
$\tilde \xi$ の像は、$\xi$ の $\tilde V$ への制限である。

\medskip\noindent
全射性の証明：困難な場合と、$W$ がアフィンの場合への帰着、その 3。
$\tilde W \to W$ の選択により、$V$ における像が選んだ点 $v \in V$ を含むような
アフィン開集合 $\tilde V_{open} \subset \tilde V$ が存在する
（記号が尽きかけている）。次段落で調べる場合と先の注意により、
$\tilde \xi|_{\tilde V_{open}}$ を、$\tilde V_{\lambda, open}$ 上の
$\tilde F$ のある元 $\tilde \xi_\lambda$ へ降下できる。ここで
$\tilde V_{\lambda, open} \to V_\lambda$ はエタール射で、その $V$ への基底変換は
$\tilde V_{open}$ である。
関手の変換 $\tilde F \to F$ を適用すると、求めていた
$F(\tilde V_{\lambda, open})$ の元を得る。これにより次段落の場合へ帰着する。

\medskip\noindent
全射性の証明：アフィン $W$ の場合。$v \in Z \subset V$ であり、$W$ はアフィン
形式代数空間である。次を思い出そう：
$$
\xi = (Z, u', \hat x) \in F(V)
$$
引き続き $V$ を $v$ のエタール近傍で置き換えてよい。特に $V$ と $V_\lambda$ は
アフィンであると仮定する。

\medskip\noindent
全射性の証明：$Z$ の降下。ある $\lambda$ と閉部分スキーム
$Z_\lambda \subset V_\lambda$ で、$Z$ が $Z_\lambda$ の $V$ への基底変換となるものを
見いだせる。極限の補題 \ref{limits-lemma-descend-finite-presentation} を参照せよ。
警告：$Z_\lambda$ が $V_\lambda$ の被約閉部分スキームであることはわからず、
一般にはそうならない。$\mu \geq \lambda$ に対して、$V_\mu$ におけるスキーム論的逆像を
$Z_\mu \subset V_\mu$ と書く。以下で、スキームとして
$Z = \lim_{\mu \geq \lambda} Z_\mu$ であることを用いる。

\medskip\noindent
全射性の証明：$u'$ の降下。$U'$ は $S$ 上局所有限型なので、$\lambda$ を大きくした後、
$V \setminus Z$ への制限が $u'$ となる射
$u'_\lambda : V_\lambda \setminus Z_\lambda \to U'$ が存在すると仮定してよい。
極限の命題 \ref{limits-proposition-characterize-locally-finite-presentation} を参照せよ。
$\mu \geq \lambda$ に対して、$u'_\lambda$ の
$V_\mu \setminus Z_\mu$ への制限を $u'_\mu$ と書く。

\medskip\noindent
全射性の証明：証示データの降下。$V' \to V$、$\hat x'$、$x'$ が $u'$ と $\hat x$ の
間の両立性を証示するとする。上と同じ参照を用いると、$\lambda$ を大きくした後、
$V$ への基底変換が $V' \to V$ となる有限型の射
$V'_\lambda \to V_\lambda$ が存在すると仮定してよい。さらに $\lambda$ を大きくすれば、
$V'_\lambda \to V_\lambda$ は固有であると仮定してよい（極限の補題
\ref{limits-lemma-eventually-proper}）。次に、$V'_\lambda \to V_\lambda$ は
$V_\lambda \setminus Z_\lambda$ 上で同型であると仮定してよい（極限の補題
\ref{limits-lemma-descend-isomorphism}）。次に、$V'$ への制限が $x'$ となる射
$x'_\lambda : V'_\lambda \to X'$ があると仮定してよい。もう一度 $\lambda$ を大きくすれば、
$x'_\lambda$ は $V_\lambda \setminus Z_\lambda$ 上で $u'_\lambda$ と一致すると
仮定してよい。$\mu \geq \lambda$ に対して、$V'_\lambda$ の基底変換と
$x'_\lambda$ の制限をそれぞれ $V'_\mu$、$x'_\mu$ と書く。

\medskip\noindent
全射性の証明：代数。$W = \text{Spf}(B)$、$V = \Spec(A)$ と書き、
$\mu \geq \lambda$ に対して $V_\mu = \Spec(A_\mu)$ と書く。
$Z_\mu$ と $Z$ を切り出すイデアルをそれぞれ $I_\mu \subset A_\mu$、$I \subset A$ と書く。
このとき $I_\lambda A_\mu = I_\mu$ および $I_\lambda A = I$ である。
射 $\hat x$ は連続環準同型
$$
(\hat x)^\sharp : B \longrightarrow A^\wedge
$$
を定め、またそれによって定まる。ここで $A^\wedge$ は $A$ の $I$-進完備化である。
証明を完了するには、この写像が十分大きなある $\mu$ に対して
$A_\mu^\wedge$ への写像へ降下することを示さなければならない。ここで
$A_\mu^\wedge$ は $A_\mu$ の $I_\mu$-進完備化である。これは自明でない事実である。
Artin は論文 \cite{ArtinII} で次のように記している：「データ (3.5) は直極限と
可換でない $I$-進完備化を含むので、確認はいくらか繊細である。これは解析における
収束証明の代数的類似物である。」

\medskip\noindent
全射性の証明：代数、さらに環。次のように書く：
$$
C_\mu = \Gamma(V'_\mu, \mathcal{O})
\quad\text{and}\quad
C = \Gamma(V', \mathcal{O})
$$
$V' \to V$ と $V'_\mu \to V_\mu$ は固有射なので、$A \to C$ と
$A_\mu \to C_\mu$ は有限環準同型であることに注意する。空間のコホモロジーの補題
\ref{spaces-cohomology-lemma-proper-over-affine-cohomology-finite} を参照せよ。
$V = \lim V_\mu$ および $V' = \lim V'_\mu$ なので、極限の補題
\ref{limits-lemma-descend-section} により
$$
A = \colim A_\mu
\quad\text{and}\quad
C = \colim C_\mu
$$
である\footnote{さまざまな射は平坦でないので、
$C_\mu = C_\lambda \otimes_{A_\lambda} A_\mu$ であるかどうかはわからない。}。
$a \in I$、それぞれ $a \in I_\mu$ である元に対して、写像 $A_a \to C_a$、
それぞれ $(A_\mu)_a \to (C_\mu)_a$ は平坦基底変換により同型である
（空間のコホモロジーの補題
\ref{spaces-cohomology-lemma-flat-base-change-cohomology}）。したがって $A \to C$ の核と
余核は $V(I)$ 上に台をもち、$A_\mu \to C_\mu$ についても同様である。
ゆえに $A \to C$ の核と余核は $I$ のある冪で零化され、$A_\mu \to C_\mu$ の核と
余核は $I_\mu$ のある冪で零化される。代数の補題
\ref{algebra-lemma-Noetherian-power-ideal-kills-module} を参照せよ。

\medskip\noindent
全射性の証明：代数、さらに環準同型。$Z_n \subset V$ を $Z$ の第 $n$ 無限小近傍とし、
$Z_{\mu, n} \subset V_\mu$ を $Z_\mu$ の第 $n$ 無限小近傍とする。
形式関数定理（空間のコホモロジーの定理
\ref{spaces-cohomology-theorem-formal-functions}）により
$$
C^\wedge = \lim_n H^0(V' \times_V Z_n, \mathcal{O})
\quad\text{and}\quad
C_\mu^\wedge =
\lim_n H^0(V'_\mu \times_{V_\mu} Z_{\mu, n}, \mathcal{O})
$$
である。ここで $C^\wedge$ と $C_\mu^\wedge$ はそれぞれ $I$ と $I_\mu$ に関する
完備化である。射 $x'_\mu : V'_\mu \to X'$ の完備化と射
$g : X'_{/T'} \to W$ を合成すると
$$
g \circ x'_{\mu, /Z_\mu} :
V'_{\mu, /Z_\mu} = \colim V_\mu' \times_{V_\mu} Z_{\mu, n}
\longrightarrow
W
$$
を得る。したがって上の完備化 $C_\mu^\wedge$ の記述により、連続環準同型
$$
(g \circ x'_{\mu, /Z_\mu})^\sharp : B \longrightarrow C_\mu^\wedge
$$
を得る。$V' \to V$、$\hat x'$、$x'$ が $u'$ と $\hat x$ の間の両立性を
証示することから、次の図式の可換性が従う：
$$
\xymatrix{
C_\mu^\wedge \ar[r] &
C^\wedge \\
B \ar[u]^{(g \circ x'_{\mu, /Z_\mu})^\sharp} \ar[r]^{(\hat x)^\sharp} &
A^\wedge \ar[u]
}
$$

\medskip\noindent
全射性の証明：さらなる代数的議論。有限 $A$-加群 $\Ker(A \to C)$ と
$\Coker(A \to C)$ は $I$ のある冪で零化され、同様に有限 $A_\mu$-加群
$\Ker(A_\mu \to C_\mu)$ と $\Coker(A_\mu \to C_\mu)$ は $I_\mu$ のある冪で
零化されることを思い出そう。このことから、これらの加群はその完備化に等しい。
有限 $A$-加群の圏上の $I$-進完備化は完全なので（代数の節
\ref{algebra-section-completion-noetherian} を参照せよ）、
$$
\Coker(A^\wedge \to C^\wedge) = \Coker(A \to C)
$$
を得る。核および写像 $A_\mu \to C_\mu$ についても同様である。もちろん
$$
\Ker(A \to C) = \colim \Ker(A_\mu \to C_\mu)
\qquad
\Coker(A \to C) = \colim \Coker(A_\mu \to C_\mu)
$$
も成り立つ。$S = \Spec(R)$ はアフィンであることを思い出そう。すべての射は $S$ 上の
射なので、上の環準同型はすべて $R$-代数準同型である。形式空間の代数化の補題
\ref{restricted-lemma-Noetherian-finite-type-red} により、$B$ は $R$ 上位相的有限型である。
$B$ は $b_1, \ldots, b_n$ で位相的に生成されるとする。ある $\mu$（たとえば
$\lambda$）を選び、元
$$
\text{images of }
(g \circ x'_{\mu, /Z_\mu})^\sharp(b_1)
, \ldots,
(g \circ x'_{\mu, /Z_\mu})^\sharp(b_n)
\text{ in }\Coker(A_\mu \to C_\mu)
$$
を考える。上の正方形の可換性により、これらの元の $\Coker(\alpha)$ における像は
零である。$\Coker(A \to C) = \colim \Coker(A_\mu \to C_\mu)$ であり、これらの
余核はその完備化に等しいので、$\mu$ を大きくした後、これらの像はすべて零であると
仮定してよい。これは連続準同型 $(g \circ x'_{\mu, /Z_\mu})^\sharp$ の像が
$\Im(A_\mu \to C_\mu)$ に含まれることを意味する。
$(C_\mu)^\wedge$ において
$(g \circ x'_{\mu, /Z_\mu})^\sharp(b_1)$ へ写る元
$a_{\mu, j} \in (A_\mu)^\wedge$ を選ぶ。このとき $a_{\mu, j} \in A_\mu^\wedge$ と
$(\hat x)^\sharp(b_j) \in A^\wedge$ は、上の正方形の可換性により $C^\wedge$ の
同じ元へ写る。$\Ker(A \to C) = \colim \Ker(A_\mu \to C_\mu)$ であり、これらの核は
その完備化に等しいので、$\mu$ を大きくした後、$a_{\mu, j}$ の $A^\wedge$ における像が
$(\hat x)^\sharp(b_j)$ に等しくなるように $a_{\mu, j}$ の選択を調整できる。

\medskip\noindent
全射性の証明：最後の代数的議論。$\mathfrak b \subset B$ を位相的冪零元のイデアルとする。
$J \subset R[x_1, \ldots, x_n]$ を、$h(b_1, \ldots, b_n) \in \mathfrak b$ となるような
$h(x_1, \ldots, x_n)$ からなるイデアルとする。このとき位相 $R$-代数の連続全射
$$
\Phi : R[x_1, \ldots, x_n]^\wedge \longrightarrow B,\quad
x_j \longmapsto b_j
$$
を得る。ここで左辺の完備化は $J$ に関するものである。
$R[x_1, \ldots, x_n]$ は Noether なので、$J$ の生成元
$h_1, \ldots, h_m$ を選べる。上の正方形の可換性により、
$h_j(a_{\mu, 1}, \ldots, a_{\mu, n})$ は $A_\mu^\wedge$ の元で、その $A^\wedge$ に
おける像は $IA^\wedge$ に含まれる。実際、環準同型 $(\hat x)^\sharp$ は連続であり、
$A^\wedge/IA^\wedge = A/I$ は被約なので、$IA^\wedge$ は $A^\wedge$ の位相的冪零元の
イデアルである。（Noether 環の完備化に関する結果については、代数の節
\ref{algebra-section-completion-noetherian} を参照せよ。）
$A/I = \colim A_\mu/I_\mu$ なので、$\mu$ を大きくした後、
$h_j(a_{\mu, 1}, \ldots, a_{\mu, n})$ は $I_\mu A_\mu^\wedge$ に属すると仮定してよい。
特に $A_\mu^\wedge$ の元 $h_j(a_{\mu, 1}, \ldots, a_{\mu, n})$ は
$A_\mu^\wedge$ において位相的冪零である。したがって連続 $R$-代数準同型
$$
\Psi : R[x_1, \ldots, x_n]^\wedge \longrightarrow A_\mu^\wedge,\quad
x_j \longmapsto a_{\mu, j}
$$
を得る。望む結論を得るには、$\Ker(\Phi)$ が $\Psi$ によって零化されるかを
確認する必要がある。これはそのままでは成り立たないかもしれないが、$\mu$ を
大きくした後には実現できる。実際、$R[x_1, \ldots, x_n]^\wedge$ は Noether なので、
イデアル $\Ker(\Phi)$ の生成元 $g_1, \ldots, g_l$ を選べる。このとき
$$
\Psi(g_1), \ldots, \Psi(g_l) \in
\Ker(A_\mu^\wedge \to C_\mu^\wedge) = \Ker(A_\mu \to C_\mu)
$$
は $\Ker(A \to C) = \colim \Ker(A_\mu \to C_\mu)$ において零へ写る。
したがって前と同様に $\mu$ を大きくすれば、望む結果を得る。

\medskip\noindent
全射性の証明：後始末。上で構成した連続環準同型 $B \to (A_\mu)^\wedge$ は射
$\hat x_\mu : V_{\mu, /Z_\mu} \to W$ を定める。$\hat x_\mu$ と $u'_\mu$ の両立性は、
環準同型 $B \to (A_\mu)^\wedge$ が構成により環準同型 $A_\mu \to C_\mu$ と
両立することから従う。実際、この両立性は、構成で用いた固有射
$V'_\mu \to V_\mu$ と射 $x'_\mu$ および
$\hat x'_\mu = x'_{\mu, /Z_\mu}$ によって証示される。これで証明が完了する。
\end{proof}

\begin{lemma}
\label{artin-lemma-rs}
状況 \ref{artin-situation-contractions} において、関手 $F$ は Rim--Schlessinger 条件 (RS) を
満たす。
\end{lemma}

\begin{proof}
この条件は、Artin 局所環のスペクトルである $S$ 上のスキーム $V$ における関手 $F$ の
値 $F(V)$ と、Artin 局所環のスペクトルである $S$ 上のスキームの射 $V_1 \to V_2$ に
対する制限写像 $F(V_2) \to F(V_1)$ だけを用いることを思い出そう。
そこで $V/S$ を Artin 局所環のスペクトルとする。
$\xi = (Z, u', \hat x) \in F(V)$ ならば、集合論的に $Z = \emptyset$ または $Z = V$ で
ある。前者の場合、$\hat x$ は空の形式代数空間から $W$ への射である。
後者の場合、$u'$ は空スキームから $X'$ への射であり、
$\hat x : V \to W$ は $W$ への射である。したがって
$$
F(V) = U'(V) \amalg W(V)
$$
であり、さらに上のような $V_1 \to V_2$ に対して、誘導される写像
$F(V_2) \to F(V_1)$ はこの分解と両立する。ゆえに $U'$ と $W$ がともに
Rim--Schlessinger 条件を満たすことを証明すれば十分である。$U'$ については補題
\ref{artin-lemma-algebraic-stack-RS} から従う。$W$ について示すため、形式空間の補題
\ref{formal-spaces-lemma-structure-locally-noetherian} におけるように
$W = \colim W_n$ と書く。$V = \Spec(A)$ とし、$(A, \mathfrak m)$ を Artin 局所環とする。
$\mathfrak m^n = 0$ となる $n \geq 1$ を選ぶ。このとき $W(V) = W_n(V)$ である。
したがって $W$ に対する Rim--Schlessinger 条件は、すべての $n$ に対する $W_n$ の
Rim--Schlessinger 条件から従う（後者は補題 \ref{artin-lemma-algebraic-stack-RS} から従う）。
\end{proof}

\begin{lemma}
\label{artin-lemma-finite-dim}
状況 \ref{artin-situation-contractions} において、関手 $F$ の接空間は有限次元である。
\end{lemma}

\begin{proof}
補題 \ref{artin-lemma-rs} の証明で、$V$ が Artin 局所環のスペクトルならば
$F(V) = U'(V) \amalg W(V)$ であることを見た。接空間は $S$ 上のこのようなスキームに
おける $F$ の値だけから計算される。したがって関手 $U'$ と $W$ の接空間が有限次元で
あることを証明すれば十分である。$U'$ については補題 \ref{artin-lemma-finite-dimension} から
従う。補題 \ref{artin-lemma-rs} の証明におけるように $W = \colim W_n$ と書く。
接空間を得るには、$\mathfrak m^2 = 0$ である Artin 局所環 $(A, \mathfrak m)$ の
スペクトルにおける $W$ の値だけが必要なので、$W$ の接空間は $W_2$ の接空間に等しい。
補題 \ref{artin-lemma-finite-dimension} により、$W_2$ の接空間も有限次元である。
\end{proof}

\begin{lemma}
\label{artin-lemma-formal-object-effective}
状況 \ref{artin-situation-contractions} において $X' \to S$ は分離であると仮定する。
このとき $F$ のすべての形式対象は有効である。
\end{lemma}

\begin{proof}
$F$ の形式対象 $\xi = (R, \xi_n)$ は、剰余体が $S$ 上有限型である Noether 完備局所
$S$-代数 $R$ と、すべての $n$ に対する元
$\xi_n \in F(\Spec(R/\mathfrak m^n))$ で、
$\xi_{n + 1}|_{\Spec(R/\mathfrak m^n)} = \xi_n$ となるものからなる。
補題 \ref{artin-lemma-rs} の証明の議論により、$\xi$ は $U'$ の形式対象であるか、$W$ の
形式対象である。前者の場合、補題 \ref{artin-lemma-effective} により $\xi$ は有効である。
後者が興味深い場合である。$V = \Spec(R)$ とおく。すべての $n \geq 1$ に対して
$F(\Spec(R/\mathfrak m^n))$ における像が $\xi_n$ となるような元
$(Z, u', \hat x) \in F(V)$ を構成する。

\medskip\noindent
元 $\xi_n$ の集まりを、$S$ 上の局所 Noether 形式代数空間の射
$$
\xi : \text{Spf}(R) \longrightarrow W
$$
とみなせる。一般には $\xi$ は {\it 進的でない}射であることに注意する。
これを直すため、$I \subset R$ を形式的閉部分空間
$$
\text{Spf}(R) \times_{\xi, W} W_{red} \subset \text{Spf}(R)
$$
に対応するイデアルとする。$I \subset \mathfrak m_R$ であることに注意する。
$Z = V(I) \subset V = \Spec(R)$ とおく。$R$ は $\mathfrak m_R$-進完備なので、
特に $I$-進完備である（代数の補題 \ref{algebra-lemma-complete-by-sub}）。
さらに各 $n \geq 1$ に対して、射
$$
\xi|_{\text{Spf}(R/I^n)} :
\text{Spf}(R/I^n)
\longrightarrow
W
$$
は実際に射
$$
\xi'_n :  \Spec(R/I^n) \longrightarrow W
$$
から生じると主張する。実際、補題 \ref{artin-lemma-rs} の証明におけるように
$W = \colim W_n$ と書き、$\xi|_{\text{Spf}(R/I^n)}$ が $W_n$ へ写ることに注意し、
形式空間の補題 \ref{formal-spaces-lemma-map-into-algebraic-space} を適用して
これを望みどおり射 $\Spec(R/I^n) \to W_n$ へ代数化すればよい。
$\text{Spf}'(R) = V_{/Z}$ を、$I$-進位相を備えた $R$ の形式スペクトル、同値な言い方では
$Z$ に沿った $V$ の形式完備化と書く。射 $\xi'_n$ を用いて、$S$ 上の局所 Noether
形式代数空間の進的射
$$
\hat x = (\xi'_n) : \text{Spf}'(R) \longrightarrow W
$$
を得る。基底変換
$$
\text{Spf}'(R) \times_{\hat x, W, g} X'_{/T'} \longrightarrow \text{Spf}'(R)
$$
を考える。これは形式空間の代数化の補題
\ref{restricted-lemma-base-change-formal-modification} により形式的修正である。
したがって膨張に関する主定理（形式空間の代数化の定理
\ref{restricted-theorem-dilatations}）により、$\Spec(R) \setminus V(I)$ 上で同型であり、
その完備化が上の形式的修正を復元するような固有射
$$
V' \longrightarrow V = \Spec(R)
$$
を得る。言い換えると
$$
V' \times_{\Spec(R)} \Spec(R/I^n) =
\Spec(R/I^n) \times_{\xi'_n, W, g} X'_{/T'}
$$
である。特に、両立する射の系
$$
V' \times_{\Spec(R)} \Spec(R/I^n) \longrightarrow X' \times_S \Spec(R/I^n)
$$
を得る。したがって Grothendieck の代数化定理（空間の射続論の補題
\ref{spaces-more-morphisms-lemma-algebraize-morphism} の形）により、上に表示した射を
復元する $S$ 上の射
$$
x' : V' \to X'
$$
を得る。最後に $u' : V \setminus Z \to X'$ を $x'$ の
$V \setminus Z \subset V'$ への制限とおけば、望む元
$(Z, u', \hat x) \in F(V)$ の第三の成分が得られる。
\end{proof}

\begin{lemma}
\label{artin-lemma-openness-smoothness}
$S$ を局所 Noether スキームとする。$V$ を $S$ 上局所有限型なスキームとし、
$Z \subset V$ を閉とする。$W$ を $S$ 上の局所 Noether 形式代数空間で、$W_{red}$ が
$S$ 上局所有限型であるものとする。$g : V_{/Z} \to W$ を $S$ 上の形式代数空間の
進的射とする。$v \in V$ を、$g$ が $v$ において半普遍的となる閉点とする
（節 \ref{artin-section-axioms-functors} における意味である）。このとき $V$ を $v$ の
開近傍で置き換えた後、射 $g$ は滑らかである（証明を参照せよ）。
\end{lemma}

\begin{proof}
$g$ は進的なので、代数空間で表現可能である（形式空間の節
\ref{formal-spaces-section-adic}）。したがって $g$ が滑らかであるとは、ブートストラップの
定義 \ref{bootstrap-definition-property-transformation} の意味で $g$ が滑らかであることをいう。

\medskip\noindent
形式空間の補題 \ref{formal-spaces-lemma-structure-locally-noetherian} におけるように
$W = \colim W_n$ と書く。$V_n = V_{/Z} \times_{\hat x, W} W_n$ とおく。
このとき $V_n$ は台集合が $Z$ である閉部分スキームである。$V \to W$ の滑らかさは、
すべての射 $V_n \to W_n$ の滑らかさと同値である（これは、$T$ が準コンパクトな
スキームである任意の射 $T \to W$ が、ある $n$ に対して $W_n$ を経由するためである）。
補題 \ref{artin-lemma-base-change-versal} により、射 $V_n \to W_n$ は $v$ において
滑らかであることがわかっている\footnote{この補題を適用できるのは、形式空間の補題
\ref{formal-spaces-lemma-diagonal-morphism-formal-algebraic-spaces} により $W$ の対角射が
代数空間で表現可能かつ局所有限型であり、さらに補題 \ref{artin-lemma-rs} の証明で
$W$ が (RS) をもつことを見たからである。}。もちろんこれは、任意の $n$ が与えられた
とき、$V_n \to W_n$ が滑らかとなるように $V$ を縮小できることを意味する。
問題は、すべての $n$ に同時に使える開集合を見いだすことである。

\medskip\noindent
問題は $V$ 上局所的なので、$S = \Spec(R)$ と $V = \Spec(A)$ はアフィンであると
仮定してよい。

\medskip\noindent
本段落では $W$ がアフィン形式代数空間である場合へ帰着する。アフィン形式スキーム
$W'$ と、$W_{red}$ における $v$ の像が $W'_{red} \to W_{red}$ の像に入るような
エタール射 $W' \to W$ を選ぶ。このとき
$V_{/Z} \times_{g, W} W' \to V_{/Z}$ は $S$ 上の形式代数空間の進的エタール射であり、
$V_{/Z} \times_{g, W} W'$ はアフィン形式代数空間である。形式空間の代数化の補題
\ref{restricted-lemma-algebraize-rig-etale-affine} により、$Z' = \varphi^{-1}(Z)$ に沿った
$V'$ の完備化が $V_{/Z}$ 上で $V_{/Z} \times_{g, W} W'$ と同型となるような、
アフィンスキームのエタール射 $\varphi : V' \to V$ が存在する。$v$ はある
$v' \in V'$ の像であることに注意する。滑らかさは基底変換で保たれるので、
$V'_n \to W'_n$ はすべての $n$ に対して滑らかである。次段落では、$V'$ を $v'$ の
開近傍で置き換えた後、射 $V'_n \to W'_n$ がすべての $n$ に対して滑らかであることを
示す。次に $V$ を $V' \to V$ の開な像で置き換えれば、滑らかさのエタール降下により
$V_n \to W_n$ が滑らかであることを得る。いくつかの詳細は省略する。

\medskip\noindent
$S = \Spec(R)$、$V = \Spec(A)$、$Z = V(I)$、$W = \text{Spf}(B)$ と仮定する。
$v$ は極大イデアル $I \subset \mathfrak m \subset A$ に対応するとする。
進的連続 $R$-代数準同型
$$
B \longrightarrow A^\wedge
$$
が与えられている。$\mathfrak b \subset B$ を位相的冪零元のイデアルとする
（これは Noether 進位相環 $B$ の最大定義イデアルである）。
$\mathfrak b A^\wedge$ と $IA^\wedge$ はともに Noether 進位相環 $A^\wedge$ の
定義イデアルであることに注意する。また $\mathfrak m A^\wedge$ は
$\mathfrak b A^\wedge$ と $IA^\wedge$ の両方を含む $A^\wedge$ の極大イデアルである。
すべての $n$ に対して
$$
B_n = B/\mathfrak b^n \to A^\wedge/\mathfrak b^n A^\wedge = A_n
$$
は $\mathfrak m$ において滑らかであることが与えられている。上の議論により、
$B_1 \to A_1$ は滑らかな環準同型であると仮定する。
$\mathfrak m_1 \subset A_1$ を $\mathfrak m$ に対応する極大イデアルと書く。
滑らかさは平坦性を含意するので、すべての $n \geq 1$ に対して写像
$$
\mathfrak b^n/\mathfrak b^{n + 1} \otimes_{B_1} (A_1)_{\mathfrak m_1}
\longrightarrow
\left(\mathfrak b^nA^\wedge/\mathfrak b^{n + 1}A^\wedge\right)_{\mathfrak m_1}
$$
は同型である（代数の補題 \ref{algebra-lemma-what-does-it-mean-again} を参照せよ）。
Rees 代数
$$
B' = \bigoplus\nolimits_{n \geq 0} \mathfrak b^n/\mathfrak b^{n + 1}
$$
を考える。これは Noether 環 $B_1$ 上の有限型次数付き代数である。また Rees 代数
$$
A' = \bigoplus\nolimits_{n \geq 0}
\mathfrak b^nA^\wedge/\mathfrak b^{n + 1}A^\wedge
$$
を考える。これは Noether 環 $A_1$ 上の有限型次数付き代数である。次数付き
$A_1$-代数の準同型
$$
\Psi : B' \otimes_{B_1} A_1 \longrightarrow A'
$$
を考える。上のことから、この写像は $A_1$ の極大イデアル $\mathfrak m_1$ において
局所化した後に同型である。したがって $\Ker(\Psi)$、それぞれ $\Coker(\Psi)$ は
$B' \otimes_{B_1} A_1$、それぞれ $A'$ 上の有限加群で、その $\mathfrak m_1$ における
局所化は零である。ゆえに $A_1$（および対応して $A$）を単項局所化で置き換えた後、
$\Psi$ は同型であると仮定してよい。（これが証明の要点である。）
すると逆向きにたどることにより、$B_n \to A_n$ は平坦である。代数の補題
\ref{algebra-lemma-what-does-it-mean-again} を参照せよ。したがって
$A_n \to B_n$ は滑らかであり（滑らかなファイバーをもつ平坦環準同型として、
代数の補題 \ref{algebra-lemma-flat-fibre-smooth} を参照せよ）、証明が完了する。
\end{proof}

\begin{lemma}
\label{artin-lemma-openness-versality}
状況 \ref{artin-situation-contractions} において、関手 $F$ は半普遍性の開性を満たす。
\end{lemma}

\begin{proof}
次を示さなければならない。$V$ を $S$ 上局所有限型なスキームとし、
$\xi \in F(V)$ とし、$v_0 \in V$ を $\xi$ が $v_0$ において半普遍的となる有限型点とする。
このとき $V$ を $v_0$ の開近傍で置き換えた後、$\xi$ は $V$ のすべての有限型点で
半普遍的である。$\xi = (Z, u', \hat x)$ と書く。

\medskip\noindent
第一の場合：$v_0 \not \in Z$。まず $V$ を $V \setminus Z$ で置き換えてよい。
したがって $\xi = (\emptyset, u', \emptyset)$ であり、射 $u' : V \to X'$ は
$v_0$ において半普遍的である。空間の射続論の補題
\ref{spaces-more-morphisms-lemma-lifting-along-artinian-at-point} により、これは
$u' : V \to X'$ が $v_0$ において滑らかであることを意味する。射が滑らかである点の
集合は開なので、$V$ を縮小した後 $u'$ は滑らかであると仮定できる。
すると同じ補題により、望みどおり $\xi$ はすべての点で半普遍的である。

\medskip\noindent
第二の場合：$v_0 \in Z$。形式空間の補題
\ref{formal-spaces-lemma-structure-locally-noetherian} におけるように
$W = \colim W_n$ と書く。補題 \ref{artin-lemma-openness-smoothness} により、
$\hat x : V_{/Z} \to W$ は形式代数空間の滑らかな射であると仮定してよい。
直ちに、$\xi = (Z, u', \hat x)$ は $Z$ のすべての有限型点で半普遍的であることが従う。
$V' \to V$、$\hat x'$、$x'$ が $u'$ と $\hat x$ の間の両立性を証示するとする。
$\hat x' : V'_{/Z} \to X'_{/T'}$ は $\hat x$ の基底変換として滑らかである。
$\hat x'$ は $x' : V' \to X'$ の完備化なので、既に用いた空間の射続論の補題
\ref{spaces-more-morphisms-lemma-lifting-along-artinian-at-point} により、
$x' : V' \to X'$ は
$(V' \to V)^{-1}(Z) = |x'|^{-1}(T') \subset |V'|$ のすべての点で滑らかである。
射の滑らかな点の集合は開なので、$x'$ が滑らかでない点からなる閉集合
$B \subset |V'|$ は $(V' \to V)^{-1}(Z)$ と交わらない。
$V' \to V$ は固有、したがって閉なので、$(V' \to V)(B) \subset V$ は $Z$ と交わらない
閉部分集合である。したがって $V$ を縮小した後、$B = \emptyset$、すなわち $x'$ は
滑らかであると仮定してよい。前段落の議論により、これはまさに $\xi$ が $Z$ に
含まれない $V$ のすべての有限型点で半普遍的であることを意味し、証明が完了する。
\end{proof}

\noindent
最後の結果を述べる。

\begin{theorem}
\label{artin-theorem-contractions}
\begin{reference}
\cite[Theorem 3.1]{ArtinII}
\end{reference}
$S$ を、すべての有限型点 $s \in S$ に対して $\mathcal{O}_{S, s}$ が G-環であるような
局所 Noether スキームとする。$X'$ を $S$ 上局所有限型な代数空間とする。
$T' \subset |X'|$ を閉部分集合とする。$W$ を $S$ 上の局所 Noether 形式代数空間で、
$W_{red}$ が $S$ 上局所有限型であるものとする。最後に
$$
g : X'_{/T'} \longrightarrow W
$$
を形式的修正とする。形式空間の代数化の定義
\ref{restricted-definition-formal-modification} を参照せよ。$X'$ と $W$ が $S$ 上
分離ならば\footnote{注意 \ref{artin-remark-separated-needed} を参照せよ。}、
$S$ 上の代数空間の固有射 $f : X' \to X$、閉部分集合 $T \subset |X|$、および
形式代数空間の同型 $a : X_{/T} \to W$ で、次を満たすものが存在する：
\begin{enumerate}
\item $T'$ は $|f| : |X'| \to |X|$ による $T$ の逆像である；
\item $f : X' \to X$ は $X' \setminus T'$ を $X \setminus T$ へ同型に写す；
\item $g = a \circ f_{/T}$ である。ここで
$f_{/T} : X'_{/T'} \to X_{/T}$ は誘導される射である。
\end{enumerate}
言い換えると、$(f : X' \to X, T, a)$ は本節で先に定義した解である。
\end{theorem}

\begin{proof}
$F$ を、本節で $X'$、$T'$、$W$、$g$ を用いて構成した関手とする。
補題 \ref{artin-lemma-functor-is-solution} により、$F$ が $S$ 上局所有限型な代数空間 $X$ に
対応することを示せば十分である。このため、命題
\ref{artin-proposition-spaces-diagonal-representable-noetherian} を適用する。
実際、補題 \ref{artin-lemma-diagonal-contractions} により $F$ の対角射は閉埋入で表現可能であり、
補題 \ref{artin-lemma-sheaf}、\ref{artin-lemma-limit-preserving}、\ref{artin-lemma-rs}、
\ref{artin-lemma-finite-dim}、\ref{artin-lemma-formal-object-effective}、および
\ref{artin-lemma-openness-versality} により、公理 [0]、[1]、[2]、[3]、[4]、[5] が成り立つ。
\end{proof}

\begin{remark}
\label{artin-remark-separated-needed}
定理 \ref{artin-theorem-contractions} の証明では、$X'$ と $W$ が $S$ 上分離であることを
二か所で用いる。第一に、$\Delta : F \to F \times F$ が代数空間で表現可能であることを
示す際に用いる。この仮定の使用は、注意 \ref{artin-remark-diagonal} で得た三つ組
$E'$、$(E' \to V)^{-1}Z$、$E'_{/Z} \to E_W$ に分離の場合の定理を適用して
$\Delta$ が表現可能であることを証明すれば、完全に避けられる
（これは対角射に対する通常のブートストラップ手続きである）。したがって
補題 \ref{artin-lemma-formal-object-effective} の証明は、われわれの
定理 \ref{artin-theorem-contractions} の証明で $X' \to S$ が分離であることを本当に使う
唯一の箇所である。読者は、この仮定を
射 $x' : V' \to X'$ を得るためにだけ用いていることを確認できる。文献の結果を用いると、
$X' \to S$ が準分離であれば $x'$ の存在を示せる。空間の射続論の注意
\ref{spaces-more-morphisms-remark-weaken-separation-axioms-question} を参照せよ。
したがって定理は、述べられた「分離」を「準分離」で置き換えても成り立つ。
これが必要になれば、ここで正確に述べ、注意深く証明することにする。
\end{remark}
